\documentclass[12pt,a4paper,openany,oneside]{book}
\usepackage[utf8]{vietnam}
\usepackage{xcolor}
\usepackage{diagbox}
\usepackage[shortlabels]{enumitem}
\usepackage[a4paper,left=1.8cm,right=1.5cm,top=1.5cm,bottom=1.6cm]{geometry}
\usepackage{setspace}
\usepackage{fontawesome}
\setstretch{1.1}
\setlength{\parskip}{6pt}
\usepackage{amsmath}
\allowdisplaybreaks[4]
\usepackage{amssymb}
\usepackage{pgfornament}
%%%============================================%%%
\DeclareSymbolFont{symbolsC}{U}{txsyc}{m}{n}
\DeclareMathSymbol{\varparallel}{\mathrel}{symbolsC}{9}
\renewcommand{\parallel}{\mathbin{\!\raisebox{0.25pt}/\mkern-2mu\raisebox{0.25pt}/\!}}
\renewcommand{\backsim}{\ \mathord{\raisebox{0.025\baselineskip}{\rotatebox{90}{\fontfamily{qag}\selectfont S}}}\ }
\setlength\arraycolsep{2pt}
\newcommand{\hoac}[1]{\renewcommand{\arraystretch}{0.85}\left[\begin{aligned}#1\end{aligned}\right.}
\newcommand{\heva}[1]{\renewcommand{\arraystretch}{0.85}\left\{\begin{aligned}#1\end{aligned}\right.}
\DeclareSymbolFont{AMSb}{U}{msb}{m}{n}
\makeatletter
\DeclareSymbolFontAlphabet{\math@bb}{AMSb}
\AtBeginDocument{\protected\def\mathbb{\math@bb}}
\g@addto@macro\normalsize{
\setlength{\abovedisplayskip}{4pt}%
\setlength{\belowdisplayskip}{2pt}%
\setlength{\abovedisplayshortskip}{4pt}%
\setlength{\belowdisplayshortskip}{1pt}%
}
\makeatother
\def\bmod{\operatorname{mod}\,}
\makeatletter
\DeclareRobustCommand
\vdots{\,\vbox{\baselineskip3\p@ \lineskiplimit\z@
\hbox{.}\hbox{.}\hbox{.}}\,}
\makeatother
\usepackage{mathrsfs}
\usepackage{esvect}
\def\vec{\vv} %vector
\def\overrightarrow{\vv}
\makeatletter
\newcommand\xleftrightarrow[2][]{%
\ext@arrow 9999{\longleftrightarrowfill@}{#1}{#2}}
\newcommand\longleftrightarrowfill@{%
\arrowfill@\leftarrow\relbar\rightarrow}
\makeatother

\usepackage{amsmath}
\allowdisplaybreaks[4]
\usepackage{amssymb}
\usepackage{graphicx}
\usepackage{fontawesome,pifont,oplotsymbl}
\usepackage{marvosym}
\usepackage{savesym}
\savesymbol{Cross}
\usepackage{bbding}
\restoresymbol{bb}{Cross}
\usepackage{tikz,pgfplots}
\usepackage{tikz-3dplot}
\usetikzlibrary{calc,intersections,shapes.geometric,arrows,decorations.markings,arrows.meta,patterns,fadings,
shadows,backgrounds}
\tikzfading[name=fade out, inner color=transparent!0,outer color=transparent!100]
\tikzfading[name=fade left, top color=transparent!0,bottom color=transparent!0,middle color=transparent!100]
\tikzfading[name=fade right,left color=transparent!0,right color=transparent!100]
\usepackage{graphicx}%chỉnh hình ảnh
\usepackage{makecell}
\usepackage{array}
\newcolumntype{L}[1]{>{\raggedright\let
\newline\\\arraybackslash\hspace{0pt}}m{#1}}
\newcolumntype{C}[1]{>{\centering\let
\newline\\\arraybackslash\hspace{0pt}}m{#1}}
\newcolumntype{R}[1]{>{\raggedleft\let
\newline\\\arraybackslash\hspace{0pt}}m{#1}}
\renewcommand{\arraystretch}{1.25}
\usepackage[most,xparse,many]{tcolorbox}
\newcommand\circlenum[2][cyan]{\tikz[baseline=(char.base)]
{\node[shape=circle,inner sep=1pt,draw=#1,fill=#1!10,
font=\bfseries\footnotesize,minimum size=14pt,outer sep=0pt] (char) {#2};}}
\newcommand\circlenumH[2][cyan]{\tikz[baseline=(char.base)]
{\node[shape=circle,inner sep=0.5pt,draw=#1,%fill=#1!10,
font=\bfseries\footnotesize,minimum size=13pt,outer sep=0pt] (char) {#2};}}
\renewenvironment{center}
{\parskip=3pt\par\nopagebreak\centering}
{\par\noindent\ignorespacesafterend}
\usepackage{multicol}
\setlength{\multicolsep}{2pt plus 1pt minus 1pt}
%%%============================================%%%
\DeclareSymbolFont{symbolsC}{U}{txsyc}{m}{n}
\DeclareMathSymbol{\varparallel}{\mathrel}{symbolsC}{9}
\renewcommand{\parallel}{\mathbin{\!\raisebox{0.25pt}/\mkern-2mu\raisebox{0.25pt}/\!}}
\renewcommand{\backsim}{\ \mathord{\raisebox{0.025\baselineskip}{\rotatebox{90}{\fontfamily{qag}\selectfont S}}}\ }
\setlength\arraycolsep{2pt}
%\newcommand{\hoac}[1]{\renewcommand{\arraystretch}{0.85}\left[\begin{aligned}#1\end{aligned}\right.}
%\newcommand{\heva}[1]{\renewcommand{\arraystretch}{0.85}\left\{\begin{aligned}#1\end{aligned}\right.}
\DeclareSymbolFont{AMSb}{U}{msb}{m}{n}
\makeatletter
\DeclareSymbolFontAlphabet{\math@bb}{AMSb}
\AtBeginDocument{\protected\def\mathbb{\math@bb}}
\g@addto@macro\normalsize{
\setlength{\abovedisplayskip}{4pt}%
\setlength{\belowdisplayskip}{2pt}%
\setlength{\abovedisplayshortskip}{4pt}%
\setlength{\belowdisplayshortskip}{1pt}%
}
\makeatother
\def\bmod{\operatorname{mod}\,}
\makeatletter
\DeclareRobustCommand
\vdots{\,\vbox{\baselineskip3\p@ \lineskiplimit\z@
\hbox{.}\hbox{.}\hbox{.}}\,}
\makeatother

\renewenvironment{center}
 {\parskip=0pt\par\nopagebreak\centering}
 {\par\noindent\ignorespacesafterend}
 
\newcommand{\indam}[1]{\textcolor{mausubsection}{\textbf{#1}}}
\newcommand{\indamm}[1]{\textcolor{mausubsection}{\textbf{#1}}}
\definecolor{mauphan}{HTML}{FFA500}
\definecolor{mausubsection}{HTML}{CD1D1D}
\definecolor{mauchuong}{HTML}{AE0E0E}
\definecolor{doimauchuong}{HTML}{0EAEAE}
%\definecolor{mausection}{HTML}{EA5548}
\definecolor{mausection}{HTML}{008bb2}

\usepackage[explicit]{titlesec} %chỉnh định của section, chapter #1
\usepackage{pgfornament}
\titlespacing*{\section}{0pt}{5pt}{3pt}
\titleformat{\section}{\filcenter}{%
	\begin{tikzpicture}[declare function={r = 0.79;},line cap=round]
		\path[rounded corners=3pt,fill=mausection!30,shift={(-15:3pt)}] (90:r)--(150:r)--(210:r)--(270:r)--(330:r)--(30:r)--cycle;
		\path[rounded corners=3pt,fill=mausection] (90:r)--(150:r)--(210:r)--(270:r)--(330:r)--(30:r)--cycle;
		\draw[line width=1pt,mausection] (-90:1.35*r) coordinate (A)--+(30:1.25*r)coordinate (B)
			--([turn]-30:0.85*\linewidth) coordinate (C);
		\fill[mausection] (A) circle (2pt);
		\fill[mausection] (C) circle (3pt);
		\draw[opacity=0.45,white,line width=3pt] (150:{0.8*r})--+(30:{0.3*r});
		\fill[opacity=0.45,white] ($(150:{0.8*r})+(30:{0.55*r})$) circle (1.65pt);
		\path let \p1=($(C)-(B)-(0.5,0)$) in  ($(B)!0.5!(C)+(0,2pt)$) 
			node[text width=\x1,anchor=south,font=\bfseries\sffamily\Large,
						inner sep=1pt,text=mausection,align=center]{\MakeUppercase{#1}};
		\path (0,0) node[font=\bfseries\sffamily,text=white]{ĐỀ \arabic{section}};		
	\end{tikzpicture}}{1pt}{\setcounter{ex}{0}\setcounter{bt}{0}}

\titleformat{name=\section,numberless}[display]{\addcontentsline{toc}{section}{#1 \arabic{chapter}}}%
	{\begin{tikzpicture}[color=teal,declare function={bk=7pt;}]
	  	\path (0,0) node[inner sep=9pt,outer sep=0pt,font=\fontfamily{pag}\bfseries\Large\selectfont](A){\phantom{I}ÔN TẬP CHƯƠNG \arabic{chapter}\phantom{I}};
	  	\path (A.north west) coordinate (Tt)
	  		(A.north east) coordinate (Tp)
	  		(A.south west) coordinate (Dt)
	  		(A.south east) coordinate (Dp);
	  	\path[ultra thick,draw=teal!35,fill] ($(Tt)!bk!(Tp)$)--($(Tp)!bk!(Tt)$)--($(Tp)!bk!(Dp)$)--($(Dp)!bk!(Tp)$)
	  		--($(Dp)!bk!(Dt)$)--($(Dt)!bk!(Dp)$)--($(Dt)!bk!(Tt)$)--($(Tt)!bk!(Dt)$)--cycle;
	 	 \path (0,0) node[inner sep=9pt,outer sep=0pt,font=\bfseries\Large,text=white]{\phantom{I}\MakeUppercase{#1} \arabic{chapter}\phantom{I}};
	  \end{tikzpicture}}{3pt}{}
	  
\titlespacing*{\subsection}{0pt}{3pt}{3pt}
\titleformat{\subsection}{}{%
	\begin{tikzpicture}[declare function={r = 0.35;},line cap=round]
		\path[fill=mauchuong!25] (-35:2pt) circle (r);
		\path[fill=mauchuong] (0,0) circle (r);
		\draw[opacity=0.45,white,line width=2pt] ({150:{0.75*r}}) arc (150:125:{0.75*r});
		\path (0,0) node[font=\bfseries\large\sffamily,text=white]{\resizebox{!}{0.27cm}{\Roman{subsection}}};
		\path ({1.25*r},0) node[anchor=west,font=\bfseries\sffamily\Large,text width=16.5cm,text=mausubsection]{#1};
	\end{tikzpicture}}{0.2em}{}

%\titlespacing*{\subsubsection}{0.65cm}{3pt}{3pt}
\titlespacing*{\subsubsection}{0pt}{3pt}{3pt}
\titleformat{\subsubsection}{}{%
	 \bfseries\sffamily\large\color{mausubsection}\arabic{subsubsection}. #1%
	}{0em}{}
	
\usepackage{titletoc} %chỉnh định dạng của mục lục

\usepackage{pdfrender}
%%Đoạn lệnh đổi mục lục
\titlespacing*{\chapter}{0pt}{60pt}{-15pt}
\def\fontchuong{23}
\titleformat{\chapter}[display]{\thispagestyle{empty}}{%
\begin{tikzpicture}[remember picture,overlay]
	\path (current page.north west) coordinate (pA);
	\begin{scope}[shift={($(pA)+(0,-8)$)}]
		\def\duong{(0,2.55){[rounded corners=19pt]--++(17,0)}--(19,8)--(0,8)--cycle;}
		\begin{scope}
			\path[fill=mauchuong,rounded corners] \duong		
			\clip \duong;
			\foreach \i in {0,...,12}{
			  \foreach \j in {-2,...,15}{
			  	\foreach \a in {0,120,-120} \draw[line width=1pt,mauchuong!87,scale=0.5] (3*\i,2*sin{60}*\j) -- +(\a:1);
			  	\foreach \a in {0,120,-120} \draw[line width=1pt,mauchuong!87,scale=0.5] (3*\i+3*cos{60},2*sin{60}*\j+sin{60}) -- +(\a:1);
			  }}				
		\end{scope}
		\def\duongtrai{(0,0)..controls +(30:3) and +(210:2)..(3,6)..controls +(30:2) and +(250:3)..(8,8)--(0,8)--cycle;}
		\begin{scope}
			\path[left color=doimauchuong!90!black,right color=doimauchuong] \duongtrai		
			\clip \duongtrai;
			\foreach \i in {-2,...,5}{
			  \foreach \j in {-7,...,15}{
			  	\foreach \a in {0,120,-120} \draw[line width=1pt,doimauchuong!97,scale=0.5] (3*\i,2*sin{60}*\j) -- +(\a:1);
			  	\foreach \a in {0,120,-120} \draw[line width=1pt,doimauchuong!97,scale=0.5] (3*\i+3*cos{60},2*sin{60}*\j+sin{60}) -- +(\a:1);
			  }}					
		\end{scope}
		\draw[white,line width=3pt] (0,0)..controls +(30:3) and +(210:2)..(3,6)..controls +(30:2) and +(250:3)..(8,8);
		\fill[white] (2.1,4.15) coordinate (O) circle (1);
%		\path[fill=green,fill opacity=0.35,even odd rule] 
%			(O) circle (1)
%			($(O)+(0:1.6)$) arc (0:270:1.6)--+(1.6,0)--cycle;
		\path[inner color=green,outer color=green!65,even odd rule,opacity=0.45] 
			(O) circle (1)
			($(O)+(0:1.6)$) arc (0:270:1.6)--+(1.6,0)--cycle;
		\path[
			postaction={
			decoration={raise=2pt,
			text along path,
			text format delimiters={|}{|},
			text={%
			|\bfseries\footnotesize\sffamily|
			{C}{H}{Ư}{Ơ}{N}{G}{}
			}},
			decorate}
			] ($(O)+(190:1.03)$) arc (190:60:1.03);
		\path (O) node[font=\bfseries\fontfamily{ugq}\fontsize{39pt}{0pt}\selectfont]{\arabic{chapter}};
   \end{scope}
   \path ($(O)+(2,0.25)$)
   node[anchor=west,text width=13cm,text=white,
   	font=\bfseries\fontfamily{pag}\fontsize{\fontchuong pt}{0pt}\selectfont]{\MakeUppercase{#1}};
\end{tikzpicture}%
}{0pt}{}

\titleformat{name=\chapter,numberless}[display]
	{\thispagestyle{empty}}%
	{\begin{tikzpicture}[remember picture,overlay]
		\path (current page.north west) coordinate (pA);
		\begin{scope}[shift={($(pA)+(0,-8)$)}]
			\def\duong{(0,2.55){[rounded corners=19pt]--++(17,0)}--(19,8)--(0,8)--cycle;}
			\begin{scope}
				\path[fill=mauchuong,rounded corners] \duong		
				\clip \duong;
				\foreach \i in {0,...,12}{
				  \foreach \j in {-2,...,15}{
				  	\foreach \a in {0,120,-120} \draw[line width=1pt,mauchuong!87,scale=0.5] (3*\i,2*sin{60}*\j) -- +(\a:1);
				  	\foreach \a in {0,120,-120} \draw[line width=1pt,mauchuong!87,scale=0.5] (3*\i+3*cos{60},2*sin{60}*\j+sin{60}) -- +(\a:1);
				  }}				
			\end{scope}
			\def\duongtrai{(0,0)..controls +(30:3) and +(210:2)..(3,6)..controls +(30:2) and +(250:3)..(8,8)--(0,8)--cycle;}
			\begin{scope}
				\path[left color=doimauchuong!90!black,right color=doimauchuong] \duongtrai		
				\clip \duongtrai;
				\foreach \i in {-2,...,5}{
				  \foreach \j in {-7,...,15}{
				  	\foreach \a in {0,120,-120} \draw[line width=1pt,doimauchuong!97,scale=0.5] (3*\i,2*sin{60}*\j) -- +(\a:1);
				  	\foreach \a in {0,120,-120} \draw[line width=1pt,doimauchuong!97,scale=0.5] (3*\i+3*cos{60},2*sin{60}*\j+sin{60}) -- +(\a:1);
				  }}					
			\end{scope}
			\draw[white,line width=3pt] (0,0)..controls +(30:3) and +(210:2)..(3,6)..controls +(30:2) and +(250:3)..(8,8);
			\path (2.1,4.15) coordinate (O) circle (1);
	   \end{scope}
	   \path ($(O)+(2,1)$)
	   node[anchor=west,text width=13cm,text=white,align=center,font=\bfseries\fontfamily{pag}
	   \fontsize{\fontchuong pt}{0pt}\selectfont]{\MakeUppercase{#1}};
	\end{tikzpicture}%
	}
	{0pt}%
	{}%
\makeatletter
\renewcommand{\tableofcontents}{%
	{\chapter*{MỤC LỤC\,}{}}
	\vspace*{0.5cm}%
	\@starttoc{toc}
	}
\makeatother
%Hết đoạn lệnh đổi mục lục


\newif\ifnotstarredversion
\renewcommand{\thepart}{}
\makeatletter
\renewcommand*\l@part[2]{%
    \ifnum \c@tocdepth >-2\relax
    \addpenalty{-\@highpenalty}%
    \addvspace{5pt \@plus\p@}%
    \setlength\@tempdima{3em}%
    \begingroup%
   \noindent\tikz[remember picture]{
	   \path (2.5,0) node[font=\bfseries\fontfamily{pag}\selectfont,inner sep=3pt] (A){#1};
	   \path ($(A.north east)+(0.15,0)$) coordinate (At)
	   			($(A.south east)+(0.15,0)$) coordinate (Ad)
	   			(A.north west) coordinate (Bt)
	   			(A.south west) coordinate (Bd);
	   \path[left color=mauchuong!85,right color=mauchuong!65] (A.south west) rectangle ($(A.north west)-(2.5,0)$)
	   ($(Bt)!0.5!(Bd)+(-2.1,0)$) node[font=\bfseries\fontfamily{pag}\selectfont,anchor=west]{Phần};
	   \draw[ultra thick,mauchuong] let \p1=($(At)-(Ad)$), \n1 = {veclen(\x1, \y1)} in (Ad) arc (-90 :{90}:{\n1/2})--(Bt)--+(-2.5,0)--($(A.south west)-(2.5,0)$)--cycle;
	}%
    \par\smallskip
	\penalty\@highpenalty
	\setcounter{chapter}{0}
    \endgroup
    \fi
}

\def\@partimage{}
\newsavebox{\Abox}
\newcommand{\partimage}[1]{%
	\savebox{\Abox}{{#1}}%
	\gdef\@partimage{ \usebox{\Abox}}%
}

\def\@part[#1]#2{%
  \thispagestyle{empty}%
  \setcounter{chapter}{0}
  \startcontents[chapters]
  \refstepcounter{part}
  \addcontentsline{toc}{part}{ #1}%
  \markboth{Phần \Roman{part}. #1}{Phần \Roman{part}. #1}%
  {\interlinepenalty \@M
    \begin{center}
    \begin{tikzpicture}[remember picture,overlay]
 		\fill[top color=cyan!15,bottom color=white] (current page.north west) rectangle ([yshift=-7.5cm]current page.north east);
 		\begin{scope}
 			\clip (current page.north west) rectangle ([yshift=-7cm]current page.north east);
 			\draw[line width=1.25pt,white,step=0.75cm,rotate=45] ([xshift=-10cm]current page.north west) grid ([yshift=-7cm]current page.north east);
 		\end{scope}	
   		\path[mausubsection] ($(current page.north)+(-3.5,-3.5)$)  to [ornament=88] +(7,0);
   		\path ($(current page.north)+(0,-3.5)$) node[anchor=south,shift={(0,6pt)}]{\textpdfrender{
   							TextRenderingMode=FillStroke,
   							StrokeColor=mausubsection!75!blue,
   							LineWidth=1pt,
   							FillColor=mausubsection,
   						}{\fontsize{27pt}{0pt}\selectfont\sffamily\bfseries Phần #2}%
   			};
    \end{tikzpicture}%
    
    \vfil%
 	\noindent\@partimage 
    \end{center}
      }%
     \vfil\par
      \@endpart
    }
\makeatother

\renewcommand{\thechapter}{\arabic{chapter}}

\titlecontents{chapter}[2.75cm]{%
\large\sffamily\bfseries\color{mauchuong}}{%
\begin{tikzpicture}[overlay,remember picture]
	\path (0,0)--+(0.975*\linewidth,0);
	\path (0,17pt) coordinate (A)--+(0.965*\linewidth,0) coordinate (B);
	\begin{scope}[transform canvas={shift={(-2.75,1pt)}}]
		\draw[thick,mauchuong] (A)--+(2.5,0) coordinate (C);
		\draw[thick,doimauchuong] (B)--(C);
		\fill[mauchuong] (A) circle (1pt);
		\fill[doimauchuong] (B) circle (1pt);
	\end{scope}
\end{tikzpicture}%
%\makebox[2.5cm][l]{\large\sffamily\bfseries\color{mauchuong}Chủ đề \thecontentslabel.}
{\contentslabel[{Chương \thecontentslabel.}]{2.75cm}}%
}{}{\large\hfill\thecontentspage}

\renewcommand{\thesection}{\arabic{section}}
\titlecontents{section}[1.25cm]{\large\sffamily\def\tcolor{\color{black}}}{Bài \thecontentslabel. }{\color{cyan!90!blue}\gdef\tcolor{\color{cyan!90!blue}}}{\tcolor\hfill\thecontentspage}

 \newcommand{\baitap}{\par%
 \noindent%
 \begin{tikzpicture}[color=teal,declare function={bk=7pt;}]
 	\path (0,0) node[inner sep=7pt,outer sep=0pt,font=\bfseries\Large](A){\phantom{I}ĐỀ TẬP\phantom{I}};
 	\path (A.north west) coordinate (Tt)
 		(A.north east) coordinate (Tp)
 		(A.south west) coordinate (Dt)
 		(A.south east) coordinate (Dp);
 	\draw[ultra thick] ($(Tt)!bk!(Tp)$)--($(Tp)!bk!(Tt)$)--($(Tp)!bk!(Dp)$)--($(Dp)!bk!(Tp)$)
 		--($(Dp)!bk!(Dt)$)--($(Dt)!bk!(Dp)$)--($(Dt)!bk!(Tt)$)--($(Tt)!bk!(Dt)$)--cycle;
 \end{tikzpicture}}
 
 \newcommand{\ontapchuong}{\setcounter{subsection}{0}\par%
	 \noindent%
	  \begin{tikzpicture}[color=teal,declare function={bk=7pt;}]
	  	\path (0,0) node[inner sep=9pt,outer sep=0pt,font=\bfseries\Large](A){\phantom{I}ÔN TẬP CHƯƠNG \arabic{chapter}\phantom{I}};
	  	\path (A.north west) coordinate (Tt)
	  		(A.north east) coordinate (Tp)
	  		(A.south west) coordinate (Dt)
	  		(A.south east) coordinate (Dp);
	  	\path[ultra thick,draw=teal!35,fill] ($(Tt)!bk!(Tp)$)--($(Tp)!bk!(Tt)$)--($(Tp)!bk!(Dp)$)--($(Dp)!bk!(Tp)$)
	  		--($(Dp)!bk!(Dt)$)--($(Dt)!bk!(Dp)$)--($(Dt)!bk!(Tt)$)--($(Tt)!bk!(Dt)$)--cycle;
	 	 \path (0,0) node[inner sep=9pt,outer sep=0pt,font=\bfseries\Large,text=white]{\phantom{I}ÔN TẬP CHƯƠNG \arabic{chapter}\phantom{I}};
	  \end{tikzpicture}
	  \addcontentsline{toc}{section}{\hypersetup{linkcolor=cyan!90!blue}Ôn tập chương \arabic{chapter}}}
	  
\newcommand{\ontaphai}[2][1]{\setcounter{subsection}{0}%
 \par%
 \noindent%
 \addcontentsline{toc}{section}{#2}
 \ifnum #1=1
	 \begin{tikzpicture}[declare function={rt=0.5;d=sqrt(2)*rt;}]
	  \path (90:rt) coordinate (A)
	  (180:rt) coordinate (B)
	  (270:rt) coordinate (C)
	  (0:rt) coordinate (D)
	  ([shift={(0.15,0)}]0:rt) coordinate (E)
	  ($ (E)+(135:rt) $) coordinate (F)
	  ($ (E)+(-135:rt) $) coordinate (G)
	  ($ (A)!0.5!(F) $) coordinate (AF)
	  ($ (D)!0.5!(E) $) coordinate (DE)
	  ($ (C)!0.5!(G) $) coordinate (CG);
	  \path[transform canvas={shift={(-45:2.5pt)}},left color=teal!50,right color=teal!5]
	  (F)--(E)--(G)--+(0.95*\linewidth,0)--([turn]90:d)--cycle;
	  \path[right color=teal!75!black,left color=teal!80,rounded corners=2pt] (A)--(B)--(C)--(D)--cycle;
	  \draw[line width=0.65pt,teal,line cap=round,rounded corners=1.5pt]
	  (AF)--(DE)--(CG);
	  \draw[transform canvas={xscale=-1},line width=0.65pt,teal,line cap=round,rounded corners=1.5pt] (AF)--(DE)--(CG);
	  \path[left color=teal!90!black,right color=teal!50] (F)--(E)--(G)--+(0.95*\linewidth,0)--([turn]90:d)--cycle;
	  \path (-0.7pt,-0.3pt) node[inner sep=0pt,font=\itshape\bfseries\color{gray}]{\faEdit};
	  \path (-1pt,0) node[inner sep=0pt,font=\itshape\bfseries\color{white}]{\faEdit};
	  \path ([shift={(0.11,-0.01)}]E) node[align=left,anchor=west,font=\sffamily\bfseries\color{gray}]{#2};
	  \path ([shift={(0.1,0)}]E) node[align=left,anchor=west,font=\sffamily\bfseries\color{white}]{#2};
	  \end{tikzpicture}%
  \else
	\begin{tikzpicture}[color=teal,declare function={bk=7pt;}]
	  	\path (0,0) node[inner sep=9pt,outer sep=0pt,font=\bfseries\Large](A){\phantom{I}\MakeUppercase{#2}\phantom{I}};
	  	\path (A.north west) coordinate (Tt)
	  		(A.north east) coordinate (Tp)
	  		(A.south west) coordinate (Dt)
	  		(A.south east) coordinate (Dp);
	  	\path[ultra thick,draw=teal!35,fill] ($(Tt)!bk!(Tp)$)--($(Tp)!bk!(Tt)$)--($(Tp)!bk!(Dp)$)--($(Dp)!bk!(Tp)$)
	  		--($(Dp)!bk!(Dt)$)--($(Dt)!bk!(Dp)$)--($(Dt)!bk!(Tt)$)--($(Tt)!bk!(Dt)$)--cycle;
	 	 \path (0,0) node[inner sep=9pt,outer sep=0pt,font=\bfseries\Large,text=white]{\phantom{I}\MakeUppercase{#2}\phantom{I}};
	  \end{tikzpicture}%
  \fi
 }
\usepackage[most,xparse]{tcolorbox}
\tcbuselibrary{skins}
%\usepackage{etoolbox}
%%%========================Hộp ví dụ========================%%%
\newcounter{vd}
\newtcolorbox{vd}[1][]{enhanced,breakable,
	attach boxed title to top left={xshift=-4.5mm,yshift=-0.5mm},
	code={\refstepcounter{vd}},
	fonttitle=\bfseries\sffamily,
	title={
	  \if\relax\detokenize{#1}\relax
	     {\color{green!75!black}VÍ DỤ \thevd}
	   \else
	     {\color{green!75!black}VÍ DỤ \thevd\ (\MakeLowercase{#1})}
	   \fi
	},
	boxed title style={empty,arc=0pt,outer arc=0pt,boxrule=0pt},
   frame empty,
   frame hidden,	
   sharp corners,
%   interior empty,
   interior style={draw=none,fill=none},
   left=3pt,right=0pt,top=0pt,bottom=0pt,
    overlay unbroken and first ={
	    \draw[line width=1pt,green!75!black] (frame.north west)--(frame.south west);
	    \fill[green!75!black] (frame.north west) circle (1.5pt);
	    \fill[green!75!black] (frame.south west) circle (1.5pt);
    },
     overlay middle and last={
	    \draw[line width=1pt,green!75!black] (frame.north west)--(frame.south west);
	    \fill[green!75!black] (frame.north west) circle (1.5pt) 
	    	+(0,6pt) node[anchor=west,text=green!75!black,font=\small\bfseries]{Tiếp theo ví dụ \thevd};
	    \fill[green!75!black] (frame.south west) circle (1.5pt);
    }   
}


\AtBeginEnvironment{vd}{%
	\def\maudapan{white}%
	\Opensolutionfile{ansbook}[Ansbook/Tam]%
	\renewcommand{\loigiai}[1]{%
	\par\noindent%
	\textcolor{green!75!black}{\textbf{Giải}}
	  #1
	}
}
\AtEndEnvironment{vd}{\Closesolutionfile{ansbook}}

%%%========================hộp định lý========================%%%
\newcounter{dl}
\newtcolorbox{dl}[1][]{enhanced,breakable,
	attach boxed title to top left={xshift=-4mm,yshift=-0.5mm},
	code={\refstepcounter{dl}},
	fonttitle=\bfseries\sffamily,
	title={
	  \if\relax\detokenize{#1}\relax
	     {\color{mausection}ĐỊNH LÝ \thedl}
	   \else
	     {\color{mausection}ĐỊNH LÝ \thedl\ (\MakeLowercase{#1})}
	   \fi
	},
	boxed title style={empty,arc=0pt,outer arc=0pt,boxrule=0pt},
   frame empty,
   frame hidden,	
   sharp corners,
%   interior empty,
   interior style={draw=none,fill=none},
   left=3pt,right=0pt,top=0pt,bottom=0pt,
    overlay unbroken and first ={
	    \draw[line width=1pt,green!75!black] (frame.north west)--(frame.south west);
	    \fill[green!75!black] (frame.north west) circle (1.5pt);
	    \fill[green!75!black] (frame.south west) circle (1.5pt);
    },
     overlay middle and last={
	    \draw[line width=1pt,green!75!black] (frame.north west)--(frame.south west);
	    \fill[mausection] (frame.north west) circle (1.5pt) 
	    	+(0,6pt) node[anchor=west,text=mausection,font=\small\bfseries]{Tiếp theo ví dụ \thevd};
	    \fill[mausection] (frame.south west) circle (1.5pt);
    }   
}

%%%========================Hộp luyện tập========================%%%
\newcounter{lt}
\newtcolorbox{lt}[1][]{enhanced,breakable,
	attach boxed title to top left={xshift=-4mm,yshift=-0.5mm},
	code={\refstepcounter{lt}},
	fonttitle=\bfseries\sffamily,
	title={
	  \if\relax\detokenize{#1}\relax
	     {\color{orange}LUYỆN TẬP \thelt}
	   \else
	     {\color{orange}LUYỆN TẬP \thelt\ (\MakeLowercase{#1})}
	   \fi
	},
	boxed title style={empty,arc=0pt,outer arc=0pt,boxrule=0pt},
   frame empty,
   frame hidden,	
   sharp corners,
   interior style={draw=none,fill=none},
   left=3pt,right=0pt,top=0pt,bottom=0pt,
    overlay unbroken and first ={
	    \draw[line width=1pt,orange] (frame.north west)--(frame.south west);
	    \fill[orange] (frame.north west) circle (1.5pt);
	    \fill[orange] (frame.south west) circle (1.5pt);
    },
     overlay middle and last={
	    \draw[line width=1pt,orange] (frame.north west)--(frame.south west);
	    \fill[orange] (frame.north west) circle (1.5pt) 
	    	+(0,6pt) node[anchor=west,text=green!75!black,font=\small\bfseries]{Tiếp theo luyện tập \thelt};
	    \fill[orange] (frame.south west) circle (1.5pt);
    }   
}

\AtBeginEnvironment{lt}{%
	\def\maudapan{white}%
	\renewcommand{\loigiai}[1]{}%
}

\newenvironment{luyentap}{\begin{lt}}{\end{lt}}
%%%========================Hộp vận dụng========================%%%
\newtcolorbox{vandung}{enhanced,breakable,
	attach boxed title to top left={xshift=-4mm,yshift=-0.5mm},
	fonttitle=\bfseries\sffamily,
	title={
	     {\color{mauchuong}VẬN DỤNG \phantom{A}}
	},
	boxed title style={empty,arc=0pt,outer arc=0pt,boxrule=0pt},
   frame empty,
   frame hidden,	
   sharp corners,
   interior style={draw=none,fill=none},
   left=3pt,right=0pt,top=0pt,bottom=0pt,
   overlay unbroken and first ={
	    \draw[line width=1pt,mauchuong] (frame.north west)--(frame.south west);
	    \fill[mauchuong] (frame.north west) circle (1.5pt);
	    \fill[mauchuong] (frame.south west) circle (1.5pt);
    },
     overlay middle and last={
	    \draw[line width=1pt,mauchuong] (frame.north west)--(frame.south west);
	    \fill[mauchuong] (frame.north west) circle (1.5pt) 
	    	+(0,6pt) node[anchor=west,text=green!75!black,font=\small\bfseries]{Vận dụng tiếp theo};
	    \fill[mauchuong] (frame.south west) circle (1.5pt);
    }   
}

%%%================Hộp định nghĩa====================%%%
\definecolor{dnxanh}{HTML}{00CC99}
\definecolor{dndo}{HTML}{CC0033}
\newcounter{dn}
%%%===================================================%%%
\NewTColorBox[use counter=dn]{dn}{+O{}}{
enhanced,
colback=white,
fonttitle=\bfseries\sffamily\color{white},
 title={\faGg\ Định nghĩa~\thetcbcounter aa},
colframe=mauchuong,
colbacktitle=white,
fonttitle=\bfseries,
attach boxed title to top left=
{yshift=-2mm-\tcboxedtitleheight,yshifttext=2mm-\tcboxedtitleheight/2},
 boxed title style={
 outer arc=0pt,
 arc=0pt,
 top=3pt,
 bottom=3pt,
 left=0pt,
 right=0pt,
 frame hidden,
 },
 overlay unbroken and first={
 \path
 ($ (title.north west) +(-2pt,0pt)$) coordinate (A)
 ($ (title.south west) +(-2pt,3pt)$) coordinate (B)
 ($ (title.south east)+(3pt,3pt) $) coordinate (C)
 ($ (title.north east)+(3pt,0) $) coordinate (D)
 (frame.north east) coordinate (E)
 (frame.south east) coordinate (F)
 (intersection of A--D and E--F) coordinate (At)
 (intersection of B--C and E--F) coordinate (Bt)
 ;
 
\path[fill=mauchuong!25,rounded corners=3pt](A)--(B)--($(Bt)!5pt!(B)$)--($(At)!5pt!(A)$)--cycle;
\path[fill=mauchuong!70!black,rounded corners=2pt](A)--(B)--([xshift=3pt]C)--([xshift=3pt]D)--cycle;   
\path[left color=dndo,right color=dndo!80,rounded corners=3pt]([xshift=-2pt]A)--([xshift=-2pt]B)--(C)--(D)--cycle;
 
 \path ($ (A)!0.5!(B) +(9pt,0)$) node[anchor=west,font=\color{white}\bfseries\sffamily\selectfont]{{\faGg\ Định nghĩa~\thetcbcounter}};
 
 \path ($ (C)!0.5!(D) +(5pt,0)$) node[anchor=west,font=\color{mauchuong}\small\bfseries]{#1};
},
top=\tcboxedtitleheight
}

%%%===================================================%%%
\newcounter{md}[section]
\NewTColorBox[use counter=md]{md}{+O{}}{
enhanced,
colback=white,
fonttitle=\bfseries\sffamily\color{white},
 title={\faGg\ Định lý~\thetcbcounter aa},
colframe=mauchuong,
colbacktitle=white,
fonttitle=\bfseries,
attach boxed title to top left=
{yshift=-2mm-\tcboxedtitleheight,yshifttext=2mm-\tcboxedtitleheight/2},
 boxed title style={
 outer arc=0pt,
 arc=0pt,
 top=3pt,
 bottom=3pt,
 left=0pt,
 right=0pt,
 frame hidden,
 },
 overlay unbroken and first={
 \path
 ($ (title.north west) +(-2pt,0pt)$) coordinate (A)
 ($ (title.south west) +(-2pt,3pt)$) coordinate (B)
 ($ (title.south east)+(3pt,3pt) $) coordinate (C)
 ($ (title.north east)+(3pt,0) $) coordinate (D)
 (frame.north east) coordinate (E)
 (frame.south east) coordinate (F)
 (intersection of A--D and E--F) coordinate (At)
 (intersection of B--C and E--F) coordinate (Bt)
 ;
 
\path[fill=mauchuong!25,rounded corners=3pt](A)--(B)--($(Bt)!5pt!(B)$)--($(At)!5pt!(A)$)--cycle;
\path[fill=mauchuong!70!black,rounded corners=2pt](A)--(B)--([xshift=3pt]C)--([xshift=3pt]D)--cycle;   
\path[left color=dndo,right color=dndo!80,rounded corners=3pt]([xshift=-2pt]A)--([xshift=-2pt]B)--(C)--(D)--cycle;
 
 \path ($ (A)!0.5!(B) +(9pt,0)$) node[anchor=west,font=\color{white}\bfseries\sffamily\selectfont]{{\faGg\ Định lý~\thetcbcounter}};
 
 \path ($ (C)!0.5!(D) +(5pt,0)$) node[anchor=west,font=\color{mauchuong}\small\bfseries]{#1};
},
top=\tcboxedtitleheight
}
%%%===============Hộp kiến thức trọng tâm===============%%%
\newtcolorbox{kt}[1][]{enhanced,breakable, notitle,
   frame empty,
   colframe=white,
   colback=cyan!10,
   sharp corners,
   rounded corners=southeast,
   arc=15pt,
   notitle,
   interior style={left color=cyan!15,right color=cyan!10},
   left=0.65cm,right=6pt,top=6pt,bottom=6pt,before skip=12pt,
    overlay unbroken and first ={
	    \draw[line width=1pt,teal] ([shift={(-1pt,-1pt)}]frame.north east) coordinate (A)--
	    ($([yshift=-1pt]frame.north west)+(0.75,0)+(1pt,0)$) coordinate (B) arc (0:-90:0.75) coordinate (C)
	    --([shift={(1pt,1pt)}]frame.south west) coordinate (D);
	    \foreach \ten in {A,B,C,D}{\fill[teal] (\ten) circle (1.75pt);}
		\path ($(B)!0.5!(C)+(135:0.4)$) node{\includegraphics[width=1cm]{Khaibao/Icon/Bulb.pdf}};
		\ifblank{#1}{}{%
			\path ($(B)+(0,-2pt)$) node[font=\footnotesize\bfseries,text=teal,anchor=south west]{#1};
		}
    },
    overlay middle and last={
	    \draw[line width=1pt,teal] ([shift={(-1pt,-1pt)}]frame.north east) coordinate (A)--
	    ($([yshift=-1pt]frame.north west)+(0.75,0)+(1pt,0)$) coordinate (B) arc (0:-90:0.75) coordinate (C)
	    --([shift={(1pt,1pt)}]frame.south west) coordinate (D);
	    \foreach \ten in {A,B,C,D}{\fill[teal] (\ten) circle (1.75pt);}
		\path ($(B)!0.5!(C)+(135:0.4)$) node{\includegraphics[width=1cm]{Khaibao/Icon/Bulb.pdf}};
    } 
}
 
%%%===============Hộp lưu ý===============%%%
\newenvironment{luuy}{\par\noindent%
\begin{tikzpicture}
\begin{scope}[rotate=-45,scale=0.5]
	\draw[line width=1.25pt,magenta] (0,0)--+(0,-0.25) arc (180:360:0.1)--+(0,0.35) arc(0:180:0.115)--+(0,-0.5) arc (180:360:0.13)--+(0,0.4);
\end{scope}
\end{tikzpicture}%
\textcolor{magenta}{\textsf{\textbf{Lưu ý:}}}%
}{}

%%%===============Hộp lưu ý===============%%%
\newenvironment{nx}{\par\noindent%
%\begin{tikzpicture}[anchor=base,baseline=-7pt]
%\begin{scope}[rotate=-45,scale=0.65]
%	\draw[line width=1.25pt,magenta] (0,0)--+(0,-0.25) arc (180:360:0.1)--+(0,0.35) arc(0:180:0.115)--+(0,-0.5) arc (180:360:0.13)--+(0,0.4);
%\end{scope}
%\path (0.25,-0.15) node[font=\bfseries,anchor=west,outer sep=0pt,magenta]{Nhận xét:};%
%\end{tikzpicture}%
%\hspace*{-6pt}
\textbf{\textcolor{magenta!75!red}{\faSpinner\ Nhận xét:}}%
}{}


\newenvironment{note}{\begin{nx}}{\end{nx}}

\usepackage{marvosym}
 \definecolor{colbox}{rgb}{0.34, 0.7, 1}
 \definecolor{mauhophd}{HTML}{E1F0EA}
 \definecolor{xanhteal}{HTML}{2A5F49}
 \newcounter{hd}
%%%=======================Hộp Hoạt động 1==================%%%
\newtcolorbox{hoatdong}{ 
enhanced, colback=mauhophd, frame hidden,
   code={\refstepcounter{hd}},
    underlay vignette={draw method=clipped, size=2.5mm, semi fade in=gray, east size=1.25mm, south size=1.25mm},
   top=20pt,left=9pt,right=6pt,bottom=5pt,before skip=12pt,
     attach boxed title to top left={xshift=-2mm,yshift=-15pt},
     fonttitle=\fontfamily{ugq}\bfseries\selectfont,
   	title={
   	     {\Coffeecup\ HOẠT ĐỘNG \thehd}
   	},   
       overlay unbroken and first ={
   	    \fill[blur shadow={shadow blur extra rounding,
   	    shadow blur steps=15,
   	    shadow xshift=0.15ex,shadow yshift=-0.15ex,
   	    shadow blur radius=0.75ex},
   	    rounded corners=5pt,
   	    white] ([shift={(145:5pt)}]title.north west) coordinate (A)
   	    rectangle 
   	    ([shift={(-35:5pt)}]title.south east) coordinate (B);
   	    \path ($(A)!0.5!(B)$) node[font=\fontfamily{ugq}\bfseries\selectfont,text=xanhteal]{\Coffeecup\ HOẠT ĐỘNG \thehd};
   	    \draw[rounded corners,teal!65] (A) rectangle (B);
       }
  }
%%%=======================Hộp Hoạt động 2==================%%%
\newtcolorbox{Hhoatdong}{ 
step=hd,  
breakable,
enhanced,
frame hidden, % hide frame
interior style={left color=colbox!10,right color=yellow!2},
toprule=0pt,rightrule=0mm,leftrule=0mm,bottomrule=0mm,
top=3pt, left=3pt, right=6pt,  
bottom=5pt,right skip=0.25cm,before skip = 12pt,after skip=3pt,
	fonttitle=\large\bfseries,
	title={\color{yellow}{\normalsize \Coffeecup} \fontsize{12pt}{0pt}\fontfamily{ugq}\selectfont HOẠT ĐỘNG \thehd},
		attach boxed title to top left,
		boxed title style={
			enhanced,
			colframe=red!70!yellow,
			colback=red!70!yellow,
			arc=5pt,
			bottomrule=0pt,
			leftrule=0pt
		},
	overlay={
		\draw[draw=colbox,line width=1.2pt,rounded corners=2mm]
		([xshift=4mm,yshift=0.5pt]title.north west)--([xshift=3mm,yshift=0.5pt]title.north east)--
		([xshift=3mm]title.south east)--(title.south west)|-
		(interior.south east)-|([yshift=-0.25cm]interior.north east) coordinate (Dt);
		\path (title.north east) coordinate (A)
			(title.south east) coordinate (C)
			($ (A)!0.5!(C) $) coordinate (Bt)
			([rotate around={90:(Bt)}]C) coordinate (B);
			\foreach[count=\j ] \x/\y in {10/0.85,15/0.75,20/0.65,25/0.5,30/0.4}{
			\draw[transform canvas={xshift=\x pt},line width=2pt,rounded corners=3pt,colbox,opacity=\y,
				shorten >={0.35*\j pt},shorten <={0.35*\j pt},line cap=round] (A)--(B)--(C);
			},
		\fill[colbox] (Dt) circle (1.75pt);
		\path ($(title.north)!0.5!(title.south)+(-0.65pt,-0.25pt)$) node[font=\large\bfseries,text=colbox]{{\normalsize \Coffeecup} \fontsize{12pt}{0pt}\fontfamily{ugq}\selectfont HOẠT ĐỘNG \thehd};
		} 
}

%%%===============Hoạt động nhóm===============%%%
\newenvironment{nhom}[1][Thảo luận]{\par\noindent%
\begin{tikzpicture}[anchor=base,baseline]
	\path (-45:1pt) node[scale=1.5,text=gray!25,inner sep=0pt,outer sep=0pt]{\faWechat};
	\path (0,0) node[scale=1.5,text=red!65!white,inner sep=0pt,outer sep=0pt](A){\faWechat};
	\path ($(A.east)+(-40:1pt)$) node[anchor=west,text=gray!35,font=\sffamily\bfseries\large]{#1};
	\path (A.east) node[anchor=west,text=red!65!white,font=\sffamily\bfseries\large]{#1};
\end{tikzpicture}%
\par
}{}


%%%==================em có biết======================%%%
\newtcolorbox{emcobiet}{%
	enhanced,breakable,
	attach boxed title to top left={xshift=3mm,yshift=-0.5mm},
	fonttitle=\bfseries\sffamily,
	title={\color{mauchuong}EM CÓ BIẾT},
	boxed title style={empty,arc=0pt,outer arc=0pt,boxrule=0pt},
   frame empty,
   frame hidden,	
   arc=7pt,
   interior style={fill=mauchuong!7,rounded corners},
   left=5pt,right=3pt,top=6pt,bottom=3pt,
   overlay unbroken and first={
 	  \path ($(title.west)+(1pt,-1pt)$) node{
     	 \begin{tikzpicture}[line cap=round,yscale=0.7,xscale=.85]
     	 	\begin{scope}[color=mauchuong]
     	 		\draw[line width=1.2pt] (-0.13,2pt) to[out=85,in=-90]+(0.05,0.1) arc(0:190:0.05);
     	 		\fill (-0.13,0) circle (0.8pt);
     	 	\end{scope}
     	 	\begin{scope}[color=mauchuong,opacity=0.4]
     	 		\draw[line width=2pt] (0,3pt) to[out=85,in=-90]+(0.1,0.16) arc(0:190:0.1);
     	 		\fill (0,0.25pt) circle (1.1pt);
     	 	\end{scope}
     	 	\end{tikzpicture}    
    };
   }
}

%%%==============================================%%%
\newtcolorbox{tomtat}[1][TÓM TẮT KIẾN THỨC]{
enhanced,breakable,
	attach boxed title to top left={xshift=9pt,yshift=-0.5mm},
	fonttitle=\bfseries\sffamily,
	title={\color{mausubsection}\MakeUppercase{#1}},
	boxed title style={empty,arc=0pt,outer arc=0pt,boxrule=0pt},
   frame empty,
   frame hidden,	
   arc=7pt,
   interior style={draw=none,fill=none},
   left=3pt,right=3pt,top=6pt,bottom=3pt,
   left skip=0.35cm, before skip=0.65cm,
   overlay unbroken and first={
   	\path ($(frame.south west)+(0,9pt)$) coordinate (A)--($(frame.north west)+(0,-3pt)$) coordinate (B)
   	    --+(0,0.35) coordinate (C)--
   	    ([turn]120:0.75) coordinate (D)--+(9,0) coordinate (E)
   	    ($(frame.south east)+(-9pt,0)$) coordinate (Bp)
   	    ($(frame.north east)+(-9pt,0)$) coordinate (Ap)
   	    ($(frame.south west)+(-9pt,0)$) coordinate (Cp)
   	    (intersection of D--E and Bp--Ap) coordinate (Ep) ;
   	\draw[thick,mauchuong] (A)--($(B)!7pt!(A)$)--($(B)!12pt!(Ep)$)--(Ep)
   		($(Cp)!0.55!(Bp)$) coordinate (Ad)--($(Bp)!17pt!(Cp)+(0.5,0)$) coordinate (Dd)--([turn]45:0.5) coordinate (Bd)
   		--($(Ap)!0.35!(Bp)+(6pt,0)$) coordinate (Cd)
   		($(Dd)+(0.12,0)$)--+(45:0.35);
   	\path ($(B)+(16pt,12pt)$) node[anchor=east,outer sep=0pt]{\includegraphics[width=1.1cm]{Khaibao/Icon/KienThuc.pdf}};
   	\foreach \diem in {A,Ep,Ad,Cd}{\fill[mauchuong] (\diem) circle (1.5pt);}
    },
      overlay middle and last={
      	\path ($(frame.south west)+(0,9pt)$) coordinate (A)--($(frame.north west)+(0,-3pt)$) coordinate (B)
      	    --+(0,0.35) coordinate (C)--
      	    ([turn]120:0.75) coordinate (D)--+(9,0) coordinate (E)
      	    ($(frame.south east)+(-9pt,0)$) coordinate (Bp)
      	    ($(frame.north east)+(-9pt,0)$) coordinate (Ap)
      	    ($(frame.south west)+(-9pt,0)$) coordinate (Cp)
      	    (intersection of D--E and Bp--Ap) coordinate (Ep) ;
      	\draw[thick,mauchuong] (A)--($(B)!7pt!(A)$)--($(B)!12pt!(Ep)$)--(Ep)
      		($(Cp)!0.55!(Bp)$) coordinate (Ad)--($(Bp)!17pt!(Cp)+(0.5,0)$) coordinate (Dd)--([turn]45:0.5) coordinate (Bd)
      		--($(Ap)!0.35!(Bp)+(6pt,0)$) coordinate (Cd)
      		($(Dd)+(0.12,0)$)--+(45:0.35);
      	\path ($(B)+(14pt,10pt)$) node[anchor=east,outer sep=0pt]{\includegraphics[width=1cm]{Khaibao/Icon/KienThuc.pdf}};
      	\foreach \diem in {A,Ep,Ad,Cd}{\fill[mauchuong] (\diem) circle (1.5pt);}
       }
}

%%%==================môi trường dạng=====================%%%
\newcounter{dang}
\newtcolorbox{dang}[1]{
fonttitle=\sffamily\bfseries,%fontupper=\itshape,
breakable,
enhanced jigsaw,
colbacktitle=mauchuong,
coltitle=white,
colframe=mauchuong,
colback=white,
sharp corners,
rounded corners=south,
borderline north={1pt}{0pt}{mausection},
drop fuzzy shadow=gray!75,
before skip=3mm,after skip=3mm,
left=2mm,right=2mm,top=3mm,bottom=2mm,
title style={left color=mauchuong,right color=mauchuong!65},
underlay unbroken and first={
\def\dai{0.75}
\pgfmathsetmacro{\daih}{\dai/cos(20)}
\def\cao{0.675}
\fill[mauchuong!85!black,rounded corners=1pt,drop shadow={shadow xshift=1pt, shadow yshift=-1pt,fill=gray!65}]
([shift={(2*\dai,5pt)}]title.north west)
coordinate (A1)
--+(2*\dai,0) coordinate (A2)
--([turn]-90:\cao) coordinate (A3)
--([turn]-70:\daih)
--([turn]-40:\daih) coordinate (A4)
--cycle;
\begin{pgfonlayer}{background}
\fill[gray!65] (A1)--++(225:0.5)--+(1,0)--cycle;
\end{pgfonlayer}
\node[circle,draw=white,fill=mausection,line width=1pt, inner sep=3pt,text=white,font=\sffamily\bfseries] at ($(A1)!0.5!(A3)+(0,-5pt)$){\thedang};
\path([yshift=-6pt]A4) node[font=\sffamily\bfseries,text=white,anchor=south east]{Dạng};
},
overlay middle and last={
\def\dai{0.75}
\pgfmathsetmacro{\daih}{\dai/cos(20)}
\def\cao{0.5}
\fill[left color=mauchuong,right color=mauchuong!65,draw=mauchuong,line width=1.4pt] ($(frame.north west)+(0.7pt,13pt)$) rectangle ($(frame.north east)+(-0.7pt,-7pt)$);
\fill[mauchuong!85!black,rounded corners=1pt,drop shadow={shadow xshift=1pt, shadow yshift=-1pt,fill=gray!65}]
([shift={(2*\dai,17pt)}]frame.north west)
coordinate (A1)
--+(2*\dai,0) coordinate (A2)
--([turn]-90:\cao) coordinate (A3)
--([turn]-70:\daih)
--([turn]-40:\daih) coordinate (A4)
--cycle;
\begin{pgfonlayer}{background}
\fill[gray!65] (A1)--++(225:0.5)--+(1,0)--cycle;
\end{pgfonlayer}
\node[circle,draw=white,fill=mausection,line width=1pt, inner sep=3pt,text=white,font=\sffamily\bfseries] at ($(A1)!0.5!(A3)+(0,-5pt)$){\thedang};
\path([yshift=-7pt]A4) node[font=\sffamily\bfseries,text=white,anchor=south east](A5){Dạng};
\path($(A5)+(2.2,0)$) node[font=\sffamily\bfseries,text=white,anchor=west]{Phần tiếp theo};
},
title={\hspace*{2.9cm}\parbox{\textwidth-3cm}{#1}\refstepcounter{dang}},%
}
\tikzset{laphuong/.pic={
	\path[pic actions,line join=round]
	(0,0)--++(27.5:r)--++(152.5:r)--+(-152.5:r)--cycle
	(0,0)--++(0,-r)--++(152.5:r)--+(0,r)--cycle
	(0,0)--++(0,-r)--++(27.5:r)--+(0,r)--cycle;
}}
\usetikzlibrary{matrix}
\newtcolorbox{kttrongtam}{
enhanced,breakable,
before skip=6pt,after skip=6pt,
boxsep=9pt,
left=1cm,
right=0.25cm,
bottom=0pt,
frame empty,
frame hidden,
arc=7pt,
colback=orange!7,
sharp corners,
rounded corners=south,
colframe=orange!35,
boxrule=0pt,
borderline north={9pt}{0pt}{orange!65},
drop fuzzy shadow=gray!75,
overlay unbroken and first ={
\fill[orange,drop shadow={opacity=0.25,shadow xshift=1pt, shadow yshift=-1.5pt,shadow scale=1.05}] (frame.north west) coordinate (A)--+(2,0)--($(A)+(0,-1.25)$)--cycle;
\path ($(A)+(0.65,-0.15)$) coordinate (B) node{
\begin{tikzpicture}[x=3.65pt,y=3.65pt,xscale=-1]
\draw[thick,orange,fill=white] (-1,-4) to[out=50,in=-90,looseness=1.4] (-1,-3) to[out=-160,in=0]
(-2,-3.2) to[out=180,in=-120] (-2.7,-2) to[out=120,in=-120] (-2.9,-1.5)
to[out=120,in=-120] (-2.95,-1.1) to[out=60,in=-30] (-3.2,-0.9)
to[out=150,in=-120] (-3,0.4) to[out=60,in=-90] (-3,1.2)
to[out=90,in=90,looseness=1.8] (3.3,1.2)
to[out=-90,in=90,looseness=0.8] (2,-2)
to[out=-90,in=150,looseness=0.8] (3.2,-4) -- cycle;
\end{tikzpicture}
};
\tikzset{declare function={r=0.275;}}
\path (B) pic[fill=orange,draw=white]{laphuong};
}
}

\newenvironment{trithuc}{\begin{kt}}{\end{kt}}
%\newenvironment{tomtath}{\begin{kttrongtam}}{\end{kttrongtam}}


%%%=======================================================%%%
\newtcolorbox{phantich}{
   		   boxsep=3pt,
   		   left=3pt,
   		   top=9pt,
   		   enhanced,
   		   right=3pt,
   		   boxrule=0.65pt, %độ dày của khung
   		   colframe=cyan!15, %màu của khung
   		   grow to left by=-3pt,
   		   before skip=5pt,
		   interior style={left color=cyan!7,right color=white},
		   frame empty,
%   		   code={\refstepcounter{ex}},
   		   title={\faGg\ Phân tích}, %nhãn của hộp
   		   coltitle=white, %màu của nhãn của hộp	
   		   fonttitle=\bfseries\sffamily,
		   attach boxed title to top left={xshift=-7pt,yshift=-\tcboxedtitleheight/2},		  
   		   underlay boxed title={
   					\path[fill=black] (title.south west)--++(7.1pt,0)--+(0,-7.1pt)--cycle;
   				},
   		 sharp corners,
   		 underlay={
   			\path[even odd rule,left color=teal,right color=orange] {($(frame.north west)+(-0.65pt,0.65pt)$) rectangle  ($(frame.south east)+(0.65pt,-0.65pt)$)}
   			{[rounded corners=3pt]($(frame.north west)+(0.65pt,-0.65pt)$) rectangle  ($(frame.south east)+(-0.65pt,0.65pt)$)};
			\path (title.north west) coordinate (A) (title.south west) coordinate (B)
					(title.north east) coordinate (C);   				
 			\path[left color=red,right color=orange] let \p1=($(A)-(B)$),  	
 					\n1={veclen(\p1)}
 					in (A) -- (C) arc (90:-90:{\n1/2})-- (B) -- cycle
 					($(B)!0.5!(C)+(0.65pt,-0.65pt)$) node[font=\bfseries\sffamily,text=gray]{\faGg\ Phân tích}
 					($(B)!0.5!(C)$) node{\tcbtitle}; 
   		},		
    }
    
  %%%============================Tính chất=========================%%%
  \newtcolorbox{hopgoc}[1][Định lý]{enhanced,breakable,notitle,
     frame empty,
     colframe=white,
     colback=cyan!10,
     sharp corners,
     rounded corners=southeast,
     arc=15pt,
     interior style={left color=cyan!15,right color=cyan!10},
     left=0.65cm,right=6pt,top=6pt,bottom=6pt,before skip=19pt,
      overlay unbroken and first ={
  	    \draw[line width=1pt,teal] ([shift={(-1pt,-1pt)}]frame.north east) coordinate (A)--
  	    ($([yshift=-1pt]frame.north west)+(0.75,0)+(1pt,0)$) coordinate (B) arc (0:-90:0.75) coordinate (C)
  	    --([shift={(1pt,1pt)}]frame.south west) coordinate (D);
  	    \foreach \ten in {A,B,C,D}{\fill[teal] (\ten) circle (1.75pt);}
  		\path ($(B)!0.5!(C)+(135:0.4)$) node{\includegraphics[width=1cm]{Khaibao/Icon/Bulb.pdf}};
  		\path ($(B)+(0,-2pt)$) node[font=\sffamily\bfseries,text=teal,anchor=south west]{#1};
      },
      overlay middle and last={
  	    \draw[line width=1pt,teal] ([shift={(-1pt,-1pt)}]frame.north east) coordinate (A)--
  	    ($([yshift=-1pt]frame.north west)+(0.75,0)+(1pt,0)$) coordinate (B) arc (0:-90:0.75) coordinate (C)
  	    --([shift={(1pt,1pt)}]frame.south west) coordinate (D);
  	    \foreach \ten in {A,B,C,D}{\fill[teal] (\ten) circle (1.75pt);}
  		\path ($(B)!0.5!(C)+(135:0.4)$) node{\includegraphics[width=1cm]{Khaibao/Icon/Bulb.pdf}};
      } 
  }
  
  \newcounter{tc}[section]
  \newenvironment{tc}{\refstepcounter{tc}\begin{hopgoc}[Tính chất \thetc]}{\end{hopgoc}}
  
  \newcounter{hq}[section]
  \newenvironment{hq}{\refstepcounter{hq}\begin{hopgoc}[Hệ quả \thehq]}{\end{hopgoc}}
 
\usepackage{ex_test}
\setlist[enumerate,1]{label*=\circlenum{\arabic*},itemsep=-1pt,topsep=0pt}
\setlist[enumerate,2]{itemsep=0pt,topsep=0pt,label=\circlenumH{\alph*},wide=0cm,leftmargin=!}
\setlist[enumerate,3]{itemsep=0pt,topsep=0pt,label=\protect{\textcolor{teal}{\roman*)}},wide=1cm,leftmargin=1cm}
\setlist[itemize,1]{itemsep=0pt,topsep=0pt,label=\protect{\small\color{teal}\rhombusdot}}
\setlist[itemize,2]{itemsep=0pt,topsep=0pt,label=\protect{\small\color{teal}\starletdot},wide=0.35cm,
leftmargin=0.35cm}
\setlist[itemize,3]{itemsep=0pt,topsep=0pt,label=\protect{\small\color{teal}\faAngellist},wide=0.5cm,leftmargin=.5cm}
%%%================Lệnh immini================%%%
\def\loaicau{}
\def\tieudehinh{}
\def\loaicau{}
\newlength{\widthPB}
\newbool{TNdungsai}
  
\newcommand{\hinhphai}[2]{%
\savebox{\imbox}{#2}%
\tcbsidebyside[
sidebyside adapt=right,
blanker,sidebyside gap=5mm,
sidebyside align=top seam,
]{%
	{\tieudehinh}%
		#1
	}{%
		\usebox{\imbox}%
	}%
%	\xdef\loaicau{\setboolean{TNdungsai}{false}}%
}

\renewcommand{\immini}[3][]{%
\savebox{\imbox}{#3}%
%\setlength{\widthPB}{0.53\linewidth}
\tcbsidebyside[
	sidebyside adapt=right,
	blanker,sidebyside gap=5mm,
	sidebyside align=top seam,before skip=0pt
	]{%
		\tieudehinh%
	#2
	}{%
		\usebox{\imbox}
	}%
}

\tikzset{daudung/.pic={\draw[line width=0.65pt] (-0.7ex,0.3ex) -- ++(0.6ex,-0.5ex) --++(1.1ex,1.7ex) --++(-1.05ex,-2ex) -- cycle;}}
\tikzset{dausai/.pic={\draw[line width=0.65pt] (50:0.5ex) -- (230:0.5ex) (130:0.5ex) -- (310:0.5ex);}}

 \newcommand{\TLdung}[1]{%
 \tikz[baseline=(char.base)]
  {\node[shape=circle,inner sep=0.5pt,draw=mausubsection,font=\bfseries\footnotesize,minimum size=12pt,outer sep=5pt] (char) {Đ};
  }\,\ignorespaces}
 
 \newcommand{\TLsai}[1]{%
 \tikz[baseline=(char.base)]
 {\node[shape=circle,inner sep=0.5pt,draw=mausubsection,font=\bfseries\footnotesize,minimum size=12pt,outer sep=5pt] (char) {S};
 }\,\ignorespaces}
 
\def\checkmarkH{{\tikz[line join=miter]{\draw (0.5ex,0.5ex) coordinate (O) circle (0.75ex);\path[preaction={draw=white,line width=0.2ex},fill=white,draw] (0,.5ex) -- (.5ex,0) -- (1.5ex,1.35ex) -- (.45ex,.25ex) -- cycle; \draw ([shift={(90:0.75ex)}]O) arc (90:270:0.75ex);}} }

\makeatletter
\RenewDocumentEnvironment{ex}{+!O{}O{}}{%
	\global\setbool{TNdungsai}{false}
	\ifthenelse{\equal{#1}{c}}{
		\ifblank{#2}{%
			\@ifnextchar\immini{\gdef\tieudehinh{{\bfseries\color{mauchuong}\checkmarkH Câu \theex. }}}%
			{{\par\noindent{\bfseries\color{mauchuong}\checkmarkH Câu \theex.}
			}}%
		 }%blank#2true
		 {\@ifnextchar\immini{\gdef\tieudehinh{{\bfseries\checkmarkH Câu \theex } {(\textit{#1}).\,}}}%
			{{\par\vspace*{-6pt}\noindent\bfseries\color{mauchuong}\checkmarkH Câu \theex} {\color{mauchuong}(\textit{#2}).\,}}%
		 }}%blank#2false	
	{%
	\ifblank{#2}{%
		\@ifnextchar\immini{\gdef\tieudehinh{{\bfseries\color{mauchuong} Câu \theex. }}}%
		{{\par\noindent{\bfseries\color{mauchuong} Câu \theex.}
		}}%
	 }%blank#2true
	 {\@ifnextchar\immini{\gdef\tieudehinh{{\bfseries Câu \theex } {(\textit{#1}).\,}}}%
		{{\par\vspace*{-6pt}\noindent\bfseries\color{mauchuong} Câu \theex} {\color{mauchuong}(\textit{#2}).\,}}%
	 }}%blank#2false	
	}{\gdef\tieudehinh{}\ignorespacesafterend}
\makeatother

\AtBeginEnvironment{ex}{%
	\refstepcounter{ex}%
	 \setcounter{numTrue}{0}%
	 \setcounter{numTrueSol}{0}%
	 \global\setbool{TNdungsai}{false}%
	 \def\dotEX{.}%
	 \setcounter{numChoice}{0}%
	 \renewcommand{\loigiai}[1]{\gdef\tieudehinh{}\par\noindent%
	 {\color{mausubsection}\reflectbox{\Large\WritingHand}\ {Lời giải.}} #1
	 \ifthenelse{\value{numTrueSol}>0}{
	 \signed{\textcolor{mauchuong}{\faKey~\circlenum{\Alph{numTrueSol}}}}
	 }{}
 }%}
} 

\renewcommand{\FalseEX}{\stepcounter{dapan}{{\Alph{dapan}}.}}
\renewcommand{\TrueEX}{\stepcounter{dapan}{\Alph{dapan}}.}

\RenewEnviron{Solution}[1]{%
	\tikz{\node[draw=mauchuong,thick,inner sep=3pt,text width=0.1\linewidth-4\fboxsep,minimum height=0.65cm,
		rounded corners,fill=yellow!5,align=center](A){
		{\resizebox{1.1cm}{!}{\textbf{#1} -	{\textbf{\BODY}}}}
		};
	 \path ($(A)+(0,-13pt)$);
	}
}

\def\bangdapan#1{
\begin{center}
	{\bfseries\sffamily BẢNG ĐÁP ÁN TRẮC NGHIỆM}
	\vspace*{-11pt}
\end{center}

\begin{flushleft}
	\input{Ans/#1}
\end{flushleft}
}

%%%===============trắc nghiệm đúng sai======================%%%
\usepackage{makecell}
\newbool{Ads}
 \newbool{Bds}
 \newbool{Cds}
 \newbool{Dds}
 \makeatletter
 \newcommand{\KiemtraA}{\@ifnextchar\True{\global\setbool{Ads}{true}\xdef\saveans{\saveans\string\TLdung{A}}}{\global\setbool{Ads}{false}\xdef\saveans{\saveans\string\TLsai{A}}}}
 \newcommand{\KiemtraB}{\@ifnextchar\True{\global\setbool{Bds}{true}\xdef\saveans{\saveans\string\TLdung{B}}}{\global\setbool{Bds}{false}\xdef\saveans{\saveans\string\TLsai{B}}}}
 \newcommand{\KiemtraC}{\@ifnextchar\True{\global\setbool{Cds}{true}\xdef\saveans{\saveans\string\TLdung{C}}}{\global\setbool{Cds}{false}\xdef\saveans{\saveans\string\TLsai{C}}}}
 \newcommand{\KiemtraD}{\@ifnextchar\True{\global\setbool{Dds}{true}\xdef\saveans{\saveans\string\TLdung{D}}}{\global\setbool{Dds}{false}\xdef\saveans{\saveans\string\TLsai{D}}}}
 \makeatother
 
 
 \def\saveans{}
 \newlength{\widthNPB}
 \setlength{\widthPB}{1.07\linewidth}
 \setlength{\widthNPB}{0.53\linewidth-1.1cm}
 \setlength\LTpre{5pt}% khoảng cách trước longtable
 \setlength\LTpost{3pt}
 \def\maudapan{black}
 
\newcommand{\choiceTF}[5][1]{%
\def\saveans{}%
\global\setbool{TNdungsai}{true}%
\loigiaiTF
\ifthenelse{\equal{#1}{1}}{%
 \begin{enumerate}[wide=\parindent,label*=\protect{\alph*)},itemsep=-3pt,topsep=-4pt]
 \item \KiemtraA#2\dotEX	
 \item \KiemtraB#3\dotEX
 \item \KiemtraC#4\dotEX
 \item \KiemtraD#5\dotEX
 \end{enumerate}
 }{
 \ifthenelse{\equal{#1}{t}}{%
 \begin{longtable}{|p{\widthPB}|C{0.45cm}|C{.45cm}|}
 \hline
 \thead{\normalsize\sffamily Phát biểu}& {\sffamily Đ} & {\sffamily S}\\
 \hline
 	{a)} \KiemtraA#2\dotEX &\ifbool{Ads}{{\color{\maudapan}\footnotesize\faCheck}}{} &\ifbool{Ads}{}{{\color{\maudapan}\footnotesize\faClose}}  \\
 \hline
	{b)} \KiemtraB#3\dotEX & \ifbool{Bds}{{\color{\maudapan}\footnotesize\faCheck}}{} & \ifbool{Bds}{}{{\color{\maudapan}\footnotesize\faClose}}\\
 \hline
	{c)} \KiemtraC#4\dotEX &\ifbool{Cds}{{\color{\maudapan}\footnotesize\faCheck}}{} &\ifbool{Cds}{}{{\color{\maudapan}\footnotesize\faClose}} \\
 \hline
	{d)} \KiemtraD#5\dotEX & \ifbool{Dds}{{\color{\maudapan}\footnotesize\faCheck}}{} & \ifbool{Dds}{}{{\color{\maudapan}\footnotesize\faClose}} \\
 \hline
 \end{longtable}
 }{%
 \ifthenelse{\equal{#1}{tt}}{%
 \begin{center}
 \vspace*{-0.65\baselineskip}
 \begin{tabular}{cc}%
 \begin{tabular}[t]{|p{\widthNPB}|C{0.45cm}|C{.45cm}|}
		 \hline
		 \thead{\normalsize\sffamily Phát biểu}& {\sffamily Đ} & {\sffamily S}\\
		 \hline
		 {a)} \KiemtraA#2\dotEX &\ifbool{Ads}{{\color{\maudapan}\footnotesize\faCheck}}{} &\ifbool{Ads}{}{{\color{\maudapan}\footnotesize\faClose}}  \\
		 \hline
		 {b)} \KiemtraB#3\dotEX &\ifbool{Bds}{{\color{\maudapan}\footnotesize\faCheck}}{} &\ifbool{Bds}{}{{\color{\maudapan}\footnotesize\faClose}}  \\
		 \hline
		 \end{tabular}~
		 \begin{tabular}[t]{|p{\widthNPB}|C{0.35cm}|C{.35cm}|}
		 \hline
		 \thead{\normalsize\sffamily Phát biểu}&{\sffamily Đ} & {\sffamily S}\\
		 \hline
		 {c)} \KiemtraC#4\dotEX &\ifbool{Cds}{{\color{\maudapan}\footnotesize\faCheck}}{} &\ifbool{Cds}{}{{\color{\maudapan}\footnotesize\faClose}}  \\
		 \hline
		 {d)} \KiemtraD#5\dotEX &\ifbool{Dds}{{\color{\maudapan}\footnotesize\faCheck}}{} &\ifbool{Dds}{}{{\color{\maudapan}\footnotesize\faClose}}  \\
		 \hline
		 \end{tabular}
		 \end{tabular}
		\end{center}
 }{%
	 \begin{multicols}{#1}
	 \begin{enumerate}[wide=0.65cm,label*=\protect{\alph*)},itemsep=-3pt,topsep=-4pt]
		 \item \KiemtraA#2\dotEX	
		 \item \KiemtraB#3\dotEX
		 \item \KiemtraC#4\dotEX
		 \item \KiemtraD#5\dotEX
	 \end{enumerate}
	 \end{multicols}	
	 }
 }
 }
\Writetofile{ansbook}{\string\def\string\writeANS{\saveans}}
  \scantokens{
	\begin{solbook}
		\writeANS
	\end{solbook}
  }%
\setcounter{numTrue}{0}%
}


\newcommand{\loigiaiTF}{
	 \renewcommand{\loigiai}[1]{%
	 \par\noindent%
	 {\xdef\tieudehinh{}\par\noindent%
	 	 {\color{mausubsection}\reflectbox{\Large\WritingHand}\ {Lời giải.}} ##1\signed{\textcolor{mauchuong}{\faKey}~\hspace*{-7pt}
		 \ifbool{Ads}{\TLdung{A}}{\TLsai{A}}~\hspace*{-7pt}
		 \ifbool{Bds}{\TLdung{B}}{\TLsai{B}}~\hspace*{-7pt}
		 \ifbool{Cds}{\TLdung{C}}{\TLsai{C}}~\hspace*{-7pt}
		 \ifbool{Dds}{\TLdung{D}}{\TLsai{D}}}}
	 }%
}
\newenvironment{itemchoice}%
	{\begin{enumerate}[wide=0.65cm,label*=\protect{\alph*)},itemsep=-3pt,topsep=-4pt]}%
	{\end{enumerate}}
\def\itemch{\item}

%\renewenvironment{Solbook}[1]{{\noindent{\textsf{Câu~#1}}} \ignorespaces}{}

\RenewEnviron{Solbook}[1]{%
	\resizebox{0.22\linewidth}{!}{\tikz{
		\node[font=\bfseries,align=center,minimum height=19pt] (A){Câu #1. \BODY};
		\draw[teal,rounded corners,thick]([yshift=-1pt]A.north west) rectangle ([yshift=1pt]A.south east);
	}}%
}

\newcommand{\bangdapanTF}[1]{
	\begin{center}
		\textbf{\textsf{BẢNG ĐÁP ÁN ĐÚNG SAI Đ/S}}
	\end{center}
	\begin{multicols}{4}
	\input{Ansbook/#1}
	\end{multicols}
}

\usepackage{xstring}
\newcommand{\khungB}[1]{%
  \StrSubstitute{#1}{$}{}[\chuoi] % Loại bỏ dấu $
  \StrLen{\chuoi}[\strlen]%
  \ifnum\strlen<4%
    \def\strlen{4}%
  \fi%
%  \centering % Căn giữa toàn bộ khối
  \foreach \i in {1,...,\strlen} {%
    \StrChar{\chuoi}{\i}[\char]%
    \tikz{\node[draw=teal,
    	minimum width=17pt,minimum height=17pt,align=center,
      	rounded corners=1pt,inner sep=1pt,anchor=south] (char) {\char};%
      }%
  }%
}

%%%================hiển thị lời giải trắc nghiệm================%%%
\def\signed#1{{\unskip\nobreak\hfil\penalty50
  \hskip2em\hbox{}\nobreak\hfil#1%
  \parfillskip=0pt \finalhyphendemerits=0 \par}}
\newcommand{\loigiaiDS}[1]{%
 \par\noindent%
 {\xdef\tieudehinh{}\par\noindent%
 	 {\color{mausubsection}\reflectbox{\Large\WritingHand}\ {\fontfamily{qzc}\selectfont Lời giải.}} #1\signed{ \textcolor{mauchuong}{\faKey}~\hspace*{-7pt}
	 \ifbool{Ads}{\TLdung{A}}{\TLsai{A}}~\hspace*{-7pt}
	 \ifbool{Bds}{\TLdung{B}}{\TLsai{B}}~\hspace*{-7pt}
	 \ifbool{Cds}{\TLdung{C}}{\TLsai{C}}~\hspace*{-7pt}
	 \ifbool{Dds}{\TLdung{D}}{\TLsai{D}}}}
}%	
				
\def\hienthiloigiai{
	\global\setbool{TNdungsai}{false}%
	\xdef\maudapan{black}%
	\AtBeginEnvironment{ex}{%
		\renewcommand{\loigiaiTF}{%
			 \renewcommand{\loigiai}[1]{%
				 \par\noindent%
				 {\xdef\tieudehinh{}\par\noindent%
				 	 {\color{mausubsection}\reflectbox{\Large\WritingHand}\ {\fontfamily{qzc}\selectfont Lời giải.}} ####1\signed{ \textcolor{mauchuong}{\faKey}~\hspace*{-7pt}
					 \ifbool{Ads}{\TLdung{A}}{\TLsai{A}}~\hspace*{-7pt}
					 \ifbool{Bds}{\TLdung{B}}{\TLsai{B}}~\hspace*{-7pt}
					 \ifbool{Cds}{\TLdung{C}}{\TLsai{C}}~\hspace*{-7pt}
					 \ifbool{Dds}{\TLdung{D}}{\TLsai{D}}}}
			}%				
		}
		\renewcommand{\shortans}[2][]{\phantom{.}\signed{{\color{mauchuong}\raisebox{6pt}{\faKey}}\;\khungB{##2}}}
		\ifbool{TNdungsai}{%
			 \renewcommand{\loigiai}[1]{%
				 \par\noindent%
				 {\xdef\tieudehinh{}\par\noindent%
				 	 {\color{mausubsection}\reflectbox{\Large\WritingHand}\ {\fontfamily{qzc}\selectfont Lời giải.}} ##1\signed{% \textcolor{mauchuong}{\faKey}~\hspace*{-7pt}
					 \ifbool{Ads}{\TLdung{A}}{\TLsai{A}}~\hspace*{-7pt}
					 \ifbool{Bds}{\TLdung{B}}{\TLsai{B}}~\hspace*{-7pt}
					 \ifbool{Cds}{\TLdung{C}}{\TLsai{C}}~\hspace*{-7pt}
					 \ifbool{Dds}{\TLdung{D}}{\TLsai{D}}}}
				}%
			}{%
				 \renewcommand{\loigiai}[1]{\gdef\tieudehinh{}\par\noindent%
				 {\color{mausubsection}\reflectbox{\Large\WritingHand}\ {\fontfamily{qzc}\selectfont Lời giải.}} ##1
				 \ifthenelse{\value{numTrueSol}>0}{%
				 \signed{{\color{mauchuong}\faKey~\circlenum{\Alph{numTrueSol}}}}%
			}{}}%
		}%					
	}
	\AtBeginEnvironment{bt}{%
		\renewcommand{\shortans}[2][]{%
			 \scantokens{%
				 \begin{giaibt}
				 	##2
				 \end{giaibt}}% 
			 \phantom{.}\signed{{\color{mauchuong}\raisebox{6pt}{\faKey}}\;\khungB{##2}}}%
			 \renewcommand{\loigiai}[1]{%
				 \par\noindent%
				 {\xdef\tieudehinh{}\par\noindent%
				 	 {\color{mausubsection}\reflectbox{\Large\WritingHand}\ {\fontfamily{qzc}\selectfont Lời giải.}} ##1
				 }%
		   }%
	}
	
}

%%%================tắt lời giải trắc nghiệm================%%%
\def\anloigiai{
	\xdef\maudapan{white}%
	\AtBeginEnvironment{ex}{%
		\renewcommand{\loigiai}[1]{}
		\renewcommand{\loigiaiTF}{}
		\renewcommand{\loigiaiDS}{}
	}
	\AtBeginEnvironment{bt}{%
		\renewcommand{\loigiai}[1]{}
		\renewcommand{\shortans}[2][]{
					 \scantokens{%
						 \begin{giaibt}
						 	##2
						 \end{giaibt}}% 
			  \signed{\foreach \i in {1,...,4} {%
			    \tikz{\node[draw=teal,
			    	minimum width=17pt,minimum height=17pt,align=center,
			      	rounded corners=1pt,inner sep=1pt,anchor=south] (char) {};%
			      }%
			  }%			
		}}
	}
}

%%%=============lưu lời giải trắc nghiệm=============%%%
\Newassociation{giaiex}{loigiaiex}{ansex}
\def\luuloigiai{
	\xdef\maudapan{white}%
	\AtBeginEnvironment{ex}{%
		\renewcommand{\loigiaiTF}{}
		\renewcommand{\loigiai}[1]{%
		 \Writetofile{ansex}{\string\def\string\writeANS{\string\faKey\, \saveans}}
		 \ifbool{TNdungsai}{}{\Writetofile{ansex}{\string\def\string\writeANS{\string\faKey\, \circlenum{\Alph{numTrueSol}}}}}
				 \scantokens{%
				 \begin{giaiex}
					 ##1
				 \end{giaiex}
				}% 
		}
	}
	\AtBeginEnvironment{bt}{%
		\renewcommand{\loigiai}[1]{%
				\Writetofile{ansex}{\string\def\string\writeANS{}}
				 \scantokens{%
				 \begin{giaiex}
					 ##1
				 \end{giaiex}
				}% 
		}
	}
}

\NewTColorBox{boxlgex}{+m}{
	 breakable,
	 enhanced,
	 outer arc=0pt,
	 arc=0pt,
	 colframe=white,
	 frame hidden,
	 left=-6pt,right=0pt,top=0pt,
	 colback=white,
	 attach boxed title to top left,
	 boxed title style={
	 colback=white,
	 outer arc=0pt,
	 arc=0pt,
	 top=1pt,
	 bottom=1pt,
	 left=3pt,
	 right=3pt,
	 colframe=white
 },
 fonttitle=\bfseries\sffamily\selectfont\color{white},
 title={\faGg\; Câu~#1},
 overlay unbroken and first={
 \path (title.north west) coordinate (A)
 ($ (title.south west) +(-4pt,0pt)$) coordinate (B)
 (title.south east) coordinate (C)
 ($ (title.north east)+(4pt,0) $) coordinate (D);
 \path[rounded corners=2pt,fill=mauchuong,
 preaction={transform canvas={shift={(-45:2pt)}},left color=mauchuong!45,right color=mauchuong!25}] 
 (A)--(B)--(C)--(D)--cycle;
 \path (A)--(C) node[midway,font=\color{white}\bfseries\sffamily\selectfont](Bai){\textsl{\faGg\ Câu~#1}};
 \draw[rounded corners=2pt,thick,mauchuong] ([xshift=-3pt]B) coordinate (Bt)
 --([shift={(-2pt,2pt)}]A)--+(\linewidth,0) coordinate (Ct);
 \path ($(Ct)+(210:3pt)$) node[anchor=north east,text=mausection]{\writeANS};
 \fill[mauchuong] (Bt) circle (1pt) (Ct) circle (2pt);
}
}
 
\renewenvironment{loigiaiex}[1]{\begin{boxlgex}{#1}}{\end{boxlgex}}

\RenewEnviron{listEX}[1][1]{
	\ifthenelse{\equal{#1}{1}}%
	{%
	\begin{enumerate}%[label*=\circlenum{\Alph*},itemsep=-1pt,topsep=0pt]%
		\BODY%
	\end{enumerate}%
	}%
	{\begin{multicols}{#1}
	\begin{enumerate}%[label*=\circlenum{\Alph*},itemsep=-1pt,topsep=0pt]%
		\BODY%
	\end{enumerate}%
	\end{multicols}}%
}

%%%================Đổi thành dòng kẻ lời giải==================%%%
\newcommand{\Pointilles}[2][1.1]{%
 \par\nobreak%
 \noindent\rule{0pt}{1.1\baselineskip}%
 \foreach \i in {1,...,#2}{%
 \ifnum\i=1
 \noindent\makebox[\linewidth]{\rule{0pt}{#1\baselineskip}\reflectbox{\Large\WritingHand}\dotfill}\endgraf
 \else
 \noindent\makebox[\linewidth]{\rule{0pt}{#1\baselineskip}\dotfill}\endgraf
 \fi
 }
}

\newcommand{\DoiThanhDongKe}[1]{
\setbox0=\vbox{\hbox{
\noindent\begin{minipage}{\linewidth}%
#1 aaaaa
\end{minipage}
}}
\linepar=\ht0
\pgfmathparse{int((\linepar-\fboxsep)/(\lineheight)+1)}
	\ifnum\pgfmathresult>3
	\noindent%
		\Pointilles{\pgfmathresult}
	\else
	\noindent%
		\Pointilles{3}
	\fi
}

\newcommand{\DoiThanhDongKeH}[1]{
\setbox0=\vbox{\hbox{
\noindent\begin{minipage}{\linewidth}%
#1 aaaaa
\end{minipage}
}}
\linepar=\ht0
\pgfmathparse{int((\linepar-\fboxsep)/(0.85\lineheight)+1)}
	\ifnum\pgfmathresult>3
	\noindent%
		\Pointilles{\pgfmathresult}
	\else
	\noindent%
		\Pointilles{4}
	\fi
}

%%%=============================================%%%
\def\dongkeloigiai{
	\AtBeginEnvironment{ex}{
		\def\maudapan{white}
		\renewcommand{\loigiaiTF}{}
		\renewcommand{\loigiai}[1]{%
		 \DoiThanhDongKe{##1}
		}
		\renewcommand{\loigiaiDS}[1]{%
		 \DoiThanhDongKe{##1}
		}
	}
}


%%%=============================================%%%
\def\dongkeHaicot{
	\AtBeginEnvironment{ex}{
		\def\maudapan{white}
		\renewcommand{\loigiaiTF}{}
		\renewcommand{\loigiai}[1]{
			\begin{multicols}{2}			
				\DoiThanhDongKeH{##1}
			\end{multicols}%
		}
		\renewcommand{\loigiaiDS}[1]{
			\begin{multicols}{2}			
				\DoiThanhDongKeH{##1}
			\end{multicols}%
		}
	}
	\AtBeginEnvironment{bt}{
		\renewcommand{\loigiai}[1]{
			\begin{multicols}{2}			
				\DoiThanhDongKeH{##1}
			\end{multicols}%
		}
	}	
}

%%%================Đổi thành dòng kẻ lời giải==================%%%

\RenewEnviron{enumEX}[2][]{
	\ifthenelse{\equal{#2}{1}}%
	{%
	\begin{enumerate}%[label*=\circlenum{\Alph*},itemsep=-1pt,topsep=0pt]%
		\BODY%
	\end{enumerate}%
	}%
	{\begin{multicols}{#2}
	\begin{enumerate}%[label*=\circlenum{\Alph*},itemsep=-1pt,topsep=0pt]%
		\BODY%
	\end{enumerate}%
	\end{multicols}}%
}

\newcounter{bt}
\makeatletter
\newenvironment{bt}[1][]{\par\noindent%
	\refstepcounter{bt}%
	\ifthenelse{\equal{#1}{c}}{%
	\@ifnextchar\immini{\vspace*{-3pt}\gdef\tieudehinh{{\bfseries\color{mauchuong}\checkmarkH Câu \arabic{bt}. }}}%
	{{\bfseries\color{mauchuong}\checkmarkH Câu \arabic{bt}. }}%\arabic{section}.
	}{%không tùy chọn
	\@ifnextchar\immini{\vspace*{-3pt}\gdef\tieudehinh{{\bfseries\color{mauchuong}Câu \arabic{bt}. }}}%
	{{\bfseries\color{mauchuong}Câu \arabic{bt}. }}%\arabic{section}.
}}%
{\xdef\tieudehinh{}}
\makeatother

 
%%%===============cau trả lời ngắn=======================%%%
\Newassociation{giaibt}{loigiaibt}{ansbt}
\tikzset{
    master/.style={
        execute at end picture={
            \coordinate (lower right) at (current bounding box.south east);
            \coordinate (upper left) at (current bounding box.north west);
        }
    },
    slave/.style={
        execute at end picture={
            \pgfresetboundingbox
            \path (upper left) rectangle (lower right);
        }
    }
}

\newcommand{\ovuong}[2][2.5]{%
 \scantokens{%
	 \begin{giaibt}
	 #2
	 \end{giaibt}}% 
	 {\linepenalty100\exhyphenpenalty0
	\signed{{\tikz{%
		\path (0,-0.15) node (char) {};
		\draw[rounded corners=2pt,thick,mauchuong] (0,-0.275) rectangle (#1,0.275);
	 }\hspace*{-9pt}%
 }}%
}}

\newcommand{\shortans}[2][2.5]{%
	\ovuong{#2}%
}

\RenewEnviron{loigiaibt}[1]{%
	\resizebox{0.235\linewidth}{!}{\textbf{Câu~#1.}\khungB{\BODY}}\;
%	\tikz{\node[rounded corners,draw=mauchuong,text width=4cm,align=center,minimum height=21pt] 
%		{\textbf{Câu~#1.} {\boldmath\textbf{\color{mausubsection}\BODY}}}%
%		}
	}%
	
\newcommand{\bangdapanSA}[1]{%
	\begin{center}
		{\sffamily\bfseries ĐÁP SỐ PHẦN TRẢ LỜI NGẮN}
	\end{center}
	
	\medskip
	\begin{flushleft}
		\begin{multicols}{4}%
			\input{Ansbook/#1}
		\end{multicols}%
	\end{flushleft}
}	


\newcommand{\vong}[1]{%
\begin{tikzpicture}[baseline=(vong.base)]
		\path (0,0) node(vong) {};
		\path (0pt,-1pt) node[anchor=south, circle,fill=mauchuong,inner sep=1pt,font=\footnotesize\bfseries,text=white,outer sep=0pt] {#1};
	\end{tikzpicture}
}
\usepackage{tkz-euclide,tkz-tab}
\usetikzlibrary{angles,intersections,shapes.geometric,arrows,decorations.markings,arrows.meta,patterns.meta}
\usetikzlibrary{shadows.blur,patterns,quotes}
\raggedbottom
\tikzset{>=stealth,line join=round,line cap=butt}
\newcommand{\khungct}[1]{\tcboxmath[colback=mausection!5,colframe=mausection,boxsep=1pt,boxrule=1pt,arc is angular]{#1}}
\newcommand{\khung}[1]{\tcbox[colback=mausection!5,colframe=mausection,boxsep=1pt,boxrule=1pt,arc is angular]{#1}}
\usetikzlibrary{bending,matrix}%Mũi tên chỗ vòng cung
\tikzset{line cap=round,line join=round,>=stealth}
\def\chongbang{\rule[-14pt]{0pt}{33pt}}
\usepackage{tasks}
\usepackage{bm}
\usepackage{paracol}
\usepackage{casio580x}
\usepackage[hidelinks]{hyperref}
\def\xuongdong{\break}
\def\hinhlaptophnih{%
	\begin{tikzpicture}[line cap=round, line join=round]
		\definecolor{c4e4c4d}{RGB}{78, 76, 77}
		\definecolor{cf6f6f6}{RGB}{246, 246, 246}
		\definecolor{cdfdfdf}{RGB}{223, 223, 223}
		\definecolor{c818488}{RGB}{129, 132, 136}
		\definecolor{cb8c5e2}{RGB}{184, 197, 226}
		\definecolor{ce6e7e8}{RGB}{230, 231, 232}
		\definecolor{cc4d1cf}{RGB}{196, 209, 207}
		\definecolor{cd3ccd2}{RGB}{211, 204, 210}
		\definecolor{cededed}{RGB}{237, 237, 237}
		\path[fill=c4e4c4d] (53.86, 7.74) -- (32.03, 7.74) .. controls +(180:0.4) and +(-127:0.2) ..
			(31.52, 8.45) -- (34.92, 12.7) .. controls +(52:0.1) and +(180:0.3) ..
			(35.56, 12.87) -- (50.33, 12.87) .. controls +(0:0.3) and +(128:0.1) ..
			(50.97, 12.7) -- (54.38, 8.44) .. controls +(-54:0.3) and +(0:0.4) ..
			(53.86, 7.74) -- cycle;

		\path[fill=cf6f6f6] (53.86, 8.14)-- (32.03, 8.14) .. controls +(180:0.4) and +(-127:0.3) ..
			(31.55, 8.5) coordinate (C) -- (34.92, 13.1) .. controls +(52:0.1) and +(180:0.3) ..
			(35.56, 13.27)  -- (50.33, 13.27) .. controls +(0:0.3) and +(128:0.1) ..
			(50.97, 13.1) -- (54.35, 8.5) .. controls +(-53:0.2) and +(0:0.4) ..
			(53.86, 8.14) -- cycle;

		\path[draw=cdfdfdf, opacity=0.5, miter limit=10] (31.91, 8.53) -- (35.14, 13.06) .. controls +(59:0.1) and +(180:0.1) ..
			(35.48, 13.15) -- (50.45, 13.15) .. controls +(0:0.1) and +(129:0.1) ..
			(50.79, 13.06) -- (54.03, 8.53) .. controls +(-54:0.1) and +(-2:0.2) ..
			(53.77, 8.4) -- (32.16, 8.4) .. controls +(175:0.2) and +(-122:0.1) ..
			(31.91, 8.53) -- cycle;

			\path[fill=c818488] (35.56, 23.75) -- (50.33, 23.75) .. controls +(0:0.3) and +(90:0.3) ..
				(50.84, 23.24) -- (50.84, 14.01) .. controls +(-90:0.3) and +(0:0.3) ..
				(50.33, 13.5) -- (35.56, 13.5) .. controls +(180:0.3) and +(-90:0.3) ..
				(35.05, 14.01) -- (35.05, 23.24) .. controls +(90:0.3) and +(180:0.3) ..
				(35.56, 23.75) -- cycle
 (35.56, 23.52) -- (50.33, 23.52) .. controls +(0:0.3) and +(90:0.3) ..
				(50.84, 23.01) -- (50.84, 13.78) .. controls +(-90:0.3) and +(0:0.3) ..
				(50.33, 13.27) -- (35.56, 13.27) .. controls +(180:0.3) and +(-90:0.3) ..
				(35.05, 13.78) coordinate (B) -- (35.05, 23.01) coordinate (A).. controls +(90:0.3) and +(180:0.3) ..
				(35.56, 23.52) -- cycle
 (35.77, 22.84) rectangle (50.09, 14.03) 
 (35.77, 22.84) rectangle (50.09, 14.03) ;
		\begin{scope}

			\path[fill=cc4d1cf] (35.93, 12.66) -- (35.41, 12.66) .. controls +(180:0.1) and +(-135:0) ..
				(35.34, 12.69) -- (35.4, 12.79) .. controls +(63:0) and +(180:0) ..
				(35.51, 12.82) -- (36.03, 12.82) .. controls +(0:0) and +(63:0) ..
				(36.1, 12.79) -- (36.04, 12.69) .. controls +(-117:0) and +(0:0) ..
				(35.93, 12.66) -- cycle
 (36.84, 12.66) -- (36.32, 12.66) .. controls +(180:0) and +(-117:0) ..
				(36.25, 12.69) -- (36.31, 12.79) .. controls +(63:0) and +(180:0.1) ..
				(36.42, 12.82) -- (36.93, 12.82) .. controls +(0:0) and +(63:0) ..
				(37.01, 12.79) -- (36.96, 12.69) .. controls +(-117:0) and +(0:0.1) ..
				(36.84, 12.66) -- cycle
 (37.76, 12.66) -- (37.24, 12.66) .. controls +(180:0.1) and +(-117:0) ..
				(37.17, 12.69) -- (37.22, 12.79) .. controls +(90:0) and +(180:0) ..
				(37.32, 12.82) -- (37.84, 12.82) .. controls +(0:0) and +(63:0) ..
				(37.91, 12.79) -- (37.87, 12.69) .. controls +(-117:0) and +(0:0.1) ..
				(37.76, 12.66) -- cycle
 (38.68, 12.66) -- (38.16, 12.66) .. controls +(180:0) and +(-117:0) ..
				(38.08, 12.69) -- (38.12, 12.79) .. controls +(63:0) and +(180:0) ..
				(38.23, 12.82) -- (38.74, 12.82) .. controls +(0:0) and +(63:0) ..
				(38.82, 12.79) -- (38.79, 12.69) .. controls +(-117:0) and +(0:0) ..
				(38.68, 12.66) -- cycle
 (39.6, 12.66) -- (39.08, 12.66) .. controls +(180:0) and +(-117:0) ..
				(39, 12.69) -- (39.03, 12.79) .. controls +(63:0) and +(180:0.1) ..
				(39.14, 12.82) -- (39.65, 12.82) .. controls +(0:0.1) and +(63:0) ..
				(39.73, 12.79) -- (39.7, 12.69) .. controls +(-90:0) and +(0:0) ..
				(39.6, 12.66) -- cycle
 (40.52, 12.66) -- (40, 12.66) .. controls +(180:0) and +(-90:0) ..
				(39.91, 12.69) -- (39.94, 12.79) .. controls +(90:0) and +(180:0) ..
				(40.04, 12.82) -- (40.55, 12.82) .. controls +(0:0.1) and +(90:0) ..
				(40.64, 12.79) -- (40.62, 12.69) .. controls +(-90:0) and +(0:0) ..
				(40.52, 12.66) -- cycle
 (41.44, 12.66) -- (40.92, 12.66) .. controls +(180:0.1) and +(-90:0) ..
				(40.83, 12.69) -- (40.85, 12.79) .. controls +(90:0) and +(180:0.1) ..
				(40.95, 12.82) -- (41.46, 12.82) .. controls +(0:0) and +(90:0) ..
				(41.55, 12.79) -- (41.53, 12.69) .. controls +(-90:0) and +(0:0.1) ..
				(41.44, 12.66) -- cycle
 (42.36, 12.66) -- (41.84, 12.66) .. controls +(180:0.1) and +(-117:0) ..
				(41.75, 12.69) -- (41.76, 12.79) .. controls +(90:0) and +(180:0.1) ..
				(41.85, 12.82) -- (42.36, 12.82) .. controls +(0:0.1) and +(63:0) ..
				(42.45, 12.79) -- (42.45, 12.69) .. controls +(-90:0) and +(0:0) ..
				(42.36, 12.66) -- cycle
 (43.27, 12.66) -- (42.75, 12.66) .. controls +(180:0) and +(-90:0) ..
				(42.66, 12.69) -- (42.66, 12.79) .. controls +(90:0) and +(180:0) ..
				(42.76, 12.82) -- (43.27, 12.82) .. controls +(0:0) and +(90:0) ..
				(43.36, 12.79) -- (43.37, 12.69) .. controls +(-90:0) and +(0:0.1) ..
				(43.27, 12.66) -- cycle
 (44.19, 12.66) -- (43.67, 12.66) .. controls +(180:0.1) and +(-90:0) ..
				(43.58, 12.69) -- (43.57, 12.79) .. controls +(90:0) and +(180:0) ..
				(43.66, 12.82) -- (44.17, 12.82) .. controls +(0:0.1) and +(90:0) ..
				(44.27, 12.79) -- (44.28, 12.69) .. controls +(-90:0) and +(0:0.1) ..
				(44.19, 12.66) -- cycle
 (45.11, 12.66) -- (44.59, 12.66) .. controls +(180:0.1) and +(-90:0) ..
				(44.49, 12.69) -- (44.48, 12.79) .. controls +(90:0) and +(180:0) ..
				(44.57, 12.82) -- (45.08, 12.82) .. controls +(0:0.1) and +(90:0) ..
				(45.18, 12.79) -- (45.2, 12.69) .. controls +(-90:0) and +(0:0) ..
				(45.11, 12.66) -- cycle
 (46.03, 12.66) -- (45.51, 12.66) .. controls +(180:0) and +(-90:0) ..
				(45.41, 12.69) -- (45.39, 12.79) .. controls +(117:0) and +(180:0) ..
				(45.47, 12.82) -- (45.99, 12.82) .. controls +(0:0) and +(117:0) ..
				(46.09, 12.79) -- (46.11, 12.69) .. controls +(-63:0) and +(0:0) ..
				(46.03, 12.66) -- cycle
 (46.95, 12.66) -- (46.43, 12.66) .. controls +(180:0) and +(-63:0) ..
				(46.32, 12.69) -- (46.3, 12.79) .. controls +(117:0) and +(180:0.1) ..
				(46.38, 12.82) -- (46.89, 12.82) .. controls +(0:0) and +(90:0) ..
				(46.99, 12.79) -- (47.03, 12.69) .. controls +(-63:0) and +(0:0) ..
				(46.95, 12.66) -- cycle
 (47.87, 12.66) -- (47.35, 12.66) .. controls +(180:0.1) and +(-63:0) ..
				(47.24, 12.69) -- (47.2, 12.79) .. controls +(90:0) and +(180:0.1) ..
				(47.28, 12.82) -- (47.8, 12.82) .. controls +(0:0.1) and +(117:0) ..
				(47.9, 12.79) -- (47.95, 12.69) .. controls +(-90:0) and +(0:0.1) ..
				(47.87, 12.66) -- cycle
 (48.78, 12.66) -- (48.26, 12.66) .. controls +(180:0) and +(-90:0) ..
				(48.16, 12.69) -- (48.11, 12.79) .. controls +(117:0) and +(180:0) ..
				(48.19, 12.82) -- (48.7, 12.82) .. controls +(0:0) and +(117:0) ..
				(48.81, 12.79) -- (48.86, 12.69) .. controls +(-63:0) and +(0:0.1) ..
				(48.78, 12.66) -- cycle
 (49.7, 12.66) -- (49.18, 12.66) .. controls +(180:0) and +(-63:0) ..
				(49.07, 12.69) -- (49.02, 12.79) .. controls +(117:0) and +(180:0.1) ..
				(49.09, 12.82) -- (49.61, 12.82) .. controls +(0:0) and +(117:0) ..
				(49.72, 12.79) -- (49.78, 12.69) .. controls +(-63:0) and +(0:0) ..
				(49.7, 12.66) -- cycle
 (50.62, 12.66) -- (50.1, 12.66) .. controls +(180:0.1) and +(-63:0) ..
				(49.99, 12.69) -- (49.93, 12.79) .. controls +(117:0) and +(180:0) ..
				(50, 12.82) -- (50.51, 12.82) .. controls +(0:0.1) and +(135:0) ..
				(50.63, 12.79) -- (50.69, 12.69) .. controls +(-63:0) and +(0:0.1) ..
				(50.62, 12.66) -- cycle
 (35.84, 12.29) -- (35.16, 12.29) .. controls +(180:0.1) and +(-135:0) ..
				(35.09, 12.32) -- (35.24, 12.56) .. controls +(27:0) and +(180:0) ..
				(35.36, 12.59) -- (36.02, 12.59) .. controls +(0:0) and +(27:0) ..
				(36.09, 12.56) -- (35.95, 12.32) .. controls +(-117:0) and +(0:0) ..
				(35.84, 12.29) -- cycle
 (36.87, 12.29) -- (36.18, 12.29) .. controls +(180:0) and +(-117:0) ..
				(36.11, 12.32) -- (36.25, 12.56) .. controls +(45:0) and +(180:0) ..
				(36.36, 12.59) -- (37.02, 12.59) .. controls +(0:0) and +(45:0) ..
				(37.1, 12.56) -- (36.98, 12.32) .. controls +(-117:0) and +(0:0.1) ..
				(36.87, 12.29) -- cycle
 (37.89, 12.29) -- (37.21, 12.29) .. controls +(180:0.1) and +(-117:0) ..
				(37.13, 12.32) -- (37.25, 12.56) .. controls +(45:0) and +(180:0) ..
				(37.36, 12.59) -- (38.02, 12.59) .. controls +(0:0) and +(45:0) ..
				(38.1, 12.56) -- (38, 12.32) .. controls +(-117:0) and +(0:0.1) ..
				(37.89, 12.29) -- cycle
 (38.92, 12.29) -- (38.24, 12.29) .. controls +(180:0.1) and +(-117:0) ..
				(38.16, 12.32) -- (38.25, 12.56) .. controls +(45:0) and +(180:0) ..
				(38.36, 12.59) -- (39.03, 12.59) .. controls +(0:0) and +(90:0) ..
				(39.11, 12.56) -- (39.03, 12.32) .. controls +(-117:0) and +(0:0) ..
				(38.92, 12.29) -- cycle
 (39.95, 12.29) -- (39.26, 12.29) .. controls +(180:0) and +(-90:0) ..
				(39.18, 12.32) -- (39.26, 12.56) .. controls +(90:0) and +(180:0) ..
				(39.36, 12.59) -- (40.03, 12.59) .. controls +(0:0) and +(90:0) ..
				(40.11, 12.56) -- (40.05, 12.32) .. controls +(-90:0) and +(0:0) ..
				(39.95, 12.29) -- cycle
 (40.98, 12.29) -- (40.29, 12.29) .. controls +(180:0) and +(-117:0) ..
				(40.21, 12.32) -- (40.26, 12.56) .. controls +(45:0) and +(180:0) ..
				(40.36, 12.59) -- (41.03, 12.59) .. controls +(0:0) and +(45:0) ..
				(41.11, 12.56) -- (41.07, 12.32) .. controls +(-90:0) and +(0:0.1) ..
				(40.98, 12.29) -- cycle
 (42, 12.29) -- (41.32, 12.29) .. controls +(180:0) and +(-90:0) ..
				(41.23, 12.32) -- (41.27, 12.56) .. controls +(90:0) and +(180:0) ..
				(41.36, 12.59) -- (42.03, 12.59) .. controls +(0:0) and +(90:0) ..
				(42.12, 12.56) -- (42.1, 12.32) .. controls +(-90:0) and +(0:0) ..
				(42, 12.29) -- cycle
 (43.03, 12.29) -- (42.35, 12.29) .. controls +(180:0.1) and +(-90:0) ..
				(42.25, 12.32) -- (42.27, 12.56) .. controls +(90:0) and +(180:0) ..
				(42.36, 12.59) -- (43.03, 12.59) .. controls +(0:0) and +(90:0) ..
				(43.12, 12.56) -- (43.12, 12.32) .. controls +(-90:0) and +(0:0) ..
				(43.03, 12.29) -- cycle
 (44.06, 12.29) -- (43.37, 12.29) .. controls +(180:0) and +(-90:0) ..
				(43.28, 12.32) -- (43.27, 12.56) .. controls +(90:0) and +(180:0) ..
				(43.36, 12.59) -- (44.03, 12.59) .. controls +(0:0) and +(90:0) ..
				(44.12, 12.56) -- (44.15, 12.32) .. controls +(-90:0) and +(0:0) ..
				(44.06, 12.29) -- cycle
 (45.09, 12.29) -- (44.4, 12.29) .. controls +(180:0) and +(-63:0) ..
				(44.3, 12.32) -- (44.28, 12.56) .. controls +(135:0) and +(180:0) ..
				(44.36, 12.59) -- (45.03, 12.59) .. controls +(0:0) and +(135:0) ..
				(45.13, 12.56) -- (45.17, 12.32) .. controls +(-63:0) and +(0:0) ..
				(45.09, 12.29) -- cycle
 (46.11, 12.29) -- (45.43, 12.29) .. controls +(180:0) and +(-90:0) ..
				(45.33, 12.32) -- (45.28, 12.56) .. controls +(90:0) and +(180:0) ..
				(45.36, 12.59) -- (46.03, 12.59) .. controls +(0:0) and +(90:0) ..
				(46.13, 12.56) -- (46.2, 12.32) .. controls +(-90:0) and +(0:0) ..
				(46.11, 12.29) -- cycle
 (47.14, 12.29) -- (46.46, 12.29) .. controls +(180:0.1) and +(-63:0) ..
				(46.35, 12.32) -- (46.28, 12.56) .. controls +(90:0) and +(180:0) ..
				(46.36, 12.59) -- (47.03, 12.59) .. controls +(0:0) and +(90:0) ..
				(47.13, 12.56) -- (47.22, 12.32) .. controls +(-63:0) and +(0:0) ..
				(47.14, 12.29) -- cycle
 (48.17, 12.29) -- (47.48, 12.29) .. controls +(180:0) and +(-90:0) ..
				(47.38, 12.32) -- (47.29, 12.56) .. controls +(135:0) and +(180:0) ..
				(47.37, 12.59) -- (48.03, 12.59) .. controls +(0:0) and +(135:0) ..
				(48.14, 12.56) -- (48.24, 12.32) .. controls +(-63:0) and +(0:0) ..
				(48.17, 12.29) -- cycle
 (49.19, 12.29) -- (48.51, 12.29) .. controls +(180:0) and +(-63:0) ..
				(48.4, 12.32) -- (48.29, 12.56) .. controls +(135:0) and +(180:0) ..
				(48.37, 12.59) -- (49.03, 12.59) .. controls +(0:0) and +(135:0) ..
				(49.14, 12.56) -- (49.27, 12.32) .. controls +(-63:0) and +(0:0.1) ..
				(49.19, 12.29) -- cycle
 (50.88, 12.29) -- (49.54, 12.29) .. controls +(180:0) and +(-45:0) ..
				(49.42, 12.32) -- (49.29, 12.56) .. controls +(135:0) and +(180:0) ..
				(49.37, 12.59) -- (50.67, 12.59) .. controls +(0:0) and +(153:0) ..
				(50.79, 12.56) -- (50.95, 12.32) .. controls +(-63:0) and +(0:0) ..
				(50.88, 12.29) -- cycle
 (35.88, 11.87) -- (34.89, 11.87) .. controls +(180:0.1) and +(-117:0) ..
				(34.81, 11.91) -- (34.98, 12.16) .. controls +(63:0) and +(180:0.1) ..
				(35.1, 12.2) -- (36.07, 12.2) .. controls +(0:0) and +(63:0) ..
				(36.15, 12.16) -- (36, 11.91) .. controls +(-117:0) and +(0:0.1) ..
				(35.88, 11.87) -- cycle
 (36.94, 11.87) -- (36.25, 11.87) .. controls +(180:0.1) and +(-135:0) ..
				(36.17, 11.91) -- (36.3, 12.16) .. controls +(45:0) and +(180:0.1) ..
				(36.43, 12.2) -- (37.1, 12.2) .. controls +(0:0.1) and +(63:0) ..
				(37.18, 12.16) -- (37.06, 11.91) .. controls +(-117:0) and +(0:0.1) ..
				(36.94, 11.87) -- cycle
 (38, 11.87) -- (37.31, 11.87) .. controls +(180:0.1) and +(-117:0) ..
				(37.23, 11.91) -- (37.34, 12.16) .. controls +(63:0) and +(180:0.1) ..
				(37.46, 12.2) -- (38.14, 12.2) .. controls +(0:0) and +(63:0) ..
				(38.22, 12.16) -- (38.12, 11.91) .. controls +(-90:0) and +(0:0.1) ..
				(38, 11.87) -- cycle
 (39.07, 11.87) -- (38.37, 11.87) .. controls +(180:0) and +(-117:0) ..
				(38.29, 11.91) -- (38.38, 12.16) .. controls +(63:0) and +(180:0.1) ..
				(38.5, 12.2) -- (39.17, 12.2) .. controls +(0:0.1) and +(63:0) ..
				(39.26, 12.16) -- (39.18, 11.91) .. controls +(-90:0) and +(0:0.1) ..
				(39.07, 11.87) -- cycle
 (40.13, 11.87) -- (39.44, 11.87) .. controls +(180:0.1) and +(-117:0) ..
				(39.35, 11.91) -- (39.42, 12.16) .. controls +(63:0) and +(180:0.1) ..
				(39.53, 12.2) -- (40.21, 12.2) .. controls +(0:0) and +(90:0) ..
				(40.3, 12.16) -- (40.24, 11.91) .. controls +(-90:0) and +(0:0.1) ..
				(40.13, 11.87) -- cycle
 (41.2, 11.87) -- (40.5, 11.87) .. controls +(180:0.1) and +(-117:0) ..
				(40.41, 11.91) -- (40.46, 12.16) .. controls +(90:0) and +(180:0.1) ..
				(40.57, 12.2) -- (41.24, 12.2) .. controls +(0:0.1) and +(90:0) ..
				(41.34, 12.16) -- (41.3, 11.91) .. controls +(-90:0) and +(0:0) ..
				(41.2, 11.87) -- cycle
 (42.26, 11.87) -- (41.56, 11.87) .. controls +(180:0.1) and +(-117:0) ..
				(41.47, 11.91) -- (41.5, 12.16) .. controls +(90:0) and +(180:0.1) ..
				(41.6, 12.2) -- (42.28, 12.2) .. controls +(0:0) and +(90:0) ..
				(42.38, 12.16) -- (42.36, 11.91) .. controls +(-90:0) and +(0:0.1) ..
				(42.26, 11.87) -- cycle
 (43.32, 11.87) -- (42.63, 11.87) .. controls +(180:0.1) and +(-117:0) ..
				(42.53, 11.91) -- (42.53, 12.16) .. controls +(90:0) and +(180:0.1) ..
				(42.64, 12.2) -- (43.31, 12.2) .. controls +(0:0.1) and +(90:0) ..
				(43.41, 12.16) -- (43.42, 11.91) .. controls +(-63:0) and +(0:0.1) ..
				(43.32, 11.87) -- cycle
 (44.39, 11.87) -- (43.69, 11.87) .. controls +(180:0.1) and +(-90:0) ..
				(43.59, 11.91) -- (43.57, 12.16) .. controls +(90:0) and +(180:0.1) ..
				(43.67, 12.2) -- (44.35, 12.2) .. controls +(0:0) and +(90:0) ..
				(44.45, 12.16) -- (44.48, 11.91) .. controls +(-63:0) and +(0:0) ..
				(44.39, 11.87) -- cycle
 (45.45, 11.87) -- (44.75, 11.87) .. controls +(180:0) and +(-90:0) ..
				(44.65, 11.91) -- (44.61, 12.16) .. controls +(90:0) and +(180:0.1) ..
				(44.71, 12.2) -- (45.38, 12.2) .. controls +(0:0.1) and +(90:0) ..
				(45.49, 12.16) -- (45.54, 11.91) .. controls +(-63:0) and +(0:0.1) ..
				(45.45, 11.87) -- cycle
 (46.51, 11.87) -- (45.82, 11.87) .. controls +(180:0.1) and +(-90:0) ..
				(45.71, 11.91) -- (45.65, 12.16) .. controls +(117:0) and +(180:0.1) ..
				(45.74, 12.2) -- (46.42, 12.2) .. controls +(0:0) and +(117:0) ..
				(46.53, 12.16) -- (46.6, 11.91) .. controls +(-63:0) and +(0:0.1) ..
				(46.51, 11.87) -- cycle
 (47.58, 11.87) -- (46.88, 11.87) .. controls +(180:0.1) and +(-63:0) ..
				(46.76, 11.91) -- (46.69, 12.16) .. controls +(117:0) and +(180:0.1) ..
				(46.78, 12.2) -- (47.45, 12.2) .. controls +(0:0.1) and +(117:0) ..
				(47.57, 12.16) -- (47.66, 11.91) .. controls +(-63:0) and +(0:0.1) ..
				(47.58, 11.87) -- cycle
 (48.64, 11.87) -- (47.94, 11.87) .. controls +(180:0) and +(-63:0) ..
				(47.82, 11.91) -- (47.73, 12.16) .. controls +(117:0) and +(180:0.1) ..
				(47.81, 12.2) -- (48.49, 12.2) .. controls +(0:0) and +(117:0) ..
				(48.61, 12.16) -- (48.72, 11.91) .. controls +(-63:0) and +(0:0.1) ..
				(48.64, 11.87) -- cycle
 (49.7, 11.87) -- (49.01, 11.87) .. controls +(180:0.1) and +(-45:0) ..
				(48.88, 11.91) -- (48.76, 12.16) .. controls +(117:0) and +(180:0.1) ..
				(48.85, 12.2) -- (49.52, 12.2) .. controls +(0:0.1) and +(117:0) ..
				(49.64, 12.16) -- (49.78, 11.91) .. controls +(-45:0) and +(0:0.1) ..
				(49.7, 11.87) -- cycle
 (51.15, 11.87) -- (50.07, 11.87) .. controls +(180:0.1) and +(-45:0) ..
				(49.94, 11.91) -- (49.8, 12.16) .. controls +(117:0) and +(180:0.1) ..
				(49.88, 12.2) -- (50.93, 12.2) .. controls +(0:0) and +(135:0) ..
				(51.06, 12.16) -- (51.22, 11.91) .. controls +(-45:0) and +(0:0.1) ..
				(51.15, 11.87) -- cycle
 (36.04, 11.44) -- (34.63, 11.44) .. controls +(180:0.1) and +(-124:0) ..
				(34.53, 11.49) -- (34.69, 11.73) .. controls +(56:0) and +(180:0.1) ..
				(34.86, 11.78) -- (36.23, 11.78) .. controls +(0:0.1) and +(72:0) ..
				(36.33, 11.73) -- (36.2, 11.49) .. controls +(-124:0) and +(0:0.1) ..
				(36.04, 11.44) -- cycle
 (37.14, 11.44) -- (36.47, 11.44) .. controls +(180:0.1) and +(-124:0) ..
				(36.37, 11.49) -- (36.49, 11.73) .. controls +(56:0) and +(180:0.1) ..
				(36.65, 11.78) -- (37.3, 11.78) .. controls +(0:0.1) and +(72:0) ..
				(37.41, 11.73) -- (37.3, 11.49) .. controls +(-124:0) and +(0:0.1) ..
				(37.14, 11.44) -- cycle
 (38.24, 11.44) -- (37.57, 11.44) .. controls +(180:0.1) and +(-108:0) ..
				(37.46, 11.49) -- (37.57, 11.73) .. controls +(72:0) and +(180:0.1) ..
				(37.72, 11.78) -- (38.37, 11.78) .. controls +(0:0.1) and +(72:0) ..
				(38.48, 11.73) -- (38.39, 11.49) .. controls +(-108:0) and +(0:0.1) ..
				(38.24, 11.44) -- cycle
 (39.34, 11.44) -- (38.68, 11.44) .. controls +(180:0.1) and +(-108:0) ..
				(38.56, 11.49) -- (38.64, 11.73) .. controls +(72:0) and +(180:0.1) ..
				(38.79, 11.78) -- (39.44, 11.78) .. controls +(0:0.1) and +(90:0) ..
				(39.56, 11.73) -- (39.49, 11.49) .. controls +(-108:0) and +(0:0.1) ..
				(39.34, 11.44) -- cycle
 (40.44, 11.44) -- (39.78, 11.44) .. controls +(180:0.1) and +(-108:0) ..
				(39.66, 11.49) -- (39.72, 11.73) .. controls +(72:0) and +(180:0.1) ..
				(39.86, 11.78) -- (40.51, 11.78) .. controls +(0:0.1) and +(72:0) ..
				(40.63, 11.73) -- (40.59, 11.49) .. controls +(-108:0) and +(0:0.1) ..
				(40.44, 11.44) -- cycle
 (41.54, 11.44) -- (40.88, 11.44) .. controls +(180:0.1) and +(-90:0) ..
				(40.75, 11.49) -- (40.8, 11.73) .. controls +(90:0) and +(180:0.1) ..
				(40.94, 11.78) -- (41.58, 11.78) .. controls +(0:0.1) and +(90:0) ..
				(41.71, 11.73) -- (41.68, 11.49) .. controls +(-90:0) and +(0:0.1) ..
				(41.54, 11.44) -- cycle
 (42.65, 11.44) -- (41.98, 11.44) .. controls +(180:0.1) and +(-90:0) ..
				(41.85, 11.49) -- (41.87, 11.73) .. controls +(90:0) and +(180:0.1) ..
				(42.01, 11.78) -- (42.65, 11.78) .. controls +(0:0.1) and +(90:0) ..
				(42.78, 11.73) -- (42.78, 11.49) .. controls +(-90:0) and +(0:0.1) ..
				(42.65, 11.44) -- cycle
 (43.75, 11.44) -- (43.08, 11.44) .. controls +(180:0.1) and +(-90:0) ..
				(42.95, 11.49) -- (42.95, 11.73) .. controls +(90:0) and +(180:0.1) ..
				(43.08, 11.78) -- (43.73, 11.78) .. controls +(0:0.1) and +(90:0) ..
				(43.86, 11.73) -- (43.88, 11.49) .. controls +(-90:0) and +(0:0.1) ..
				(43.75, 11.44) -- cycle
 (44.85, 11.44) -- (44.18, 11.44) .. controls +(180:0.1) and +(-72:0) ..
				(44.04, 11.49) -- (44.02, 11.73) .. controls +(90:0) and +(180:0.1) ..
				(44.15, 11.78) -- (44.8, 11.78) .. controls +(0:0.1) and +(108:0) ..
				(44.94, 11.73) -- (44.97, 11.49) .. controls +(-72:0) and +(0:0.1) ..
				(44.85, 11.44) -- cycle
 (45.95, 11.44) -- (45.29, 11.44) .. controls +(180:0.1) and +(-72:0) ..
				(45.14, 11.49) -- (45.1, 11.73) .. controls +(108:0) and +(180:0.1) ..
				(45.22, 11.78) -- (45.87, 11.78) .. controls +(0:0.1) and +(108:0) ..
				(46.01, 11.73) -- (46.07, 11.49) .. controls +(-72:0) and +(0:0.1) ..
				(45.95, 11.44) -- cycle
 (47.05, 11.44) -- (46.39, 11.44) .. controls +(180:0.1) and +(-72:0) ..
				(46.24, 11.49) -- (46.17, 11.73) .. controls +(90:0) and +(180:0.1) ..
				(46.29, 11.78) -- (46.94, 11.78) .. controls +(0:0.1) and +(108:0) ..
				(47.09, 11.73) -- (47.17, 11.49) .. controls +(-72:0) and +(0:0.1) ..
				(47.05, 11.44) -- cycle
 (48.15, 11.44) -- (47.49, 11.44) .. controls +(180:0.1) and +(-72:0) ..
				(47.34, 11.49) -- (47.25, 11.73) .. controls +(108:0) and +(180:0.1) ..
				(47.36, 11.78) -- (48.01, 11.78) .. controls +(0:0.1) and +(108:0) ..
				(48.16, 11.73) -- (48.27, 11.49) .. controls +(-72:0) and +(0:0.1) ..
				(48.15, 11.44) -- cycle
 (49.26, 11.44) -- (48.59, 11.44) .. controls +(180:0.1) and +(-56:0) ..
				(48.43, 11.49) -- (48.33, 11.73) .. controls +(124:0) and +(180:0.1) ..
				(48.43, 11.78) -- (49.08, 11.78) .. controls +(0:0.1) and +(124:0) ..
				(49.24, 11.73) -- (49.36, 11.49) .. controls +(-56:0) and +(0:0.1) ..
				(49.26, 11.44) -- cycle
 (51.42, 11.44) -- (49.72, 11.44) .. controls +(180:0.1) and +(-72:0) ..
				(49.56, 11.49) -- (49.43, 11.73) .. controls +(124:0) and +(180:0.1) ..
				(49.53, 11.78) -- (51.18, 11.78) .. controls +(0:0.1) and +(124:0) ..
				(51.35, 11.73) -- (51.51, 11.49) .. controls +(-56:0) and +(0:0.1) ..
				(51.42, 11.44) -- cycle
 (35.14, 10.96) -- (34.36, 10.96) .. controls +(180:0.1) and +(-121:0.1) ..
				(34.22, 11.04) -- (34.37, 11.26) .. controls +(53:0.1) and +(180:0.1) ..
				(34.6, 11.33) -- (35.36, 11.33) .. controls +(0:0.1) and +(53:0.1) ..
				(35.5, 11.26) -- (35.37, 11.04) .. controls +(-112:0.1) and +(0:0.1) ..
				(35.14, 10.96) -- cycle
 (36.28, 10.96) -- (35.69, 10.96) .. controls +(180:0.1) and +(-112:0.1) ..
				(35.54, 11.04) -- (35.67, 11.26) .. controls +(53:0.1) and +(180:0.1) ..
				(35.89, 11.33) -- (36.48, 11.33) .. controls +(0:0.1) and +(63:0) ..
				(36.62, 11.26) -- (36.51, 11.04) .. controls +(-112:0.1) and +(0:0.1) ..
				(36.28, 10.96) -- cycle
 (37.43, 10.96) -- (36.83, 10.96) .. controls +(180:0.1) and +(-112:0.1) ..
				(36.68, 11.04) -- (36.79, 11.26) .. controls +(63:0) and +(180:0.1) ..
				(37, 11.33) -- (37.59, 11.33) .. controls +(0:0.1) and +(76:0) ..
				(37.74, 11.26) -- (37.64, 11.04) .. controls +(-101:0.1) and +(0:0.1) ..
				(37.43, 10.96) -- cycle
 (38.57, 10.96) -- (37.97, 10.96) .. controls +(180:0.1) and +(-112:0.1) ..
				(37.82, 11.04) -- (37.91, 11.26) .. controls +(76:0) and +(180:0.1) ..
				(38.12, 11.33) -- (38.7, 11.33) .. controls +(0:0.1) and +(63:0) ..
				(38.85, 11.26) -- (38.78, 11.04) .. controls +(-101:0.1) and +(0:0.1) ..
				(38.57, 10.96) -- cycle
 (39.71, 10.96) -- (39.12, 10.96) .. controls +(180:0.1) and +(-101:0.1) ..
				(38.95, 11.04) -- (39.02, 11.26) .. controls +(63:0) and +(180:0.1) ..
				(39.23, 11.33) -- (39.81, 11.33) .. controls +(0:0.1) and +(76:0) ..
				(39.97, 11.26) -- (39.92, 11.04) .. controls +(-101:0.1) and +(0:0.1) ..
				(39.71, 10.96) -- cycle
 (40.86, 10.96) -- (40.26, 10.96) .. controls +(180:0.1) and +(-101:0.1) ..
				(40.09, 11.04) -- (40.14, 11.26) .. controls +(76:0) and +(180:0.1) ..
				(40.34, 11.33) -- (40.92, 11.33) .. controls +(0:0.1) and +(90:0) ..
				(41.09, 11.26) -- (41.06, 11.04) .. controls +(-101:0.1) and +(0:0.1) ..
				(40.86, 10.96) -- cycle
 (42, 10.96) -- (41.4, 10.96) .. controls +(180:0.1) and +(-101:0.1) ..
				(41.23, 11.04) -- (41.26, 11.26) .. controls +(90:0) and +(180:0.1) ..
				(41.45, 11.33) -- (42.03, 11.33) .. controls +(0:0.1) and +(76:0) ..
				(42.2, 11.26) -- (42.19, 11.04) .. controls +(-90:0) and +(0:0.1) ..
				(42, 10.96) -- cycle
 (43.15, 10.96) -- (42.55, 10.96) .. controls +(180:0.1) and +(-90:0) ..
				(42.36, 11.04) -- (42.38, 11.26) .. controls +(90:0) and +(180:0.1) ..
				(42.56, 11.33) -- (43.14, 11.33) .. controls +(0:0.1) and +(90:0) ..
				(43.32, 11.26) -- (43.33, 11.04) .. controls +(-90:0) and +(0:0.1) ..
				(43.15, 10.96) -- cycle
 (44.29, 10.96) -- (43.69, 10.96) .. controls +(180:0.1) and +(-90:0) ..
				(43.5, 11.04) -- (43.49, 11.26) .. controls +(90:0) and +(180:0.1) ..
				(43.67, 11.33) -- (44.25, 11.33) .. controls +(0:0.1) and +(104:0) ..
				(44.44, 11.26) -- (44.47, 11.04) .. controls +(-90:0) and +(0:0.1) ..
				(44.29, 10.96) -- cycle
 (45.43, 10.96) -- (44.83, 10.96) .. controls +(180:0.1) and +(-90:0) ..
				(44.64, 11.04) -- (44.61, 11.26) .. controls +(104:0) and +(180:0.1) ..
				(44.78, 11.33) -- (45.36, 11.33) .. controls +(0:0.1) and +(104:0) ..
				(45.56, 11.26) -- (45.6, 11.04) .. controls +(-79:0.1) and +(0:0.1) ..
				(45.43, 10.96) -- cycle
 (46.58, 10.96) -- (45.98, 10.96) .. controls +(180:0.1) and +(-90:0) ..
				(45.78, 11.04) -- (45.73, 11.26) .. controls +(104:0) and +(180:0.1) ..
				(45.89, 11.33) -- (46.47, 11.33) .. controls +(0:0.1) and +(104:0) ..
				(46.67, 11.26) -- (46.74, 11.04) .. controls +(-79:0.1) and +(0:0.1) ..
				(46.58, 10.96) -- cycle
 (47.72, 10.96) -- (47.12, 10.96) .. controls +(180:0.1) and +(-79:0.1) ..
				(46.91, 11.04) -- (46.84, 11.26) .. controls +(104:0) and +(180:0.1) ..
				(47, 11.33) -- (47.58, 11.33) .. controls +(0:0.1) and +(117:0) ..
				(47.79, 11.26) -- (47.88, 11.04) .. controls +(-79:0.1) and +(0:0.1) ..
				(47.72, 10.96) -- cycle
 (35.73, 10.43) -- (33.92, 10.43) .. controls +(180:0.1) and +(-117:0) ..
				(33.84, 10.47) -- (34.06, 10.79) .. controls +(63:0) and +(180:0.1) ..
				(34.19, 10.83) -- (35.95, 10.83) .. controls +(0:0) and +(63:0) ..
				(36.03, 10.79) -- (35.86, 10.47) .. controls +(-117:0) and +(0:0.1) ..
				(35.73, 10.43) -- cycle
 (36.92, 10.43) -- (36.12, 10.43) .. controls +(180:0) and +(-117:0) ..
				(36.04, 10.47) -- (36.2, 10.79) .. controls +(45:0) and +(180:0.1) ..
				(36.33, 10.83) -- (37.1, 10.83) .. controls +(0:0.1) and +(63:0) ..
				(37.19, 10.79) -- (37.05, 10.47) .. controls +(-117:0) and +(0:0.1) ..
				(36.92, 10.43) -- cycle
 (38.11, 10.43) -- (37.31, 10.43) .. controls +(180:0.1) and +(-117:0) ..
				(37.23, 10.47) -- (37.36, 10.79) .. controls +(63:0) and +(180:0) ..
				(37.48, 10.83) -- (38.26, 10.83) .. controls +(0:0.1) and +(63:0) ..
				(38.34, 10.79) -- (38.23, 10.47) .. controls +(-90:0) and +(0:0.1) ..
				(38.11, 10.43) -- cycle
 (39.3, 10.43) -- (38.5, 10.43) .. controls +(180:0) and +(-90:0) ..
				(38.41, 10.47) -- (38.52, 10.79) .. controls +(63:0) and +(180:0.1) ..
				(38.64, 10.83) -- (39.41, 10.83) .. controls +(0:0.1) and +(63:0) ..
				(39.5, 10.79) -- (39.42, 10.47) .. controls +(-117:0) and +(0:0.1) ..
				(39.3, 10.43) -- cycle
 (45.25, 10.43) -- (39.69, 10.43) .. controls +(180:0.1) and +(-90:0) ..
				(39.59, 10.47) -- (39.67, 10.79) .. controls +(63:0) and +(180:0.1) ..
				(39.78, 10.83) -- (45.18, 10.83) .. controls +(0:0.1) and +(90:0) ..
				(45.29, 10.79) -- (45.35, 10.47) .. controls +(-90:0) and +(0:0.1) ..
				(45.25, 10.43) -- cycle
 (46.44, 10.43) -- (45.64, 10.43) .. controls +(180:0) and +(-63:0) ..
				(45.53, 10.47) -- (45.47, 10.79) .. controls +(90:0) and +(180:0.1) ..
				(45.56, 10.83) -- (46.34, 10.83) .. controls +(0:0.1) and +(90:0) ..
				(46.45, 10.79) -- (46.54, 10.47) .. controls +(-90:0) and +(0:0.1) ..
				(46.44, 10.43) -- cycle
 (47.64, 10.43) -- (46.84, 10.43) .. controls +(180:0.1) and +(-63:0) ..
				(46.72, 10.47) -- (46.63, 10.79) .. controls +(90:0) and +(180:0.1) ..
				(46.73, 10.83) -- (47.5, 10.83) .. controls +(0:0.1) and +(117:0) ..
				(47.62, 10.79) -- (47.73, 10.47) .. controls +(-63:0) and +(0:0.1) ..
				(47.64, 10.43) -- cycle
 (48.86, 10.96) -- (48.26, 10.96) .. controls +(180:0.1) and +(-68:0.1) ..
				(48.05, 11.04) -- (47.96, 11.26) .. controls +(117:0) and +(180:0.1) ..
				(48.11, 11.33) -- (48.69, 11.33) .. controls +(0:0.1) and +(117:0) ..
				(48.91, 11.26) -- (49.01, 11.04) .. controls +(-68:0.1) and +(0:0.1) ..
				(48.86, 10.96) -- cycle
 (51.69, 10.96) -- (49.44, 10.96) .. controls +(180:0.1) and +(-68:0.1) ..
				(49.21, 11.04) -- (49.11, 11.26) .. controls +(117:0) and +(180:0.1) ..
				(49.25, 11.33) -- (51.44, 11.33) .. controls +(0:0.1) and +(127:0.1) ..
				(51.67, 11.26) -- (51.83, 11.04) .. controls +(-68:0.1) and +(0:0.1) ..
				(51.69, 10.96) -- cycle
 (50.57, 10.66) -- (49.53, 10.66) .. controls +(180:0.1) and +(-72:0) ..
				(49.4, 10.71) -- (49.36, 10.79) .. controls +(117:0) and +(180:0.1) ..
				(49.45, 10.83) -- (50.47, 10.83) .. controls +(0:0.1) and +(135:0) ..
				(50.6, 10.79) -- (50.64, 10.71) .. controls +(-56:0) and +(0:0) ..
				(50.57, 10.66) -- cycle
 (50.7, 10.43) -- (49.65, 10.43) .. controls +(180:0.1) and +(-63:0) ..
				(49.52, 10.47) -- (49.48, 10.56) .. controls +(117:0) and +(180:0.1) ..
				(49.56, 10.6) -- (50.6, 10.6) .. controls +(0:0.1) and +(117:0) ..
				(50.73, 10.56) -- (50.78, 10.47) .. controls +(-45:0) and +(0:0.1) ..
				(50.7, 10.43) -- cycle
 (52.14, 10.43) -- (51.08, 10.43) .. controls +(180:0.1) and +(-63:0) ..
				(50.95, 10.47) -- (50.9, 10.56) .. controls +(135:0) and +(180:0.1) ..
				(50.98, 10.6) -- (52.02, 10.6) .. controls +(0:0.1) and +(117:0) ..
				(52.15, 10.56) -- (52.21, 10.47) .. controls +(-45:0) and +(0:0) ..
				(52.14, 10.43) -- cycle
 (49.27, 10.43) -- (48.21, 10.43) .. controls +(180:0.1) and +(-63:0) ..
				(48.09, 10.47) -- (48.06, 10.56) .. controls +(117:0) and +(180:0.1) ..
				(48.15, 10.6) -- (49.19, 10.6) .. controls +(0:0.1) and +(117:0) ..
				(49.31, 10.56) -- (49.35, 10.47) .. controls +(-63:0) and +(0:0.1) ..
				(49.27, 10.43) -- cycle;
		\end{scope}

		\path[fill=cd3ccd2] (50.98, 13.1) -- (50.97, 13.07) .. controls +(125:0.1) and +(0:0.3) ..
			(50.33, 13.2) -- (35.56, 13.2) .. controls +(180:0.3) and +(55:0.1) ..
			(34.92, 13.07) -- (34.91, 13.1) -- cycle;
		\begin{scope}
			\path[fill=c4e4c4d] (35.95, 12.7) -- (35.43, 12.7) .. controls +(180:0) and +(-117:0) ..
				(35.36, 12.73) -- (35.42, 12.82) .. controls +(45:0) and +(180:0) ..
				(35.54, 12.86) -- (36.05, 12.86) .. controls +(0:0.1) and +(63:0) ..
				(36.12, 12.82) -- (36.06, 12.73) .. controls +(-117:0) and +(0:0) ..
				(35.95, 12.7) -- cycle
 (36.87, 12.7) -- (36.35, 12.7) .. controls +(180:0.1) and +(-117:0) ..
				(36.27, 12.73) -- (36.33, 12.82) .. controls +(63:0) and +(180:0) ..
				(36.44, 12.86) -- (36.95, 12.86) .. controls +(0:0) and +(63:0) ..
				(37.03, 12.82) -- (36.98, 12.73) .. controls +(-117:0) and +(0:0.1) ..
				(36.87, 12.7) -- cycle
 (37.78, 12.7) -- (37.26, 12.7) .. controls +(180:0) and +(-117:0) ..
				(37.19, 12.73) -- (37.23, 12.82) .. controls +(63:0) and +(180:0.1) ..
				(37.34, 12.86) -- (37.85, 12.86) .. controls +(0:0) and +(63:0) ..
				(37.93, 12.82) -- (37.89, 12.73) .. controls +(-117:0) and +(0:0) ..
				(37.78, 12.7) -- cycle
 (38.7, 12.7) -- (38.18, 12.7) .. controls +(180:0) and +(-117:0) ..
				(38.1, 12.73) -- (38.14, 12.82) .. controls +(63:0) and +(180:0.1) ..
				(38.25, 12.86) -- (38.76, 12.86) .. controls +(0:0.1) and +(90:0) ..
				(38.84, 12.82) -- (38.8, 12.73) .. controls +(-117:0) and +(0:0) ..
				(38.7, 12.7) -- cycle
 (39.61, 12.7) -- (39.09, 12.7) .. controls +(180:0.1) and +(-117:0) ..
				(39.01, 12.73) -- (39.04, 12.82) .. controls +(63:0) and +(180:0) ..
				(39.15, 12.86) -- (39.66, 12.86) .. controls +(0:0.1) and +(63:0) ..
				(39.74, 12.82) -- (39.71, 12.73) .. controls +(-90:0) and +(0:0) ..
				(39.61, 12.7) -- cycle
 (40.53, 12.7) -- (40.01, 12.7) .. controls +(180:0) and +(-90:0) ..
				(39.92, 12.73) -- (39.95, 12.82) .. controls +(90:0) and +(180:0) ..
				(40.05, 12.86) -- (40.56, 12.86) .. controls +(0:0) and +(90:0) ..
				(40.65, 12.82) -- (40.63, 12.73) .. controls +(-117:0) and +(0:0) ..
				(40.53, 12.7) -- cycle
 (41.44, 12.7) -- (40.92, 12.7) .. controls +(180:0.1) and +(-117:0) ..
				(40.84, 12.73) -- (40.85, 12.82) .. controls +(63:0) and +(180:0.1) ..
				(40.95, 12.86) -- (41.46, 12.86) .. controls +(0:0) and +(90:0) ..
				(41.55, 12.82) -- (41.54, 12.73) .. controls +(-90:0) and +(0:0.1) ..
				(41.44, 12.7) -- cycle
 (42.36, 12.7) -- (41.84, 12.7) .. controls +(180:0.1) and +(-90:0) ..
				(41.75, 12.73) -- (41.76, 12.82) .. controls +(90:0) and +(180:0.1) ..
				(41.85, 12.86) -- (42.37, 12.86) .. controls +(0:0.1) and +(90:0) ..
				(42.46, 12.82) -- (42.45, 12.73) .. controls +(-90:0) and +(0:0) ..
				(42.36, 12.7) -- cycle
 (43.27, 12.7) -- (42.75, 12.7) .. controls +(180:0) and +(-90:0) ..
				(42.66, 12.73) -- (42.66, 12.82) .. controls +(90:0) and +(180:0) ..
				(42.76, 12.86) -- (43.27, 12.86) .. controls +(0:0) and +(90:0) ..
				(43.36, 12.82) -- (43.36, 12.73) .. controls +(-63:0) and +(0:0) ..
				(43.27, 12.7) -- cycle
 (44.19, 12.7) -- (43.67, 12.7) .. controls +(180:0.1) and +(-63:0) ..
				(43.57, 12.73) -- (43.57, 12.82) .. controls +(90:0) and +(180:0) ..
				(43.66, 12.86) -- (44.17, 12.86) .. controls +(0:0) and +(117:0) ..
				(44.27, 12.82) -- (44.28, 12.73) .. controls +(-90:0) and +(0:0.1) ..
				(44.19, 12.7) -- cycle
 (45.1, 12.7) -- (44.59, 12.7) .. controls +(180:0.1) and +(-90:0) ..
				(44.49, 12.73) -- (44.47, 12.82) .. controls +(90:0) and +(180:0.1) ..
				(44.56, 12.86) -- (45.07, 12.86) .. controls +(0:0) and +(90:0) ..
				(45.17, 12.82) -- (45.19, 12.73) .. controls +(-90:0) and +(0:0) ..
				(45.1, 12.7) -- cycle
 (46.02, 12.7) -- (45.5, 12.7) .. controls +(180:0) and +(-90:0) ..
				(45.4, 12.73) -- (45.38, 12.82) .. controls +(90:0) and +(180:0.1) ..
				(45.46, 12.86) -- (45.98, 12.86) .. controls +(0:0.1) and +(117:0) ..
				(46.08, 12.82) -- (46.1, 12.73) .. controls +(-63:0) and +(0:0) ..
				(46.02, 12.7) -- cycle
 (46.93, 12.7) -- (46.42, 12.7) .. controls +(180:0.1) and +(-63:0) ..
				(46.31, 12.73) -- (46.28, 12.82) .. controls +(90:0) and +(180:0) ..
				(46.37, 12.86) -- (46.88, 12.86) .. controls +(0:0) and +(117:0) ..
				(46.98, 12.82) -- (47.02, 12.73) .. controls +(-90:0) and +(0:0.1) ..
				(46.93, 12.7) -- cycle
 (47.85, 12.7) -- (47.33, 12.7) .. controls +(180:0) and +(-90:0) ..
				(47.23, 12.73) -- (47.19, 12.82) .. controls +(117:0) and +(180:0.1) ..
				(47.27, 12.86) -- (47.78, 12.86) .. controls +(0:0) and +(117:0) ..
				(47.89, 12.82) -- (47.93, 12.73) .. controls +(-63:0) and +(0:0) ..
				(47.85, 12.7) -- cycle
 (48.76, 12.7) -- (48.25, 12.7) .. controls +(180:0.1) and +(-63:0) ..
				(48.14, 12.73) -- (48.09, 12.82) .. controls +(90:0) and +(180:0.1) ..
				(48.17, 12.86) -- (48.68, 12.86) .. controls +(0:0) and +(117:0) ..
				(48.79, 12.82) -- (48.84, 12.73) .. controls +(-63:0) and +(0:0.1) ..
				(48.76, 12.7) -- cycle
 (49.68, 12.7) -- (49.16, 12.7) .. controls +(180:0) and +(-63:0) ..
				(49.05, 12.73) -- (49, 12.82) .. controls +(117:0) and +(180:0) ..
				(49.07, 12.86) -- (49.58, 12.86) .. controls +(0:0.1) and +(117:0) ..
				(49.7, 12.82) -- (49.75, 12.73) .. controls +(-63:0) and +(0:0) ..
				(49.68, 12.7) -- cycle
 (50.6, 12.7) -- (50.08, 12.7) .. controls +(180:0) and +(-63:0) ..
				(49.96, 12.73) -- (49.9, 12.82) .. controls +(117:0) and +(180:0) ..
				(49.98, 12.86) -- (50.49, 12.86) .. controls +(0:0) and +(117:0) ..
				(50.6, 12.82) -- (50.67, 12.73) .. controls +(-63:0) and +(0:0) ..
				(50.6, 12.7) -- cycle
 (35.86, 12.33) -- (35.18, 12.33) .. controls +(180:0) and +(-117:0) ..
				(35.11, 12.36) -- (35.27, 12.59) .. controls +(63:0) and +(180:0.1) ..
				(35.38, 12.63) -- (36.05, 12.63) .. controls +(0:0.1) and +(63:0) ..
				(36.12, 12.59) -- (35.98, 12.36) .. controls +(-117:0) and +(0:0.1) ..
				(35.86, 12.33) -- cycle
 (36.89, 12.33) -- (36.21, 12.33) .. controls +(180:0.1) and +(-117:0) ..
				(36.13, 12.36) -- (36.27, 12.59) .. controls +(63:0) and +(180:0.1) ..
				(36.38, 12.63) -- (37.04, 12.63) .. controls +(0:0.1) and +(63:0) ..
				(37.12, 12.59) -- (37, 12.36) .. controls +(-117:0) and +(0:0) ..
				(36.89, 12.33) -- cycle
 (37.91, 12.33) -- (37.23, 12.33) .. controls +(180:0) and +(-117:0) ..
				(37.15, 12.36) -- (37.27, 12.59) .. controls +(63:0) and +(180:0.1) ..
				(37.38, 12.63) -- (38.04, 12.63) .. controls +(0:0.1) and +(63:0) ..
				(38.12, 12.59) -- (38.02, 12.36) .. controls +(-117:0) and +(0:0.1) ..
				(37.91, 12.33) -- cycle
 (38.94, 12.33) -- (38.25, 12.33) .. controls +(180:0) and +(-90:0) ..
				(38.17, 12.36) -- (38.27, 12.59) .. controls +(63:0) and +(180:0) ..
				(38.37, 12.63) -- (39.04, 12.63) .. controls +(0:0.1) and +(90:0) ..
				(39.12, 12.59) -- (39.04, 12.36) .. controls +(-117:0) and +(0:0.1) ..
				(38.94, 12.33) -- cycle
 (39.96, 12.33) -- (39.28, 12.33) .. controls +(180:0.1) and +(-117:0) ..
				(39.2, 12.36) -- (39.27, 12.59) .. controls +(63:0) and +(180:0) ..
				(39.37, 12.63) -- (40.04, 12.63) .. controls +(0:0.1) and +(90:0) ..
				(40.12, 12.59) -- (40.06, 12.36) .. controls +(-90:0) and +(0:0) ..
				(39.96, 12.33) -- cycle
 (40.98, 12.33) -- (40.3, 12.33) .. controls +(180:0) and +(-117:0) ..
				(40.22, 12.36) -- (40.27, 12.59) .. controls +(90:0) and +(180:0) ..
				(40.37, 12.63) -- (41.03, 12.63) .. controls +(0:0) and +(90:0) ..
				(41.12, 12.59) -- (41.08, 12.36) .. controls +(-90:0) and +(0:0.1) ..
				(40.98, 12.33) -- cycle
 (42.01, 12.33) -- (41.32, 12.33) .. controls +(180:0) and +(-117:0) ..
				(41.24, 12.36) -- (41.27, 12.59) .. controls +(90:0) and +(180:0) ..
				(41.37, 12.63) -- (42.03, 12.63) .. controls +(0:0) and +(90:0) ..
				(42.12, 12.59) -- (42.1, 12.36) .. controls +(-90:0) and +(0:0.1) ..
				(42.01, 12.33) -- cycle
 (43.03, 12.33) -- (42.35, 12.33) .. controls +(180:0.1) and +(-90:0) ..
				(42.26, 12.36) -- (42.27, 12.59) .. controls +(90:0) and +(180:0) ..
				(42.36, 12.63) -- (43.03, 12.63) .. controls +(0:0) and +(90:0) ..
				(43.12, 12.59) -- (43.12, 12.36) .. controls +(-90:0) and +(0:0) ..
				(43.03, 12.33) -- cycle
 (44.05, 12.33) -- (43.37, 12.33) .. controls +(180:0) and +(-90:0) ..
				(43.28, 12.36) -- (43.27, 12.59) .. controls +(90:0) and +(180:0) ..
				(43.36, 12.63) -- (44.03, 12.63) .. controls +(0:0) and +(90:0) ..
				(44.12, 12.59) -- (44.14, 12.36) .. controls +(-63:0) and +(0:0.1) ..
				(44.05, 12.33) -- cycle
 (45.08, 12.33) -- (44.4, 12.33) .. controls +(180:0.1) and +(-90:0) ..
				(44.3, 12.36) -- (44.27, 12.59) .. controls +(90:0) and +(180:0) ..
				(44.36, 12.63) -- (45.02, 12.63) .. controls +(0:0) and +(90:0) ..
				(45.12, 12.59) -- (45.16, 12.36) .. controls +(-63:0) and +(0:0.1) ..
				(45.08, 12.33) -- cycle
 (46.1, 12.33) -- (45.42, 12.33) .. controls +(180:0.1) and +(-90:0) ..
				(45.32, 12.36) -- (45.27, 12.59) .. controls +(90:0) and +(180:0) ..
				(45.36, 12.63) -- (46.02, 12.63) .. controls +(0:0) and +(90:0) ..
				(46.12, 12.59) -- (46.19, 12.36) .. controls +(-90:0) and +(0:0) ..
				(46.1, 12.33) -- cycle
 (47.13, 12.33) -- (46.44, 12.33) .. controls +(180:0) and +(-63:0) ..
				(46.34, 12.36) -- (46.27, 12.59) .. controls +(90:0) and +(180:0.1) ..
				(46.35, 12.63) -- (47.02, 12.63) .. controls +(0:0) and +(117:0) ..
				(47.12, 12.59) -- (47.21, 12.36) .. controls +(-90:0) and +(0:0) ..
				(47.13, 12.33) -- cycle
 (48.15, 12.33) -- (47.47, 12.33) .. controls +(180:0) and +(-63:0) ..
				(47.36, 12.36) -- (47.27, 12.59) .. controls +(90:0) and +(180:0.1) ..
				(47.35, 12.63) -- (48.02, 12.63) .. controls +(0:0) and +(117:0) ..
				(48.12, 12.59) -- (48.23, 12.36) .. controls +(-90:0) and +(0:0.1) ..
				(48.15, 12.33) -- cycle
 (49.17, 12.33) -- (48.49, 12.33) .. controls +(180:0.1) and +(-63:0) ..
				(48.38, 12.36) -- (48.27, 12.59) .. controls +(117:0) and +(180:0.1) ..
				(48.35, 12.63) -- (49.01, 12.63) .. controls +(0:0.1) and +(117:0) ..
				(49.12, 12.59) -- (49.25, 12.36) .. controls +(-63:0) and +(0:0) ..
				(49.17, 12.33) -- cycle
 (50.85, 12.33) -- (49.51, 12.33) .. controls +(180:0) and +(-63:0) ..
				(49.4, 12.36) -- (49.27, 12.59) .. controls +(117:0) and +(180:0.1) ..
				(49.35, 12.63) -- (50.65, 12.63) .. controls +(0:0.1) and +(117:0) ..
				(50.76, 12.59) -- (50.92, 12.36) .. controls +(-63:0) and +(0:0) ..
				(50.85, 12.33) -- cycle
 (35.9, 11.92) -- (34.92, 11.92) .. controls +(180:0.1) and +(-108:0) ..
				(34.84, 11.96) -- (35.01, 12.2) .. controls +(63:0) and +(180:0.1) ..
				(35.13, 12.24) -- (36.09, 12.24) .. controls +(0:0.1) and +(63:0) ..
				(36.17, 12.2) -- (36.03, 11.96) .. controls +(-108:0) and +(0:0.1) ..
				(35.9, 11.92) -- cycle
 (36.96, 11.92) -- (36.27, 11.92) .. controls +(180:0.1) and +(-108:0) ..
				(36.19, 11.96) -- (36.33, 12.2) .. controls +(63:0) and +(180:0.1) ..
				(36.45, 12.24) -- (37.12, 12.24) .. controls +(0:0.1) and +(63:0) ..
				(37.2, 12.2) -- (37.09, 11.96) .. controls +(-108:0) and +(0:0.1) ..
				(36.96, 11.92) -- cycle
 (38.02, 11.92) -- (37.33, 11.92) .. controls +(180:0.1) and +(-108:0) ..
				(37.25, 11.96) -- (37.36, 12.2) .. controls +(63:0) and +(180:0.1) ..
				(37.48, 12.24) -- (38.15, 12.24) .. controls +(0:0.1) and +(63:0) ..
				(38.24, 12.2) -- (38.14, 11.96) .. controls +(-108:0) and +(0:0.1) ..
				(38.02, 11.92) -- cycle
 (39.08, 11.92) -- (38.39, 11.92) .. controls +(180:0.1) and +(-108:0) ..
				(38.3, 11.96) -- (38.4, 12.2) .. controls +(90:0) and +(180:0) ..
				(38.51, 12.24) -- (39.18, 12.24) .. controls +(0:0.1) and +(63:0) ..
				(39.27, 12.2) -- (39.2, 11.96) .. controls +(-108:0) and +(0:0.1) ..
				(39.08, 11.92) -- cycle
 (40.14, 11.92) -- (39.45, 11.92) .. controls +(180:0.1) and +(-108:0) ..
				(39.36, 11.96) -- (39.43, 12.2) .. controls +(63:0) and +(180:0) ..
				(39.54, 12.24) -- (40.22, 12.24) .. controls +(0:0.1) and +(90:0) ..
				(40.31, 12.2) -- (40.25, 11.96) .. controls +(-90:0) and +(0:0.1) ..
				(40.14, 11.92) -- cycle
 (41.2, 11.92) -- (40.51, 11.92) .. controls +(180:0.1) and +(-90:0) ..
				(40.41, 11.96) -- (40.47, 12.2) .. controls +(90:0) and +(180:0) ..
				(40.57, 12.24) -- (41.25, 12.24) .. controls +(0:0) and +(63:0) ..
				(41.34, 12.2) -- (41.31, 11.96) .. controls +(-90:0) and +(0:0.1) ..
				(41.2, 11.92) -- cycle
 (42.26, 11.92) -- (41.57, 11.92) .. controls +(180:0.1) and +(-90:0) ..
				(41.47, 11.96) -- (41.5, 12.2) .. controls +(90:0) and +(180:0.1) ..
				(41.61, 12.24) -- (42.28, 12.24) .. controls +(0:0) and +(90:0) ..
				(42.38, 12.2) -- (42.37, 11.96) .. controls +(-90:0) and +(0:0.1) ..
				(42.26, 11.92) -- cycle
 (43.32, 11.92) -- (42.63, 11.92) .. controls +(180:0.1) and +(-90:0) ..
				(42.53, 11.96) -- (42.54, 12.2) .. controls +(90:0) and +(180:0.1) ..
				(42.64, 12.24) -- (43.31, 12.24) .. controls +(0:0.1) and +(90:0) ..
				(43.41, 12.2) -- (43.42, 11.96) .. controls +(-90:0) and +(0:0.1) ..
				(43.32, 11.92) -- cycle
 (44.38, 11.92) -- (43.69, 11.92) .. controls +(180:0.1) and +(-90:0) ..
				(43.58, 11.96) -- (43.57, 12.2) .. controls +(90:0) and +(180:0.1) ..
				(43.67, 12.24) -- (44.34, 12.24) .. controls +(0:0.1) and +(117:0) ..
				(44.45, 12.2) -- (44.48, 11.96) .. controls +(-90:0) and +(0:0.1) ..
				(44.38, 11.92) -- cycle
 (45.44, 11.92) -- (44.75, 11.92) .. controls +(180:0.1) and +(-90:0) ..
				(44.64, 11.96) -- (44.6, 12.2) .. controls +(90:0) and +(180:0.1) ..
				(44.7, 12.24) -- (45.37, 12.24) .. controls +(0:0.1) and +(90:0) ..
				(45.48, 12.2) -- (45.53, 11.96) .. controls +(-72:0) and +(0:0.1) ..
				(45.44, 11.92) -- cycle
 (46.5, 11.92) -- (45.81, 11.92) .. controls +(180:0.1) and +(-90:0) ..
				(45.7, 11.96) -- (45.64, 12.2) .. controls +(117:0) and +(180:0) ..
				(45.73, 12.24) -- (46.4, 12.24) .. controls +(0:0.1) and +(117:0) ..
				(46.52, 12.2) -- (46.59, 11.96) .. controls +(-72:0) and +(0:0.1) ..
				(46.5, 11.92) -- cycle
 (47.56, 11.92) -- (46.87, 11.92) .. controls +(180:0.1) and +(-72:0) ..
				(46.75, 11.96) -- (46.67, 12.2) .. controls +(90:0) and +(180:0) ..
				(46.76, 12.24) -- (47.44, 12.24) .. controls +(0:0.1) and +(117:0) ..
				(47.55, 12.2) -- (47.65, 11.96) .. controls +(-72:0) and +(0:0.1) ..
				(47.56, 11.92) -- cycle
 (48.62, 11.92) -- (47.93, 11.92) .. controls +(180:0.1) and +(-72:0) ..
				(47.81, 11.96) -- (47.71, 12.2) .. controls +(117:0) and +(180:0) ..
				(47.79, 12.24) -- (48.47, 12.24) .. controls +(0:0.1) and +(117:0) ..
				(48.59, 12.2) -- (48.7, 11.96) .. controls +(-72:0) and +(0:0.1) ..
				(48.62, 11.92) -- cycle
 (49.68, 11.92) -- (48.99, 11.92) .. controls +(180:0.1) and +(-72:0) ..
				(48.86, 11.96) -- (48.74, 12.2) .. controls +(117:0) and +(180:0.1) ..
				(48.83, 12.24) -- (49.5, 12.24) .. controls +(0:0) and +(117:0) ..
				(49.62, 12.2) -- (49.76, 11.96) .. controls +(-72:0) and +(0:0) ..
				(49.68, 11.92) -- cycle
 (51.12, 11.92) -- (50.05, 11.92) .. controls +(180:0.1) and +(-72:0) ..
				(49.92, 11.96) -- (49.78, 12.2) .. controls +(117:0) and +(180:0.1) ..
				(49.86, 12.24) -- (50.9, 12.24) .. controls +(0:0.1) and +(135:0) ..
				(51.03, 12.2) -- (51.2, 11.96) .. controls +(-72:0) and +(0:0.1) ..
				(51.12, 11.92) -- cycle
 (36.06, 11.48) -- (34.66, 11.48) .. controls +(180:0.1) and +(-124:0) ..
				(34.56, 11.54) -- (34.72, 11.78) .. controls +(45:0) and +(180:0.1) ..
				(34.88, 11.83) -- (36.25, 11.83) .. controls +(0:0.1) and +(45:0) ..
				(36.35, 11.78) -- (36.22, 11.54) .. controls +(-108:0) and +(0:0.1) ..
				(36.06, 11.48) -- cycle
 (37.16, 11.48) -- (36.5, 11.48) .. controls +(180:0.1) and +(-124:0) ..
				(36.39, 11.54) -- (36.52, 11.78) .. controls +(63:0) and +(180:0.1) ..
				(36.67, 11.83) -- (37.32, 11.83) .. controls +(0:0.1) and +(63:0) ..
				(37.43, 11.78) -- (37.32, 11.54) .. controls +(-124:0) and +(0:0.1) ..
				(37.16, 11.48) -- cycle
 (38.26, 11.48) -- (37.59, 11.48) .. controls +(180:0.1) and +(-108:0) ..
				(37.48, 11.54) -- (37.59, 11.78) .. controls +(63:0) and +(180:0.1) ..
				(37.74, 11.83) -- (38.39, 11.83) .. controls +(0:0.1) and +(63:0) ..
				(38.5, 11.78) -- (38.41, 11.54) .. controls +(-108:0) and +(0:0.1) ..
				(38.26, 11.48) -- cycle
 (39.35, 11.48) -- (38.69, 11.48) .. controls +(180:0.1) and +(-108:0) ..
				(38.58, 11.54) -- (38.66, 11.78) .. controls +(63:0) and +(180:0.1) ..
				(38.81, 11.83) -- (39.45, 11.83) .. controls +(0:0.1) and +(63:0) ..
				(39.57, 11.78) -- (39.5, 11.54) .. controls +(-108:0) and +(0:0.1) ..
				(39.35, 11.48) -- cycle
 (40.45, 11.48) -- (39.79, 11.48) .. controls +(180:0.1) and +(-108:0) ..
				(39.67, 11.54) -- (39.73, 11.78) .. controls +(63:0) and +(180:0.1) ..
				(39.88, 11.83) -- (40.52, 11.83) .. controls +(0:0.1) and +(63:0) ..
				(40.64, 11.78) -- (40.6, 11.54) .. controls +(-108:0) and +(0:0.1) ..
				(40.45, 11.48) -- cycle
 (41.55, 11.48) -- (40.89, 11.48) .. controls +(180:0.1) and +(-90:0) ..
				(40.76, 11.54) -- (40.8, 11.78) .. controls +(63:0) and +(180:0.1) ..
				(40.94, 11.83) -- (41.59, 11.83) .. controls +(0:0.1) and +(63:0) ..
				(41.71, 11.78) -- (41.69, 11.54) .. controls +(-90:0) and +(0:0.1) ..
				(41.55, 11.48) -- cycle
 (42.65, 11.48) -- (41.98, 11.48) .. controls +(180:0.1) and +(-90:0) ..
				(41.85, 11.54) -- (41.88, 11.78) .. controls +(90:0) and +(180:0.1) ..
				(42.01, 11.83) -- (42.65, 11.83) .. controls +(0:0.1) and +(90:0) ..
				(42.78, 11.78) -- (42.78, 11.54) .. controls +(-90:0) and +(0:0.1) ..
				(42.65, 11.48) -- cycle
 (43.74, 11.48) -- (43.08, 11.48) .. controls +(180:0.1) and +(-90:0) ..
				(42.95, 11.54) -- (42.95, 11.78) .. controls +(90:0) and +(180:0.1) ..
				(43.08, 11.83) -- (43.72, 11.83) .. controls +(0:0.1) and +(117:0) ..
				(43.86, 11.78) -- (43.87, 11.54) .. controls +(-72:0) and +(0:0.1) ..
				(43.74, 11.48) -- cycle
 (44.84, 11.48) -- (44.18, 11.48) .. controls +(180:0.1) and +(-90:0) ..
				(44.04, 11.54) -- (44.02, 11.78) .. controls +(90:0) and +(180:0.1) ..
				(44.15, 11.83) -- (44.79, 11.83) .. controls +(0:0.1) and +(117:0) ..
				(44.93, 11.78) -- (44.97, 11.54) .. controls +(-90:0) and +(0:0.1) ..
				(44.84, 11.48) -- cycle
 (45.94, 11.48) -- (45.28, 11.48) .. controls +(180:0.1) and +(-72:0) ..
				(45.13, 11.54) -- (45.09, 11.78) .. controls +(90:0) and +(180:0.1) ..
				(45.21, 11.83) -- (45.86, 11.83) .. controls +(0:0.1) and +(117:0) ..
				(46, 11.78) -- (46.06, 11.54) .. controls +(-72:0) and +(0:0.1) ..
				(45.94, 11.48) -- cycle
 (47.04, 11.48) -- (46.37, 11.48) .. controls +(180:0.1) and +(-90:0) ..
				(46.23, 11.54) -- (46.16, 11.78) .. controls +(90:0) and +(180:0.1) ..
				(46.28, 11.83) -- (46.92, 11.83) .. controls +(0:0.1) and +(117:0) ..
				(47.07, 11.78) -- (47.15, 11.54) .. controls +(-72:0) and +(0:0.1) ..
				(47.04, 11.48) -- cycle
 (48.13, 11.48) -- (47.47, 11.48) .. controls +(180:0.1) and +(-72:0) ..
				(47.32, 11.54) -- (47.23, 11.78) .. controls +(90:0) and +(180:0.1) ..
				(47.35, 11.83) -- (47.99, 11.83) .. controls +(0:0.1) and +(117:0) ..
				(48.14, 11.78) -- (48.25, 11.54) .. controls +(-72:0) and +(0:0.1) ..
				(48.13, 11.48) -- cycle
 (49.23, 11.48) -- (48.57, 11.48) .. controls +(180:0.1) and +(-56:0) ..
				(48.41, 11.54) -- (48.31, 11.78) .. controls +(135:0) and +(180:0.1) ..
				(48.42, 11.83) -- (49.06, 11.83) .. controls +(0:0.1) and +(135:0) ..
				(49.22, 11.78) -- (49.34, 11.54) .. controls +(-72:0) and +(0:0.1) ..
				(49.23, 11.48) -- cycle
 (51.39, 11.48) -- (49.7, 11.48) .. controls +(180:0.1) and +(-56:0) ..
				(49.53, 11.54) -- (49.41, 11.78) .. controls +(135:0) and +(180:0.1) ..
				(49.51, 11.83) -- (51.15, 11.83) .. controls +(0:0.1) and +(135:0) ..
				(51.32, 11.78) -- (51.48, 11.54) .. controls +(-56:0) and +(0:0.1) ..
				(51.38, 11.48) -- cycle
 (35.17, 11.01) -- (34.39, 11.01) .. controls +(180:0.1) and +(-117:0) ..
				(34.25, 11.08) -- (34.4, 11.31) .. controls +(53:0.1) and +(180:0.1) ..
				(34.63, 11.38) -- (35.39, 11.38) .. controls +(0:0.1) and +(63:0) ..
				(35.53, 11.31) -- (35.4, 11.08) .. controls +(-117:0) and +(0:0.1) ..
				(35.17, 11.01) -- cycle
 (36.31, 11.01) -- (35.71, 11.01) .. controls +(180:0.1) and +(-117:0) ..
				(35.57, 11.08) -- (35.7, 11.31) .. controls +(63:0) and +(180:0.1) ..
				(35.92, 11.38) -- (36.5, 11.38) .. controls +(0:0.1) and +(63:0) ..
				(36.64, 11.31) -- (36.53, 11.08) .. controls +(-117:0) and +(0:0.1) ..
				(36.31, 11.01) -- cycle
 (37.45, 11.01) -- (36.85, 11.01) .. controls +(180:0.1) and +(-117:0) ..
				(36.7, 11.08) -- (36.81, 11.31) .. controls +(63:0) and +(180:0.1) ..
				(37.03, 11.38) -- (37.61, 11.38) .. controls +(0:0.1) and +(76:0) ..
				(37.76, 11.31) -- (37.66, 11.08) .. controls +(-104:0) and +(0:0.1) ..
				(37.45, 11.01) -- cycle
 (38.59, 11.01) -- (37.99, 11.01) .. controls +(180:0.1) and +(-117:0) ..
				(37.84, 11.08) -- (37.93, 11.31) .. controls +(76:0) and +(180:0.1) ..
				(38.13, 11.38) -- (38.71, 11.38) .. controls +(0:0.1) and +(76:0) ..
				(38.87, 11.31) -- (38.8, 11.08) .. controls +(-117:0) and +(0:0.1) ..
				(38.59, 11.01) -- cycle
 (39.73, 11.01) -- (39.13, 11.01) .. controls +(180:0.1) and +(-104:0) ..
				(38.97, 11.08) -- (39.04, 11.31) .. controls +(76:0) and +(180:0.1) ..
				(39.24, 11.38) -- (39.82, 11.38) .. controls +(0:0.1) and +(76:0) ..
				(39.98, 11.31) -- (39.93, 11.08) .. controls +(-104:0) and +(0:0.1) ..
				(39.73, 11.01) -- cycle
 (40.87, 11.01) -- (40.27, 11.01) .. controls +(180:0.1) and +(-104:0) ..
				(40.1, 11.08) -- (40.15, 11.31) .. controls +(76:0) and +(180:0.1) ..
				(40.35, 11.38) -- (40.93, 11.38) .. controls +(0:0.1) and +(76:0) ..
				(41.09, 11.31) -- (41.06, 11.08) .. controls +(-90:0) and +(0:0.1) ..
				(40.87, 11.01) -- cycle
 (42.01, 11.01) -- (41.41, 11.01) .. controls +(180:0.1) and +(-90:0) ..
				(41.23, 11.08) -- (41.26, 11.31) .. controls +(76:0) and +(180:0.1) ..
				(41.45, 11.38) -- (42.03, 11.38) .. controls +(0:0.1) and +(90:0) ..
				(42.21, 11.31) -- (42.19, 11.08) .. controls +(-90:0) and +(0:0.1) ..
				(42.01, 11.01) -- cycle
 (43.14, 11.01) -- (42.55, 11.01) .. controls +(180:0.1) and +(-90:0) ..
				(42.37, 11.08) -- (42.38, 11.31) .. controls +(90:0) and +(180:0.1) ..
				(42.56, 11.38) -- (43.14, 11.38) .. controls +(0:0.1) and +(90:0) ..
				(43.32, 11.31) -- (43.33, 11.08) .. controls +(-90:0) and +(0:0.1) ..
				(43.14, 11.01) -- cycle
 (44.28, 11.01) -- (43.69, 11.01) .. controls +(180:0.1) and +(-90:0) ..
				(43.5, 11.08) -- (43.49, 11.31) .. controls +(90:0) and +(180:0.1) ..
				(43.67, 11.38) -- (44.25, 11.38) .. controls +(0:0.1) and +(90:0) ..
				(44.43, 11.31) -- (44.46, 11.08) .. controls +(-90:0) and +(0:0.1) ..
				(44.28, 11.01) -- cycle
 (45.42, 11.01) -- (44.83, 11.01) .. controls +(180:0.1) and +(-76:0) ..
				(44.63, 11.08) -- (44.6, 11.31) .. controls +(90:0) and +(180:0.1) ..
				(44.77, 11.38) -- (45.35, 11.38) .. controls +(0:0.1) and +(104:0) ..
				(45.55, 11.31) -- (45.59, 11.08) .. controls +(-76:0) and +(0:0.1) ..
				(45.42, 11.01) -- cycle
 (46.56, 11.01) -- (45.97, 11.01) .. controls +(180:0.1) and +(-76:0) ..
				(45.76, 11.08) -- (45.72, 11.31) .. controls +(104:0) and +(180:0.1) ..
				(45.88, 11.38) -- (46.46, 11.38) .. controls +(0:0.1) and +(104:0) ..
				(46.66, 11.31) -- (46.72, 11.08) .. controls +(-63:0) and +(0:0.1) ..
				(46.56, 11.01) -- cycle
 (47.7, 11.01) -- (47.1, 11.01) .. controls +(180:0.1) and +(-76:0) ..
				(46.9, 11.08) -- (46.83, 11.31) .. controls +(104:0) and +(180:0.1) ..
				(46.99, 11.38) -- (47.56, 11.38) .. controls +(0:0.1) and +(104:0) ..
				(47.77, 11.31) -- (47.86, 11.08) .. controls +(-76:0) and +(0:0.1) ..
				(47.7, 11.01) -- cycle
 (35.76, 10.48) -- (33.95, 10.48) .. controls +(180:0.1) and +(-124:0) ..
				(33.88, 10.53) -- (34.09, 10.84) .. controls +(45:0) and +(180:0.1) ..
				(34.22, 10.88) -- (35.97, 10.88) .. controls +(0:0.1) and +(45:0) ..
				(36.05, 10.84) -- (35.89, 10.53) .. controls +(-108:0) and +(0:0.1) ..
				(35.76, 10.48) -- cycle
 (36.95, 10.48) -- (36.15, 10.48) .. controls +(180:0.1) and +(-108:0) ..
				(36.07, 10.53) -- (36.23, 10.84) .. controls +(63:0) and +(180:0.1) ..
				(36.35, 10.88) -- (37.12, 10.88) .. controls +(0:0.1) and +(63:0) ..
				(37.21, 10.84) -- (37.07, 10.53) .. controls +(-108:0) and +(0:0) ..
				(36.95, 10.48) -- cycle
 (38.13, 10.48) -- (37.34, 10.48) .. controls +(180:0.1) and +(-108:0) ..
				(37.25, 10.53) -- (37.38, 10.84) .. controls +(63:0) and +(180:0) ..
				(37.5, 10.88) -- (38.27, 10.88) .. controls +(0:0.1) and +(63:0) ..
				(38.36, 10.84) -- (38.25, 10.53) .. controls +(-108:0) and +(0:0.1) ..
				(38.13, 10.48) -- cycle
 (39.32, 10.48) -- (38.52, 10.48) .. controls +(180:0.1) and +(-108:0) ..
				(38.43, 10.53) -- (38.54, 10.84) .. controls +(63:0) and +(180:0) ..
				(38.65, 10.88) -- (39.42, 10.88) .. controls +(0:0.1) and +(90:0) ..
				(39.52, 10.84) -- (39.43, 10.53) .. controls +(-90:0) and +(0:0) ..
				(39.32, 10.48) -- cycle
 (45.24, 10.48) -- (39.7, 10.48) .. controls +(180:0.1) and +(-108:0) ..
				(39.61, 10.53) -- (39.68, 10.84) .. controls +(63:0) and +(180:0.1) ..
				(39.8, 10.88) -- (45.18, 10.88) .. controls +(0:0) and +(117:0) ..
				(45.29, 10.84) -- (45.34, 10.53) .. controls +(-72:0) and +(0:0.1) ..
				(45.24, 10.48) -- cycle
 (46.43, 10.48) -- (45.63, 10.48) .. controls +(180:0.1) and +(-72:0) ..
				(45.52, 10.53) -- (45.46, 10.84) .. controls +(90:0) and +(180:0) ..
				(45.55, 10.88) -- (46.33, 10.88) .. controls +(0:0.1) and +(117:0) ..
				(46.44, 10.84) -- (46.52, 10.53) .. controls +(-72:0) and +(0:0.1) ..
				(46.43, 10.48) -- cycle
 (47.62, 10.48) -- (46.83, 10.48) .. controls +(180:0.1) and +(-72:0) ..
				(46.71, 10.53) -- (46.62, 10.84) .. controls +(117:0) and +(180:0.1) ..
				(46.71, 10.88) -- (47.48, 10.88) .. controls +(0:0.1) and +(117:0) ..
				(47.6, 10.84) -- (47.71, 10.53) .. controls +(-72:0) and +(0:0.1) ..
				(47.62, 10.48) -- cycle
 (48.84, 11.01) -- (48.24, 11.01) .. controls +(180:0.1) and +(-63:0) ..
				(48.03, 11.08) -- (47.94, 11.31) .. controls +(104:0) and +(180:0.1) ..
				(48.09, 11.38) -- (48.67, 11.38) .. controls +(0:0.1) and +(104:0) ..
				(48.88, 11.31) -- (48.99, 11.08) .. controls +(-63:0) and +(0:0.1) ..
				(48.84, 11.01) -- cycle
 (51.66, 11.01) -- (49.41, 11.01) .. controls +(180:0.1) and +(-63:0) ..
				(49.19, 11.08) -- (49.08, 11.31) .. controls +(117:0) and +(180:0.1) ..
				(49.23, 11.38) -- (51.41, 11.38) .. controls +(0:0.1) and +(127:0.1) ..
				(51.64, 11.31) -- (51.79, 11.08) .. controls +(-53:0.1) and +(0:0.1) ..
				(51.66, 11.01) -- cycle
 (50.54, 10.72) -- (49.5, 10.72) .. controls +(180:0) and +(-63:0) ..
				(49.38, 10.76) -- (49.34, 10.84) .. controls +(117:0) and +(180:0.1) ..
				(49.42, 10.88) -- (50.44, 10.88) .. controls +(0:0.1) and +(117:0) ..
				(50.57, 10.84) -- (50.61, 10.76) .. controls +(-45:0) and +(0:0.1) ..
				(50.54, 10.72) -- cycle
 (50.67, 10.48) -- (49.62, 10.48) .. controls +(180:0.1) and +(-72:0) ..
				(49.49, 10.53) -- (49.45, 10.61) .. controls +(117:0) and +(180:0.1) ..
				(49.53, 10.65) -- (50.57, 10.65) .. controls +(0:0.1) and +(117:0) ..
				(50.7, 10.61) -- (50.75, 10.53) .. controls +(-72:0) and +(0:0.1) ..
				(50.67, 10.48) -- cycle
 (52.1, 10.48) -- (51.05, 10.48) .. controls +(180:0.1) and +(-56:0) ..
				(50.91, 10.53) -- (50.87, 10.61) .. controls +(135:0) and +(180:0) ..
				(50.94, 10.65) -- (51.98, 10.65) .. controls +(0:0.1) and +(135:0) ..
				(52.12, 10.61) -- (52.17, 10.53) .. controls +(-56:0) and +(0:0.1) ..
				(52.1, 10.48) -- cycle
 (49.24, 10.48) -- (48.19, 10.48) .. controls +(180:0.1) and +(-72:0) ..
				(48.07, 10.53) -- (48.04, 10.61) .. controls +(117:0) and +(180:0) ..
				(48.12, 10.65) -- (49.16, 10.65) .. controls +(0:0.1) and +(117:0) ..
				(49.29, 10.61) -- (49.33, 10.53) .. controls +(-72:0) and +(0:0.1) ..
				(49.24, 10.48) -- cycle;
		\end{scope}

		\path[fill=cededed] (45.46, 8.66) -- (40.62, 8.66) .. controls +(180:0.2) and +(-100:0.1) ..
			(40.22, 8.87) -- (40.43, 10.08) .. controls +(78:0.1) and +(180:0.2) ..
			(40.86, 10.25) -- (45.2, 10.25) .. controls +(0:0.2) and +(96:0.1) ..
			(45.62, 10.08) -- (45.86, 8.87) .. controls +(-80:0.1) and +(0:0.2) ..
			(45.46, 8.66) -- cycle;
		\draw[ultra thick,red] (C) node[shift={(145:12pt)},font=\Huge]{$C$} --($(B)+(0,-16pt)$) coordinate (B) node[shift={(170:11pt)},text=red!90!black,font=\Huge]{$A$}  --(A) node[left,font=\Huge]{$B$};	
	\foreach \x/\y/\z in {A/B/C}{
			\path (\y) pic[draw=red,angle radius=22pt,angle eccentricity=1.25,text=red,"$\alpha$",font=\Large]{angle = \x--\y--\z};
	}
	
	\path (current bounding box.south) node[below=7pt,font=\fontsize{40pt}{0pt}\selectfont]{Hình 7};
	\end{tikzpicture}
}
\usepackage{tikz-3dplot}
\usetikzlibrary{positioning,calc}
%==================== ĐỀ CHÍNH THỨC: in đen trắng, đầu đề Sở/Trường ====================
\definecolor{mauchuong}{HTML}{000000}
\definecolor{mausubsection}{HTML}{000000}
\definecolor{mausection}{HTML}{000000}
\usepackage{fancyhdr}
\pagestyle{fancy}\fancyhf{}
\renewcommand{\headrulewidth}{0pt}
\def\madehientai{}
\newcommand{\dauDe}[1]{%
  \setcounter{page}{1}\gdef\madehientai{#1}%
  \fancyfoot[C]{\small Trang \thepage/\pageref{cuoide\madehientai} -- Mã đề \madehientai}%
  \noindent\begin{minipage}[t]{0.47\textwidth}\centering
    {\small SỞ GD\&ĐT TP. HỒ CHÍ MINH}\\[1pt]
    {\small\bfseries TRƯỜNG THPT NGUYỄN HỮU CẢNH}\\[-6pt]
    \rule{3.2cm}{0.5pt}\\[1pt]
    {\small\itshape (Đề thi có \pageref{cuoide#1} trang)}
  \end{minipage}\hfill
  \begin{minipage}[t]{0.52\textwidth}\centering
    {\bfseries ĐỀ KIỂM TRA GIỮA HỌC KỲ I}\\[1pt]
    {\bfseries NĂM HỌC 2026 -- 2027}\\[1pt]
    {\bfseries Môn: TOÁN -- Lớp 12}\\[1pt]
    {\small\itshape Thời gian làm bài: 90 phút, không kể thời gian phát đề}
  \end{minipage}\par\medskip
  \noindent Họ, tên thí sinh: \makebox[7.2cm]{\dotfill}\hfill
  \fbox{\bfseries\,Mã đề #1\,}\par\noindent
  Số báo danh: \makebox[5.5cm]{\dotfill}\par\medskip}
\newcommand{\phanDe}[1]{\par\medskip\noindent #1\par\smallskip}
\newcommand{\hetDe}[1]{\par\bigskip\begin{center}{\bfseries ------ HẾT ------}\end{center}\label{cuoide#1}}
\newcommand{\phanDapAn}[1]{\newpage\fancyfoot[C]{\small Đáp án -- Mã đề #1}%
  \begin{center}{\large\bfseries ĐÁP ÁN -- MÃ ĐỀ #1}\end{center}}
\begin{document}
\setcounter{tocdepth}{1}
\setcounter{secnumdepth}{3}
\hienthiloigiai
\dauDe{101}

\phanDe{\textbf{PHẦN I. Câu trắc nghiệm nhiều phương án lựa chọn.} Thí sinh trả lời từ câu 1 đến câu 12. Mỗi câu hỏi thí sinh chỉ chọn một phương án.}

\setcounter{ex}{0}

\Opensolutionfile{ans}[Ans/GHK1_12_101_TN]

\begin{ex}%[0D1B1-3]%[ma:5c0c30c1]
	Mệnh đề nào sau đây là phủ định của  mệnh đề “Mọi động vật đều di chuyển” ?
	 \choice
		{ \True Có  ít nhất một động vật không di chuyển}
		{ Mọi động vật đều không di  chuyển}
		{ Có ít nhất một động vật di  chuyển}
		{ Mọi động vật đều đứng  yên}
	 \loigiai{
		Theo định nghĩa, phủ định của mệnh đề “Mọi động vật đều di chuyển” là mệnh đề “Có ít nhất một động vật không di  chuyển”.
	 }
\end{ex}

\begin{ex}%[0D1B3-3]%[ma:5f619861]
	Để thành lập đội tuyển học sinh giỏi khối $ 10 $, nhà trường tổ chức thi chọn các môn Toán, Văn, Anh trên tổng số $ 111 $ học sinh. Kết quả  có: $ 70 $ học sinh giỏi Toán, $ 65 $ học sinh giỏi Văn, $ 62 $ học sinh giỏi Anh. Trong đó có $ 49 $ học sinh giỏi cả hai môn Văn và Toán, $ 32 $ học sinh giỏi cả hai môn Toán và Anh, $ 34 $ học sinh giỏi cả hai môn Văn và Anh. Xác định số học sinh giỏi cả ba môn Văn, Toán, Anh. Biết rằng có $ 6 $ học sinh không đạt yêu cầu cả ba môn.
	\choice 
		 {$  32 $}
		{$ 34  $}
		{$ 36 $}
		 {\True $ 23 $}
	 \loigiai{
	 \immini
		{ 
		Có $ 111-6=105 $ học sinh thi đạt ít nhất $ 1 $ môn.\\ 
		Gọi $ A $ là số học sinh giỏi môn  Toán và Tiếng Anh nhưng không giỏi Văn.\\
		Gọi $ B $ là số học sinh giỏi môn Toán và Văn nhưng không giỏi Tiếng Anh.\\ 
		Gọi $ C  $ là số học sinh giỏi môn Văn và Tiếng Anh nhưng không giỏi Toán.\\
		Gọi $ D $ là  số học sinh giỏi cả ba môn. Ta có hệ:\\
		$\begin{cases}
		 B+D=49\\
		A+D=32\\ 
		 C+D=34\\
		70+65+62-(A+B+C+2D)=105 
	\end{cases}\\
	\Rightarrow  92=32-D+49-D+34-D+2D\\
	\Rightarrow D=23$.\\ 
	Vậy có $  23 $ học sinh giỏi cả ba môn.
	 }
	{ 
	\begin{tikzpicture}[scale=.8]
			\def\radius{2cm} 
			\coordinate  (ceni);
			\coordinate[xshift=.8*\radius,yshift=.2*\radius] (cenii); 
			\coordinate[xshift=.6*\radius,yshift=-.8*\radius] (ceniii); 
			\draw (ceni)  circle (.9*\radius);
			\draw (cenii)  circle (\radius);
			\draw (ceniii)  circle (\radius);
			 \node[xshift=.3*\radius,yshift=.25*\radius] at (ceni) {A};
			\node[xshift=.1*\radius,yshift=-.5*\radius] at (ceni) {B}; 
			\node[xshift=.45*\radius,yshift=-.25*\radius] at (ceni) {D}; 
			\node[xshift=.1*\radius,yshift=-.55*\radius] at (cenii)  {C};
		\end{tikzpicture}
	} 
} 
\end{ex}

\begin{ex}%[0D1N2-3]%[ma:b26dadc3]
	Trong các tập hợp sau, tập nào là  tập rỗng?
	 \choice
		 {\True $C=\{ x\in \mathbb{Z} |~ x^2+x-1=0  \}$}
		{$D=\{ x \in \mathbb{R} |~ x^2+5x-1=0 \}$}
		{$B=\{  x \in \mathbb{Q} |~ 3x^2-5x+2=0 \}$}
		{$A= \{ x \in \mathbb{R} |~ x^2+5x-6=0 \}$} 
	 \loigiai{
		Ta có\\ 
		 $x^2+5x-6=0
		\Leftrightarrow \hoac{&x=1 \in \mathbb{R}\\&x=-6 \in \mathbb{R}}$. Vậy $A=\{ -6;1  \} \neq \varnothing$.\\
		$3x^2-5x+2=0 
		\Leftrightarrow \hoac{&x=1 \in \mathbb{Q}\\&x=\dfrac{2}{3} \in \mathbb{Q}}$. Vậy $B=\left\{ 1;\dfrac{2}{3}  \right\} \neq \varnothing$.\\
		 $x^2+x-1=0
		\Leftrightarrow \hoac{&x=\dfrac{-1+\sqrt{5}}{2} \notin \mathbb{Z}\\&x=\dfrac{-1-\sqrt{5}}{2} \notin \mathbb{Z}}$. Vì $x \in \mathbb{Z}$  nên $C=\varnothing$.\\
		$x^2+5x-1=0 
		\Leftrightarrow \hoac{&x=\dfrac{-5+\sqrt{29}}{2} \in \mathbb{R}\\&x=\dfrac{-5-\sqrt{29}}{2} \in \mathbb{R}}$.  Vậy $D=\left\{ \dfrac{-5+\sqrt{29}}{2};\dfrac{-5-\sqrt{29}}{2} \right\} \neq \varnothing$.
	} 
\end{ex}

\begin{ex}%[0D1B1-5]%[ma:ca8f23aa]
	Cho mệnh đề ``$\forall x \in  \mathbb{R} : x^2+x+1>0$''. Mệnh đề phủ định của mệnh đề đã cho là
	\choice{``$\exists x  \in \mathbb{R} :x^2+x+1 < 0$''} 
		 {``$\forall x \in \mathbb{R} :x^2+x+1 = 0$''}
		{``$\forall x \in \mathbb{R} :x^2+x+1 < 0$''}
		{\True ``$\exists x \in \mathbb{R} :x^2+x+1  \le 0$''}
	\loigiai{ Mệnh đề phủ định của mệnh đề ``$\forall x \in \mathbb{R} : x^2+x+1>0$'' là mệnh đề  ``$\exists x \in \mathbb{R} :x^2+x+1 \le 0$''.
	 }
\end{ex}

\begin{ex}%[0D1B2-1]%[ma:f3f80b92]
	Cho tập hợp $A = \left\{n \in \mathbb{Z}\big| \,6 \, \vdots \, n\right\}$. Hãy  viết tập hợp $A$ dưới dạng liệt kê?
	\choice{$A = \left\{1;2;3;6\right\}$}
		 {\True $A =  \left\{-6;-3;-2;-1;1;2;3;6\right\}$}
		{$A = \left\{-6;-3;-2;-1\right\}$} 
		{$A  = \left\{-6;-1;1;6\right\}$}
	 \loigiai{
		$A$ là tập các ước nguyên  của $6$, do đó $A = \left\{-6;-3;-2;-1;1;2;3;6\right\}$.
	} 
\end{ex}

\begin{ex}%[0D1B3-1]%[ma:a97e4b36]
	Cho hai tập hợp $A=\left\{1;2;3;4\right\}$, $B=\left\{ 2;4;6;8 \right\}$. Tập hợp nào sau đây  bằng tập hợp $A\cap B$?
	 \choice
		{\True$\left\{ 2;4  \right\}$}
		 {$\left\{6;8  \right\}$}
		{$\left\{1;2;3;4;6;8 \right\}$}
		 {$\left\{1;3 \right\}$}
	\loigiai{ 
		Ta có $A\cap B=\left\{ 2;4 \right\}.$ 
	 }
\end{ex}

\begin{ex}%[0D1H3-5]%[ma:1d662c9d]
	Lớp 10A có $45$ học sinh trong đó có $25$ em học giỏi môn Toán, $23$ em học giỏi  môn Lý, $20$ em học giỏi môn Hóa, $11$ em học giỏi cả môn Toán và môn Lý, $8$ em học giỏi cả môn Lý và môn Hóa, $9$ em học giỏi cả môn Toán và môn Học. Hỏi lớp $10$ A có bao nhiêu bạn học giỏi cả ba môn Toán, Lý, Hóa, biết rằng mỗi học sinh trong lớp học giỏi ít nhất một trong $3$ môn Toán, Lý, Hóa?
	\choice 
		{$6$} 
		 {$3$}
		 {$4$}
		{\True $5$} 
	 \loigiai{
		\begin{center}
			\begin{tikzpicture}[scale=.8,font=\footnotesize,line join=round,line cap=round,>=stealth]
					\coordinate  (A) at (0,0);
					 \coordinate (B) at (2,0);
					\coordinate  (C) at (1,-1.5);
					\begin{scope}
						 \clip (C) circle (1.5);
						 \clip (A) circle (1.5);
						\fill[pattern=north west lines]  (B) circle (1.5);
					\end{scope}
					\draw (A)  circle (1.5);
					\draw (B)  circle (1.5);
					\draw (C) circle  (1.5);
					 \path
						 (A)++(90:.6) node[fill=white,above]{$L$}
						(B)++(90:.6)  node[fill=white,above]{$T$}
						(C)++(-90:1)  node[fill=white,above]{$H$};
				\end{tikzpicture}
		\end{center}
		Gọi $T$, $L$, $H$ lần lượt là tập hợp các học sinh giỏi môn Toán,  Lý, Hóa.\\
		Khi đó 
		\begin{eqnarray*}
			|T \cup L \cup H|&=&|T|+|L|+|H|-|T \cap L|-|L \cap H|-|H \cap T|+|T \cap L  \cap H| \\
			 45&=&25+23+20-11-8-9+|T \cap L \cap H|\\
			|T  \cap L \cap H|&=&5
		\end{eqnarray*}
		Vậy có $5$ học  sinh giỏi cả $3$ môn.
	 }
\end{ex}

\begin{ex}%[0D1H2-3]%[ma:e60745e5]
	Tìm tất cả các giá trị của  tham số $m$ sao cho $(m;+\infty)\subset(1;+\infty)$?
	 \choice
		{$m<1$} 
		 {$m>1$}
		{$m\leq 1$} 
		 {\True $m\geq 1$}
	 \loigiai{
		Do $(m;+\infty)\subset(1;+\infty)$  nên $m\geq 1$.
	} 
\end{ex}

\begin{ex}%[0D1H3-3]%[ma:76b2c2a1]
	Tổng các  số nguyên của tập $[-3;2]\setminus(1;2)$ bằng
	 \choice
		 {$4$}
		 {$-5$}
		 {$6$}
		{\True$-3$} 
	 \loigiai{
		Ta có $[-3;2]\setminus(1;2)=[3;1]\cup\{2\}$ nên chứa các số nguyên  là $-3;-2;-1;0;1;2$.\\
		Vậy tổng các số nguyên là $-3$. 
	 }
\end{ex}

\begin{ex}%[0D1H1-5]%[ma:9a292bba]
	Trong các mệnh đề dưới đây, mệnh đề nào  là mệnh đề đúng?
	 \choice
		{\lq\lq $\exists x\in \mathbb{R}\mid x^2+1=0$\rq\rq} 
		{\lq\lq $\exists x\in \mathbb{N}\mid x^2-2=0$\rq\rq}
		{\lq\lq $\exists x\in  \mathbb{R}\mid x^2< 0$\rq\rq}
		{\True \lq\lq $\exists x\in \mathbb{N}\mid 2x^2-1<  0$\rq\rq} 
	 \loigiai{
		Ta  có $x^2\ge 0\Leftrightarrow x^2+1\ge 1$, $\forall x\in \mathbb{R}$.
		Nên  $\exists x\in \mathbb{R}\mid x^2+1=0$ là mệnh đề sai.\\
		Ta có $x^2\ge 0$, $\forall x\in \mathbb{R}$. Nên $\exists x\in \mathbb{R}\mid x^2< 0$ là  mệnh đề sai.\\
		$2x^2-1< 0\Leftrightarrow x^2< \dfrac{1}{2} \Leftrightarrow-\dfrac{\sqrt{2}}{2} < x < \dfrac{\sqrt{2}}{2}$, mà $x\in \mathbb{N}\Rightarrow x=0$. Vậy $\exists x\in \mathbb{N}\mid 2x^2-1< 0$ là mệnh  đề đúng.\\
		$x^2-2=0\Leftrightarrow x=\pm \sqrt{2} \notin \mathbb{N}$. Vậy $\exists x\in  \mathbb{N}\mid x^2-2=0$ là mệnh đề sai.
	 }
\end{ex}

\begin{ex}%[0D1K1-3]%[ma:a990ffed]
	Phát biểu mệnh đề  phủ định của mệnh đề "Tồn tại một loài thú có nhiệt độ thân thể nhỏ hơn $35^\circ$C".
	\choice 
		{"Mọi loài thú có nhiệt độ thân thể lớn hơn $35^\circ$C"}
		{"Tồn tại một loài thú có nhiệt  độ thân thể lớn hơn $35^\circ$C"}
		{"Mọi loài thú có nhiệt độ thân thể  nhỏ hơn $35^\circ$C"}
		{\True "Mọi loài thú có nhiệt độ thân thể  không nhỏ hơn $35^\circ$C"} 
	\loigiai{Mệnh đề phủ định của mệnh đề "Tồn tại một loài thú có nhiệt độ thân thể nhỏ hơn $35^\circ$C" là  mệnh đề "Mọi loài thú có nhiệt độ thân thể không nhỏ hơn $35^\circ$C".}
\end{ex}

\begin{ex}%[0D1K1-5]%[ma:2b89aaa2]
	Trong các mệnh đề sau,  mệnh đề nào \textbf{sai}?
	\choice 
		{$\exists n\in \mathbb{N}$, $n^2+11n+2$ chia  hết cho $11$}
		{\True $\exists n\in \mathbb{N}$, $n^2+1$ chia hết cho $4$} 
		{Tồn tại số nguyên  tố chia hết cho $5$}
		 {$\exists n\in \mathbb{Z}$, $2x^2-8=0$}
	 \loigiai{
		\begin{itemize}
			\item Xét  mệnh đề ``$\exists n\in \mathbb{N}$, $n^2+11n+2$ chia hết cho $11$''.\\
			Khi $n=3$ thì  giá trị của $(n^2+11n+2)$ bằng $44\vdots 11$ nên mệnh đề đúng.
			\item Xét mệnh đề ``$\exists n\in \mathbb{N}$, $n^2+1$  chia hết cho $4$''.\\
			Khi  $n=2k,k\in \mathbb{N} \Rightarrow n^2+1=4k^2+1$ không chia hết cho $4$, $k\in \mathbb{N}$.\\
			Khi $n=2k+1,k\in \mathbb{N} \Rightarrow n^2+1=(2k+1)^2+1=4k^2+4k+2$ không  chia hết cho $4$, $k\in \mathbb{N}$.
			\item Xét mệnh đề ``Tồn tại  số nguyên tố chia hết cho $5$''.\\
			Tồn tại số nguyên tố $5$ chia hết cho $5$  nên mệnh đề đúng.
			\item  Xét mệnh đề ``$\exists n\in \mathbb{Z}$, $2x^2-8=0$ ''.\\
			Phương trình $2x^2-8=0 \Leftrightarrow x^2=4 \Leftrightarrow x=-2;x=2\in  \Bbb{Z}$ nên mệnh đề đúng.
		\end{itemize}
	 }
\end{ex}

\Closesolutionfile{ans}

\phanDe{\textbf{PHẦN II. Câu trắc nghiệm đúng sai.} Thí sinh trả lời từ câu 1 đến câu 4. Trong mỗi ý a), b), c), d) ở mỗi câu, thí sinh chọn đúng hoặc sai.}

\setcounter{ex}{0}

\Opensolutionfile{ansbook}[Ansbook/GHK1_12_101_DS]

\begin{ex}%[0D1H1-3]%[ma:f76fc24c]
	Cho  mệnh đề $P(x):$ \lq\lq $\forall x\in \mathbb{R},x^2-4x+1> 0$\rq\rq.
	 \choiceTF
		{$P(x)$: \lq\lq Có ít nhất một số thực $x$ thoả mãn $x^2-4x+1> 0$\rq\rq} 
		{$P(x)$ là mệnh đề chứa biến} 
		{\True  $\overline{P(x)}$: \lq\lq $\exists x\in \mathbb{R}, x^2-4x+1\le 0$\rq\rq}
		{Mệnh đề $\overline{P(x)}$ là mệnh  đề sai}
	 \loigiai{
		\begin{itemchoice}
			\itemch  SAI.\\
			Ta có phát biểu đúng của mệnh đề $P(x):$\lq\lq Với mọi số thực  $x$ thoả mãn $x^2-4x+1> 0$\rq\rq.
			\itemch SAI.\\ 
			$P(x)$ là mệnh đề sai nên không phải là  mệnh đề chứa biến.
			\itemch  ĐÚNG.\\
			Mệnh đề phủ định  của $P(x)$ là $\overline{P(x)}:$ \lq\lq $\exists x\in \mathbb{R}, x^2-4x+1\le 0$\rq\rq.
			\itemch SAI.\\ 
			Ta có $\overline{P(x)}:$ \lq\lq $\exists x\in \mathbb{R}, x^2-4x+1\le 0$\rq\rq. Mệnh đề $\overline{P(x)}$ đúng vì giả sử ta  chọn $x=1$ thì $1^2-4.1+1=-2< 0$.
		\end{itemchoice}
	} 
\end{ex}

\begin{ex}%[0D1H1-2]%[ma:bff4f470]
	Cho $\Delta ABC$ có $A'$, $B'$, $C'$ lần lượt là các trung điểm của các  cạnh $BC$, $CA$, $AB$. Khi đó
	 \choiceTF
		{\True $B{C'}=C' A=A'{B'}=\dfrac{AB}{2}{\rm}$} 
		{Hai vectơ $\overrightarrow{BC'}$,  $\overrightarrow{A'B'}$ ngược hướng}
		{\True $\overrightarrow{B{C'}}=\overrightarrow{C' A} =\overrightarrow{A'B'}$} 
		{\True $\overrightarrow{B'  C'}=\overrightarrow{CA'}$}
	\loigiai{ 
		\begin{center}
			\begin{tikzpicture}[scale=1,line join=round,line cap=round,font=\footnotesize,>=stealth]
					\def\r{3} 
					\path (180:\r)  coordinate (B)
						(0:\r)  coordinate (C)
						(120:\r) coordinate  (A)
						($(B)!0.5!(C)$) coordinate  (A')
						($(C)!0.5!(A)$) coordinate (B') 
						($(B)!0.5!(A)$)  coordinate (C');
					\draw  (A)--(B)--(C)--cycle;
					\foreach \x/ \goc in  {A/135,B/-135,C/-45,A'/-90,B'/30,C'/180}
					 \fill (\x) circle (1pt)
						 ($(\x)+(\goc:3mm)$) node {$\x$};
				\end{tikzpicture}
		\end{center}
		\begin{itemchoice}
			\itemch Ta có $C'$ là trung điểm của $AB$ và $A'B'$ là đường trung bình của tam giác $ABC$ ứng với  cạnh đáy $AB$ nên $BC'=C'A=A'B'=\dfrac{AB}{2}$.
			\itemch  Vì Hai vectơ $\overrightarrow{BC'}$, $\overrightarrow{A'{B'}}$ cùng hướng.
			\itemch Vì $C'$ là trung điểm của $AB$ và $A'B'$ là đường trung bình của tam giác $ABC$ ứng với cạnh đáy $AB$ nên $\overrightarrow{BC'}=\overrightarrow{C'A}  =\overrightarrow{A'B'}$.
			\itemch Vì $A'$ là trung điểm của $BC$ và $B'C'$ là đường trung bình  của tam giác $ABC$ ứng với cạnh đáy $BC$ nên $\overrightarrow{B'C'}=\overrightarrow{CA'}$.
		\end{itemchoice}
	 }
\end{ex}

\begin{ex}%[0D1V2-1]%[ma:161eac9d]
	 Cho các tập hợp sau: A các số nguyên tố nhỏ hơn 11; $B=\left\{\left. x\in \mathbb{R}\right|3x^2-4x+1=0\right\}$;
	$C=\left\{\left. x\in  \mathbb{N}\right|\left(x^2-5x+6\right)\left(2x+1\right)=0\right\}$; $D=\{\left. x\in \mathbb{Z}\right|\left|x+1\right|<3\}$.
	 \choiceTF
		{\True Tập hợp $A$ có $4$ phần tử} 
		{Tập hợp $B$ có  $3$ phần tử}
		{Tập hợp  $C$ có $3$ phần tử}
		{\True Tập  hợp $D$ có $3$ phần tử}
	 \loigiai{
		\begin{itemchoice}
			\itemch Ta  có các số nguyên tố nhỏ hơn $11$ là $2$; $3$; $5$; $7$. \\
			Vậy $A=\{2; 3; 5; 7\}$. 
			\itemch Ta có $3x^2-4x+1=0\Leftrightarrow \hoac{ 
			&x=1\in \mathbb{R} \\ 
			 &x=\dfrac{1}{3}\in \mathbb{R}.}$ \\
			Vậy $B=\left\{\dfrac{1}{3};  1\right\}$.
			 \itemch Ta có $\left(x^2-5x+6\right)\left(2x+1\right)=0\Leftrightarrow \hoac{
			& x^2-5x+6=0 \\ 
			& 2x+1=0}\Leftrightarrow  \hoac{
			 & x=2\in \mathbb{N} \\
			& x=3\in \mathbb{N}  \\
			&  x=\dfrac{-1}{2}\notin \mathbb{N}.} $\\
			Vậy  $C=\left\{2; 3\right\}$.
			 \itemch Ta có $\heva{
			&x\in \mathbb{Z}  \\
			 &|x+1|<2}\Rightarrow \hoac{
			&x=-2  \\
			 &x=-1 \\
			 &x=0.} $\\
			Vậy $D=\{-2;-1; 0\}$. 
		\end{itemchoice}
	} 
\end{ex}

\begin{ex}%[0D1V1-2]%[ma:d97640b0]
	Cho hình thang $ABCD$ vuông tại $A$ và có $AB=AD=\dfrac{1}{2}DC=a$. Gọi $BF$ là  đường phân giác trong của tam giác $ABD$ ($F\in AD$). Khi đó
	 \choiceTF
		{\True  $CA^2=DA^2+DC^2$}
		{$\left|\overrightarrow{CA}\right|=a\sqrt  3 $}
		 {$\widehat{ABF}=45^\circ$}
		{$\left|\overrightarrow{BF}\right|\approx 2{,}08a$} 
	\loigiai{\begin{center}
		\begin{tikzpicture}[>=stealth,line join=round,line cap=round,font=\footnotesize,scale=0.8]
				\def\a{2.4} 
				\path 
					 (0,4) coordinate (A)
					(0,0) coordinate  (D)
					(8,0) coordinate (C) 
					(4,4)  coordinate (B)
					($(A)!1/\a!(D)$) coordinate  (F);
				 \draw (A)--(B)--(C)--(D)--(A)
					 (A)--(C)
					(F)--(B)--(D); 
				\foreach \x/\g in  {B/60,A/120,C/30,D/-60,F/-140} \fill (\x) circle (1pt) +(\g:3mm) node[scale=.8] {$\x$};
				 \draw pic[thin,draw,angle radius=2mm] {right angle= B--A--D};
				 \draw pic[thin,draw,angle radius=2mm] {right angle= A--D--C};
			\end{tikzpicture}
	\end{center}
	\begin{itemchoice}
		\itemch Áp dụng định lí Pytago cho tam giác $ADC$ vuông  tại $D$.
		\itemch Áp dụng  định lí Pytago cho tam giác $ADC$ vuông tại $D$, ta có\\
		$C{A^2}=D{A^2}+D{C^2}=a^2+(2a)^2=5a^2$\\ 
		 $\Rightarrow \left|\overrightarrow{CA}\right|=CA=a\sqrt 5$.
		\itemch Ta có  $\Delta ABD$ vuông cân tại $A$, do đó\\
		$\widehat{ABD}=45^\circ  \Rightarrow \widehat{ABF}=22{,}5^\circ$.
		\itemch Xét  $\Delta ABF$ vuông tại $A$, ta có\\
		$ \left| \overrightarrow{BF} \right| =BF=\dfrac{AB}{\cos \left(  ABF \right)}=\dfrac{a}{\cos{22{,}5^\circ}}\approx 1{,}08a$.
	\end{itemchoice}
 }
\end{ex}

\Closesolutionfile{ansbook}

\phanDe{\textbf{PHẦN III. Câu trắc nghiệm trả lời ngắn.} Thí sinh trả lời từ câu 1 đến câu 5.}

\setcounter{bt}{0}

\Opensolutionfile{ansbt}[Ansbook/GHK1_12_101_TLN]

\begin{bt}%[0D1H1-3]%[ma:d05bebba]
	Cho mệnh đề $P$: \lq\lq $x^2+x+1$ là số dương với $x$ là số thực\rq\rq\ và các  mệnh đề sau.
	\begin{itemize}
		\item \lq\lq $x^2+x+1$  là số không dương với $x$ là số thực\rq\rq.
		\item \lq\lq $x^2+x+1$ là khác $0$ với $x$ là số thực\rq\rq. 
		\item \lq\lq $x^2+x+1$ là số âm với $x$ là số thực\rq\rq. 
		\item \lq\lq $x^2+x+1$  là số không âm với $x$ là số thực\rq\rq.
	\end{itemize}
	Có bao nhiêu phát biểu là phủ định của mệnh  đề $P$?
	\shortans{1} 
	 \loigiai{
		Phủ định của mệnh đề $P$ là $\overline{P}$: \lq\lq $x^2+x+1$ là số không dương  với $x$ là số thực\rq\rq.\\
		Vậy  có $1$ mệnh đề thỏa. Các mệnh đề còn lại không phải là phủ định của mệnh đề $P$.
	 }
\end{bt}

\begin{bt}%[0D1H1-6]%[ma:deff5465]
	Cho phát biểu $\forall x\in\mathbb{R}$, $x^{2}+a-5>0$ với $a$ là số thực cho trước. Có bao giá trị  nguyên của $a<20$ để phát biểu đã cho là mệnh đề đúng?
	\shortans{14} 
	 \loigiai{
		Mệnh đề $\forall x\in\mathbb{R}$, $x^{2}+a-5>0$ đúng $\Leftrightarrow  x^{2}>5-a$ đúng $\forall x\in\mathbb{R}$.\\
		Điều này xảy ra khi và chỉ khi  $\max(5-a, 0) < 0$.\\
		Vì $x^2 \ge 0$, nên ta cần  $5-a$ phải nhỏ hơn $0$.\\
		 $5-a<0\Leftrightarrow a>5$.\\
		Kết hợp  với điều kiện $a<20$, ta có $5<a<20$.\\
		Các giá trị nguyên của $a$ là $6, 7, 8, \dots,  19$.\\
		Số giá trị nguyên của  $a$ là $19-6+1=14$.
	} 
\end{bt}

\begin{bt}%[0D1V2-2]%[ma:8c3a2e78]
	Lớp $10$A có $10$ học sinh giỏi Toán, $10$ học sinh giỏi Lý, $11$ học sinh giỏi hóa, $6$ học sinh giỏi cả Toán và Lý, $5$ học sinh giỏi cả Hóa và Lý, $4$ học sinh giỏi cả Toán và Hóa, $3$ học sinh giỏi cả ba môn Toán, Lý,  Hóa. Tính học sinh giỏi ít nhất một trong ba môn (Toán, Lý, Hóa) của lớp $10$A?
	 \shortans{18}
	\loigiai{ 
		Theo giả thiết đề bài cho, ta có  biểu đồ Ven
		\begin{center}
			\begin{tikzpicture}[>=stealth,line join=round,line cap=round,font=\footnotesize,scale=1]
					\draw [thick,green!60!black] (0,0) circle (1.5) 
						(1.5,-0.5)  circle (1)
						(1.5,0.5) circle  (0.8);
					\draw (1,.8) node {$6$} 
						 (1.2,.2) node {$3$}
						(2,.8) node {Lý} 
						 (1.8,.1) node {$5$}
						(1.8,-1)  node {Hoá}
						(.8,-.8)  node {$4$}
						(-1,0) node  {Toán};
				\end{tikzpicture}
		\end{center}
		Dựa vào biểu đồ Ven, ta có học sinh giỏi ít nhất một trong ba môn (Toán, Lý, Hóa)  của lớp $10$A là
		\begin{itemize}
			\item Số học  sinh giỏi Toán $6+4+3=13$.
			\item Số học sinh  giỏi Lý $6+5+3=14$.
			\item Số học  sinh giỏi Hóa $4+5+3=12$.
		\end{itemize}
		Ta lại có 
		\begin{itemize}
			\item Số học  sinh giỏi cả Toán và Lý $6$.
			\item Số học sinh giỏi cả Toán và Hóa $4$. 
			\item Số học sinh giỏi cả Hóa  và Lý $5$.
			\item Và số học  sinh giỏi cả Toán, Lý và Hóa là $3$.
			\item Số học sinh giỏi hơn một môn  là $4+6+5+3=18$.
		\end{itemize}
	 }
\end{bt}

\begin{bt}%[0D1V3-4]%[ma:bd954b33]
	Cho các tập hợp $A=\left(2;+\infty \right)$ và $B=\left[m^2-7;+\infty \right)$ với  $m > 0$. Số giá trị nguyên của $m$ để $A\setminus B$ là một khoảng $\left(a; b\right)$ thoả mãn $b-a$ thuộc đoạn $\left[3; 16\right]$ là?
	\shortans{2} 
	 \loigiai{
		Điều kiện để $A\setminus  B\ne \varnothing$ là $\heva{&m^2-7> 2\\
		&m > 0} \Leftrightarrow  \heva{&m^2> 9\\
		&m > 0} 
		\Leftrightarrow m  > 3$.\\
		 Khi đó $A\setminus B=\left(2;m^2-7\right)\Rightarrow a=2$, $b=m^2-7$.\\
		Ta có $b-a$ thuộc đoạn $\left[3; 16\right]\Leftrightarrow 3\le m^2-9\le 16\Leftrightarrow 12\le  m^2\le 25$.\\
		Do $m > 3$ và $m\in \mathbb{Z}$ nên $m\in \left\{4; 5\right\}$.\\ Vậy có  $2$ giá trị $m$ nguyên thoả mãn.
	} 
\end{bt}

\begin{bt}%[0D1K1-3]%[ma:7f419606]
	Cho mệnh đề $P\colon$ \lq\lq Bình phương của một số thực lớn hơn $ 2 $ \rq\rq, tìm số số  nguyên thỏa mãn mệnh đề phủ định của mệnh đề $P$.
	\shortans{3} 
	 \loigiai{
		Ta  có: $P\colon$ \lq\lq $x^2>2$\rq\rq.
		Suy ra $\overline{P}\colon$  \lq\lq $ x^2 \leq 2 $\rq\rq.\\
		Vì $x \in \mathbb{Z}$ nên $x$ là $ -1 $; $ 0 $;  $ 1 $.\\
	Vậy có $ 3 $ số nguyên  thỏa mãn mệnh đề $\overline{P}$.}
\end{bt}

\Closesolutionfile{ansbt}

\hetDe{101}

\phanDapAn{101}
\par\noindent\textbf{PHẦN I.}\par
\begin{flushleft}
\begin{Solution}{1}
A
\end{Solution}
\begin{Solution}{2}
D
\end{Solution}
\begin{Solution}{3}
A
\end{Solution}
\begin{Solution}{4}
D
\end{Solution}
\begin{Solution}{5}
B
\end{Solution}
\begin{Solution}{6}
A
\end{Solution}
\begin{Solution}{7}
D
\end{Solution}
\begin{Solution}{8}
D
\end{Solution}
\begin{Solution}{9}
D
\end{Solution}
\begin{Solution}{10}
D
\end{Solution}
\begin{Solution}{11}
D
\end{Solution}
\begin{Solution}{12}
B
\end{Solution}

\end{flushleft}
\par\noindent\textbf{PHẦN II.}\par
\begin{flushleft}
\def\writeANS{\TLsai{A}\TLsai{B}\TLdung{C}\TLsai{D}}
\begin{Solbook}{1}
 \writeANS 
\end{Solbook}
\def\writeANS{\TLdung{A}\TLsai{B}\TLdung{C}\TLdung{D}}
\begin{Solbook}{2}
 \writeANS 
\end{Solbook}
\def\writeANS{\TLdung{A}\TLsai{B}\TLsai{C}\TLdung{D}}
\begin{Solbook}{3}
 \writeANS 
\end{Solbook}
\def\writeANS{\TLdung{A}\TLsai{B}\TLsai{C}\TLsai{D}}
\begin{Solbook}{4}
 \writeANS 
\end{Solbook}

\end{flushleft}
\par\noindent\textbf{PHẦN III.}\par
\begin{flushleft}
\begin{loigiaibt}{1}
 1 
\end{loigiaibt}
\begin{loigiaibt}{2}
 14 
\end{loigiaibt}
\begin{loigiaibt}{3}
 18 
\end{loigiaibt}
\begin{loigiaibt}{4}
 2 
\end{loigiaibt}
\begin{loigiaibt}{5}
 3 
\end{loigiaibt}

\end{flushleft}

\newpage
\dauDe{102}

\phanDe{\textbf{PHẦN I. Câu trắc nghiệm nhiều phương án lựa chọn.} Thí sinh trả lời từ câu 1 đến câu 12. Mỗi câu hỏi thí sinh chỉ chọn một phương án.}

\setcounter{ex}{0}

\Opensolutionfile{ans}[Ans/GHK1_12_102_TN]

\begin{ex}%[0D1B2-1]%[ma:f3f80b92]
	Cho tập hợp $A = \left\{n \in \mathbb{Z}\big| \,6 \, \vdots \, n\right\}$. Hãy  viết tập hợp $A$ dưới dạng liệt kê?
	\choice{$A = \left\{1;2;3;6\right\}$}
		 {\True $A =  \left\{-6;-3;-2;-1;1;2;3;6\right\}$}
		{$A = \left\{-6;-3;-2;-1\right\}$} 
		{$A  = \left\{-6;-1;1;6\right\}$}
	 \loigiai{
		$A$ là tập các ước nguyên  của $6$, do đó $A = \left\{-6;-3;-2;-1;1;2;3;6\right\}$.
	} 
\end{ex}

\begin{ex}%[0D1B1-3]%[ma:5c0c30c1]
	Mệnh đề nào sau đây là phủ định của  mệnh đề “Mọi động vật đều di chuyển” ?
	 \choice
		{ Mọi động vật đều không di  chuyển}
		{ Mọi động vật đều đứng  yên}
		{ \True Có  ít nhất một động vật không di chuyển}
		{ Có ít nhất một động vật di  chuyển}
	 \loigiai{
		Theo định nghĩa, phủ định của mệnh đề “Mọi động vật đều di chuyển” là mệnh đề “Có ít nhất một động vật không di  chuyển”.
	 }
\end{ex}

\begin{ex}%[0D1B3-1]%[ma:a97e4b36]
	Cho hai tập hợp $A=\left\{1;2;3;4\right\}$, $B=\left\{ 2;4;6;8 \right\}$. Tập hợp nào sau đây  bằng tập hợp $A\cap B$?
	 \choice
		{$\left\{1;2;3;4;6;8 \right\}$}
		 {$\left\{6;8  \right\}$}
		{$\left\{1;3 \right\}$}
		 {\True$\left\{ 2;4  \right\}$}
	\loigiai{ 
		Ta có $A\cap B=\left\{ 2;4 \right\}.$ 
	 }
\end{ex}

\begin{ex}%[0D1B1-5]%[ma:ca8f23aa]
	Cho mệnh đề ``$\forall x \in  \mathbb{R} : x^2+x+1>0$''. Mệnh đề phủ định của mệnh đề đã cho là
	\choice{\True ``$\exists x \in \mathbb{R} :x^2+x+1  \le 0$''} 
		 {``$\forall x \in \mathbb{R} :x^2+x+1 < 0$''}
		{``$\forall x \in \mathbb{R} :x^2+x+1 = 0$''}
		{``$\exists x  \in \mathbb{R} :x^2+x+1 < 0$''}
	\loigiai{ Mệnh đề phủ định của mệnh đề ``$\forall x \in \mathbb{R} : x^2+x+1>0$'' là mệnh đề  ``$\exists x \in \mathbb{R} :x^2+x+1 \le 0$''.
	 }
\end{ex}

\begin{ex}%[0D1B3-3]%[ma:5f619861]
	Để thành lập đội tuyển học sinh giỏi khối $ 10 $, nhà trường tổ chức thi chọn các môn Toán, Văn, Anh trên tổng số $ 111 $ học sinh. Kết quả  có: $ 70 $ học sinh giỏi Toán, $ 65 $ học sinh giỏi Văn, $ 62 $ học sinh giỏi Anh. Trong đó có $ 49 $ học sinh giỏi cả hai môn Văn và Toán, $ 32 $ học sinh giỏi cả hai môn Toán và Anh, $ 34 $ học sinh giỏi cả hai môn Văn và Anh. Xác định số học sinh giỏi cả ba môn Văn, Toán, Anh. Biết rằng có $ 6 $ học sinh không đạt yêu cầu cả ba môn.
	\choice 
		 {$ 36 $}
		{$  32 $}
		{\True $ 23 $}
		 {$ 34  $}
	 \loigiai{
	 \immini
		{ 
		Có $ 111-6=105 $ học sinh thi đạt ít nhất $ 1 $ môn.\\ 
		Gọi $ A $ là số học sinh giỏi môn  Toán và Tiếng Anh nhưng không giỏi Văn.\\
		Gọi $ B $ là số học sinh giỏi môn Toán và Văn nhưng không giỏi Tiếng Anh.\\ 
		Gọi $ C  $ là số học sinh giỏi môn Văn và Tiếng Anh nhưng không giỏi Toán.\\
		Gọi $ D $ là  số học sinh giỏi cả ba môn. Ta có hệ:\\
		$\begin{cases}
		 B+D=49\\
		A+D=32\\ 
		 C+D=34\\
		70+65+62-(A+B+C+2D)=105 
	\end{cases}\\
	\Rightarrow  92=32-D+49-D+34-D+2D\\
	\Rightarrow D=23$.\\ 
	Vậy có $  23 $ học sinh giỏi cả ba môn.
	 }
	{ 
	\begin{tikzpicture}[scale=.8]
			\def\radius{2cm} 
			\coordinate  (ceni);
			\coordinate[xshift=.8*\radius,yshift=.2*\radius] (cenii); 
			\coordinate[xshift=.6*\radius,yshift=-.8*\radius] (ceniii); 
			\draw (ceni)  circle (.9*\radius);
			\draw (cenii)  circle (\radius);
			\draw (ceniii)  circle (\radius);
			 \node[xshift=.3*\radius,yshift=.25*\radius] at (ceni) {A};
			\node[xshift=.1*\radius,yshift=-.5*\radius] at (ceni) {B}; 
			\node[xshift=.45*\radius,yshift=-.25*\radius] at (ceni) {D}; 
			\node[xshift=.1*\radius,yshift=-.55*\radius] at (cenii)  {C};
		\end{tikzpicture}
	} 
} 
\end{ex}

\begin{ex}%[0D1N2-3]%[ma:b26dadc3]
	Trong các tập hợp sau, tập nào là  tập rỗng?
	 \choice
		 {$D=\{ x \in \mathbb{R} |~ x^2+5x-1=0 \}$}
		{$A= \{ x \in \mathbb{R} |~ x^2+5x-6=0 \}$}
		{$B=\{  x \in \mathbb{Q} |~ 3x^2-5x+2=0 \}$}
		{\True $C=\{ x\in \mathbb{Z} |~ x^2+x-1=0  \}$} 
	 \loigiai{
		Ta có\\ 
		 $x^2+5x-6=0
		\Leftrightarrow \hoac{&x=1 \in \mathbb{R}\\&x=-6 \in \mathbb{R}}$. Vậy $A=\{ -6;1  \} \neq \varnothing$.\\
		$3x^2-5x+2=0 
		\Leftrightarrow \hoac{&x=1 \in \mathbb{Q}\\&x=\dfrac{2}{3} \in \mathbb{Q}}$. Vậy $B=\left\{ 1;\dfrac{2}{3}  \right\} \neq \varnothing$.\\
		 $x^2+x-1=0
		\Leftrightarrow \hoac{&x=\dfrac{-1+\sqrt{5}}{2} \notin \mathbb{Z}\\&x=\dfrac{-1-\sqrt{5}}{2} \notin \mathbb{Z}}$. Vì $x \in \mathbb{Z}$  nên $C=\varnothing$.\\
		$x^2+5x-1=0 
		\Leftrightarrow \hoac{&x=\dfrac{-5+\sqrt{29}}{2} \in \mathbb{R}\\&x=\dfrac{-5-\sqrt{29}}{2} \in \mathbb{R}}$.  Vậy $D=\left\{ \dfrac{-5+\sqrt{29}}{2};\dfrac{-5-\sqrt{29}}{2} \right\} \neq \varnothing$.
	} 
\end{ex}

\begin{ex}%[0D1H2-3]%[ma:e60745e5]
	Tìm tất cả các giá trị của  tham số $m$ sao cho $(m;+\infty)\subset(1;+\infty)$?
	 \choice
		{\True $m\geq 1$} 
		 {$m<1$}
		{$m>1$} 
		 {$m\leq 1$}
	 \loigiai{
		Do $(m;+\infty)\subset(1;+\infty)$  nên $m\geq 1$.
	} 
\end{ex}

\begin{ex}%[0D1H3-3]%[ma:76b2c2a1]
	Tổng các  số nguyên của tập $[-3;2]\setminus(1;2)$ bằng
	 \choice
		 {\True$-3$}
		 {$6$}
		 {$-5$}
		{$4$} 
	 \loigiai{
		Ta có $[-3;2]\setminus(1;2)=[3;1]\cup\{2\}$ nên chứa các số nguyên  là $-3;-2;-1;0;1;2$.\\
		Vậy tổng các số nguyên là $-3$. 
	 }
\end{ex}

\begin{ex}%[0D1H1-5]%[ma:9a292bba]
	Trong các mệnh đề dưới đây, mệnh đề nào  là mệnh đề đúng?
	 \choice
		{\True \lq\lq $\exists x\in \mathbb{N}\mid 2x^2-1<  0$\rq\rq} 
		{\lq\lq $\exists x\in  \mathbb{R}\mid x^2< 0$\rq\rq}
		{\lq\lq $\exists x\in \mathbb{R}\mid x^2+1=0$\rq\rq}
		{\lq\lq $\exists x\in \mathbb{N}\mid x^2-2=0$\rq\rq} 
	 \loigiai{
		Ta  có $x^2\ge 0\Leftrightarrow x^2+1\ge 1$, $\forall x\in \mathbb{R}$.
		Nên  $\exists x\in \mathbb{R}\mid x^2+1=0$ là mệnh đề sai.\\
		Ta có $x^2\ge 0$, $\forall x\in \mathbb{R}$. Nên $\exists x\in \mathbb{R}\mid x^2< 0$ là  mệnh đề sai.\\
		$2x^2-1< 0\Leftrightarrow x^2< \dfrac{1}{2} \Leftrightarrow-\dfrac{\sqrt{2}}{2} < x < \dfrac{\sqrt{2}}{2}$, mà $x\in \mathbb{N}\Rightarrow x=0$. Vậy $\exists x\in \mathbb{N}\mid 2x^2-1< 0$ là mệnh  đề đúng.\\
		$x^2-2=0\Leftrightarrow x=\pm \sqrt{2} \notin \mathbb{N}$. Vậy $\exists x\in  \mathbb{N}\mid x^2-2=0$ là mệnh đề sai.
	 }
\end{ex}

\begin{ex}%[0D1H3-5]%[ma:1d662c9d]
	Lớp 10A có $45$ học sinh trong đó có $25$ em học giỏi môn Toán, $23$ em học giỏi  môn Lý, $20$ em học giỏi môn Hóa, $11$ em học giỏi cả môn Toán và môn Lý, $8$ em học giỏi cả môn Lý và môn Hóa, $9$ em học giỏi cả môn Toán và môn Học. Hỏi lớp $10$ A có bao nhiêu bạn học giỏi cả ba môn Toán, Lý, Hóa, biết rằng mỗi học sinh trong lớp học giỏi ít nhất một trong $3$ môn Toán, Lý, Hóa?
	\choice 
		{$4$} 
		 {$3$}
		 {$6$}
		{\True $5$} 
	 \loigiai{
		\begin{center}
			\begin{tikzpicture}[scale=.8,font=\footnotesize,line join=round,line cap=round,>=stealth]
					\coordinate  (A) at (0,0);
					 \coordinate (B) at (2,0);
					\coordinate  (C) at (1,-1.5);
					\begin{scope}
						 \clip (C) circle (1.5);
						 \clip (A) circle (1.5);
						\fill[pattern=north west lines]  (B) circle (1.5);
					\end{scope}
					\draw (A)  circle (1.5);
					\draw (B)  circle (1.5);
					\draw (C) circle  (1.5);
					 \path
						 (A)++(90:.6) node[fill=white,above]{$L$}
						(B)++(90:.6)  node[fill=white,above]{$T$}
						(C)++(-90:1)  node[fill=white,above]{$H$};
				\end{tikzpicture}
		\end{center}
		Gọi $T$, $L$, $H$ lần lượt là tập hợp các học sinh giỏi môn Toán,  Lý, Hóa.\\
		Khi đó 
		\begin{eqnarray*}
			|T \cup L \cup H|&=&|T|+|L|+|H|-|T \cap L|-|L \cap H|-|H \cap T|+|T \cap L  \cap H| \\
			 45&=&25+23+20-11-8-9+|T \cap L \cap H|\\
			|T  \cap L \cap H|&=&5
		\end{eqnarray*}
		Vậy có $5$ học  sinh giỏi cả $3$ môn.
	 }
\end{ex}

\begin{ex}%[0D1K1-3]%[ma:a990ffed]
	Phát biểu mệnh đề  phủ định của mệnh đề "Tồn tại một loài thú có nhiệt độ thân thể nhỏ hơn $35^\circ$C".
	\choice 
		{"Mọi loài thú có nhiệt độ thân thể  nhỏ hơn $35^\circ$C"}
		{"Tồn tại một loài thú có nhiệt  độ thân thể lớn hơn $35^\circ$C"}
		{\True "Mọi loài thú có nhiệt độ thân thể  không nhỏ hơn $35^\circ$C"}
		{"Mọi loài thú có nhiệt độ thân thể lớn hơn $35^\circ$C"} 
	\loigiai{Mệnh đề phủ định của mệnh đề "Tồn tại một loài thú có nhiệt độ thân thể nhỏ hơn $35^\circ$C" là  mệnh đề "Mọi loài thú có nhiệt độ thân thể không nhỏ hơn $35^\circ$C".}
\end{ex}

\begin{ex}%[0D1K1-5]%[ma:2b89aaa2]
	Trong các mệnh đề sau,  mệnh đề nào \textbf{sai}?
	\choice 
		{$\exists n\in \mathbb{N}$, $n^2+11n+2$ chia  hết cho $11$}
		{$\exists n\in \mathbb{Z}$, $2x^2-8=0$} 
		{Tồn tại số nguyên  tố chia hết cho $5$}
		 {\True $\exists n\in \mathbb{N}$, $n^2+1$ chia hết cho $4$}
	 \loigiai{
		\begin{itemize}
			\item Xét  mệnh đề ``$\exists n\in \mathbb{N}$, $n^2+11n+2$ chia hết cho $11$''.\\
			Khi $n=3$ thì  giá trị của $(n^2+11n+2)$ bằng $44\vdots 11$ nên mệnh đề đúng.
			\item Xét mệnh đề ``$\exists n\in \mathbb{N}$, $n^2+1$  chia hết cho $4$''.\\
			Khi  $n=2k,k\in \mathbb{N} \Rightarrow n^2+1=4k^2+1$ không chia hết cho $4$, $k\in \mathbb{N}$.\\
			Khi $n=2k+1,k\in \mathbb{N} \Rightarrow n^2+1=(2k+1)^2+1=4k^2+4k+2$ không  chia hết cho $4$, $k\in \mathbb{N}$.
			\item Xét mệnh đề ``Tồn tại  số nguyên tố chia hết cho $5$''.\\
			Tồn tại số nguyên tố $5$ chia hết cho $5$  nên mệnh đề đúng.
			\item  Xét mệnh đề ``$\exists n\in \mathbb{Z}$, $2x^2-8=0$ ''.\\
			Phương trình $2x^2-8=0 \Leftrightarrow x^2=4 \Leftrightarrow x=-2;x=2\in  \Bbb{Z}$ nên mệnh đề đúng.
		\end{itemize}
	 }
\end{ex}

\Closesolutionfile{ans}

\phanDe{\textbf{PHẦN II. Câu trắc nghiệm đúng sai.} Thí sinh trả lời từ câu 1 đến câu 4. Trong mỗi ý a), b), c), d) ở mỗi câu, thí sinh chọn đúng hoặc sai.}

\setcounter{ex}{0}

\Opensolutionfile{ansbook}[Ansbook/GHK1_12_102_DS]

\begin{ex}%[0D1H1-3]%[ma:f76fc24c]
	Cho  mệnh đề $P(x):$ \lq\lq $\forall x\in \mathbb{R},x^2-4x+1> 0$\rq\rq.
	 \choiceTF
		{$P(x)$: \lq\lq Có ít nhất một số thực $x$ thoả mãn $x^2-4x+1> 0$\rq\rq} 
		{$P(x)$ là mệnh đề chứa biến} 
		{\True  $\overline{P(x)}$: \lq\lq $\exists x\in \mathbb{R}, x^2-4x+1\le 0$\rq\rq}
		{Mệnh đề $\overline{P(x)}$ là mệnh  đề sai}
	 \loigiai{
		\begin{itemchoice}
			\itemch  SAI.\\
			Ta có phát biểu đúng của mệnh đề $P(x):$\lq\lq Với mọi số thực  $x$ thoả mãn $x^2-4x+1> 0$\rq\rq.
			\itemch SAI.\\ 
			$P(x)$ là mệnh đề sai nên không phải là  mệnh đề chứa biến.
			\itemch  ĐÚNG.\\
			Mệnh đề phủ định  của $P(x)$ là $\overline{P(x)}:$ \lq\lq $\exists x\in \mathbb{R}, x^2-4x+1\le 0$\rq\rq.
			\itemch SAI.\\ 
			Ta có $\overline{P(x)}:$ \lq\lq $\exists x\in \mathbb{R}, x^2-4x+1\le 0$\rq\rq. Mệnh đề $\overline{P(x)}$ đúng vì giả sử ta  chọn $x=1$ thì $1^2-4.1+1=-2< 0$.
		\end{itemchoice}
	} 
\end{ex}

\begin{ex}%[0D1H1-2]%[ma:bff4f470]
	Cho $\Delta ABC$ có $A'$, $B'$, $C'$ lần lượt là các trung điểm của các  cạnh $BC$, $CA$, $AB$. Khi đó
	 \choiceTF
		{\True $B{C'}=C' A=A'{B'}=\dfrac{AB}{2}{\rm}$} 
		{Hai vectơ $\overrightarrow{BC'}$,  $\overrightarrow{A'B'}$ ngược hướng}
		{\True $\overrightarrow{B{C'}}=\overrightarrow{C' A} =\overrightarrow{A'B'}$} 
		{\True $\overrightarrow{B'  C'}=\overrightarrow{CA'}$}
	\loigiai{ 
		\begin{center}
			\begin{tikzpicture}[scale=1,line join=round,line cap=round,font=\footnotesize,>=stealth]
					\def\r{3} 
					\path (180:\r)  coordinate (B)
						(0:\r)  coordinate (C)
						(120:\r) coordinate  (A)
						($(B)!0.5!(C)$) coordinate  (A')
						($(C)!0.5!(A)$) coordinate (B') 
						($(B)!0.5!(A)$)  coordinate (C');
					\draw  (A)--(B)--(C)--cycle;
					\foreach \x/ \goc in  {A/135,B/-135,C/-45,A'/-90,B'/30,C'/180}
					 \fill (\x) circle (1pt)
						 ($(\x)+(\goc:3mm)$) node {$\x$};
				\end{tikzpicture}
		\end{center}
		\begin{itemchoice}
			\itemch Ta có $C'$ là trung điểm của $AB$ và $A'B'$ là đường trung bình của tam giác $ABC$ ứng với  cạnh đáy $AB$ nên $BC'=C'A=A'B'=\dfrac{AB}{2}$.
			\itemch  Vì Hai vectơ $\overrightarrow{BC'}$, $\overrightarrow{A'{B'}}$ cùng hướng.
			\itemch Vì $C'$ là trung điểm của $AB$ và $A'B'$ là đường trung bình của tam giác $ABC$ ứng với cạnh đáy $AB$ nên $\overrightarrow{BC'}=\overrightarrow{C'A}  =\overrightarrow{A'B'}$.
			\itemch Vì $A'$ là trung điểm của $BC$ và $B'C'$ là đường trung bình  của tam giác $ABC$ ứng với cạnh đáy $BC$ nên $\overrightarrow{B'C'}=\overrightarrow{CA'}$.
		\end{itemchoice}
	 }
\end{ex}

\begin{ex}%[0D1V1-2]%[ma:d97640b0]
	Cho hình thang $ABCD$ vuông tại $A$ và có $AB=AD=\dfrac{1}{2}DC=a$. Gọi $BF$ là  đường phân giác trong của tam giác $ABD$ ($F\in AD$). Khi đó
	 \choiceTF
		{\True  $CA^2=DA^2+DC^2$}
		{$\left|\overrightarrow{CA}\right|=a\sqrt  3 $}
		 {$\widehat{ABF}=45^\circ$}
		{$\left|\overrightarrow{BF}\right|\approx 2{,}08a$} 
	\loigiai{\begin{center}
		\begin{tikzpicture}[>=stealth,line join=round,line cap=round,font=\footnotesize,scale=0.8]
				\def\a{2.4} 
				\path 
					 (0,4) coordinate (A)
					(0,0) coordinate  (D)
					(8,0) coordinate (C) 
					(4,4)  coordinate (B)
					($(A)!1/\a!(D)$) coordinate  (F);
				 \draw (A)--(B)--(C)--(D)--(A)
					 (A)--(C)
					(F)--(B)--(D); 
				\foreach \x/\g in  {B/60,A/120,C/30,D/-60,F/-140} \fill (\x) circle (1pt) +(\g:3mm) node[scale=.8] {$\x$};
				 \draw pic[thin,draw,angle radius=2mm] {right angle= B--A--D};
				 \draw pic[thin,draw,angle radius=2mm] {right angle= A--D--C};
			\end{tikzpicture}
	\end{center}
	\begin{itemchoice}
		\itemch Áp dụng định lí Pytago cho tam giác $ADC$ vuông  tại $D$.
		\itemch Áp dụng  định lí Pytago cho tam giác $ADC$ vuông tại $D$, ta có\\
		$C{A^2}=D{A^2}+D{C^2}=a^2+(2a)^2=5a^2$\\ 
		 $\Rightarrow \left|\overrightarrow{CA}\right|=CA=a\sqrt 5$.
		\itemch Ta có  $\Delta ABD$ vuông cân tại $A$, do đó\\
		$\widehat{ABD}=45^\circ  \Rightarrow \widehat{ABF}=22{,}5^\circ$.
		\itemch Xét  $\Delta ABF$ vuông tại $A$, ta có\\
		$ \left| \overrightarrow{BF} \right| =BF=\dfrac{AB}{\cos \left(  ABF \right)}=\dfrac{a}{\cos{22{,}5^\circ}}\approx 1{,}08a$.
	\end{itemchoice}
 }
\end{ex}

\begin{ex}%[0D1V2-1]%[ma:161eac9d]
	 Cho các tập hợp sau: A các số nguyên tố nhỏ hơn 11; $B=\left\{\left. x\in \mathbb{R}\right|3x^2-4x+1=0\right\}$;
	$C=\left\{\left. x\in  \mathbb{N}\right|\left(x^2-5x+6\right)\left(2x+1\right)=0\right\}$; $D=\{\left. x\in \mathbb{Z}\right|\left|x+1\right|<3\}$.
	 \choiceTF
		{\True Tập hợp $A$ có $4$ phần tử} 
		{Tập hợp $B$ có  $3$ phần tử}
		{Tập hợp  $C$ có $3$ phần tử}
		{\True Tập  hợp $D$ có $3$ phần tử}
	 \loigiai{
		\begin{itemchoice}
			\itemch Ta  có các số nguyên tố nhỏ hơn $11$ là $2$; $3$; $5$; $7$. \\
			Vậy $A=\{2; 3; 5; 7\}$. 
			\itemch Ta có $3x^2-4x+1=0\Leftrightarrow \hoac{ 
			&x=1\in \mathbb{R} \\ 
			 &x=\dfrac{1}{3}\in \mathbb{R}.}$ \\
			Vậy $B=\left\{\dfrac{1}{3};  1\right\}$.
			 \itemch Ta có $\left(x^2-5x+6\right)\left(2x+1\right)=0\Leftrightarrow \hoac{
			& x^2-5x+6=0 \\ 
			& 2x+1=0}\Leftrightarrow  \hoac{
			 & x=2\in \mathbb{N} \\
			& x=3\in \mathbb{N}  \\
			&  x=\dfrac{-1}{2}\notin \mathbb{N}.} $\\
			Vậy  $C=\left\{2; 3\right\}$.
			 \itemch Ta có $\heva{
			&x\in \mathbb{Z}  \\
			 &|x+1|<2}\Rightarrow \hoac{
			&x=-2  \\
			 &x=-1 \\
			 &x=0.} $\\
			Vậy $D=\{-2;-1; 0\}$. 
		\end{itemchoice}
	} 
\end{ex}

\Closesolutionfile{ansbook}

\phanDe{\textbf{PHẦN III. Câu trắc nghiệm trả lời ngắn.} Thí sinh trả lời từ câu 1 đến câu 5.}

\setcounter{bt}{0}

\Opensolutionfile{ansbt}[Ansbook/GHK1_12_102_TLN]

\begin{bt}%[0D1H1-6]%[ma:deff5465]
	Cho phát biểu $\forall x\in\mathbb{R}$, $x^{2}+a-5>0$ với $a$ là số thực cho trước. Có bao giá trị  nguyên của $a<20$ để phát biểu đã cho là mệnh đề đúng?
	\shortans{14} 
	 \loigiai{
		Mệnh đề $\forall x\in\mathbb{R}$, $x^{2}+a-5>0$ đúng $\Leftrightarrow  x^{2}>5-a$ đúng $\forall x\in\mathbb{R}$.\\
		Điều này xảy ra khi và chỉ khi  $\max(5-a, 0) < 0$.\\
		Vì $x^2 \ge 0$, nên ta cần  $5-a$ phải nhỏ hơn $0$.\\
		 $5-a<0\Leftrightarrow a>5$.\\
		Kết hợp  với điều kiện $a<20$, ta có $5<a<20$.\\
		Các giá trị nguyên của $a$ là $6, 7, 8, \dots,  19$.\\
		Số giá trị nguyên của  $a$ là $19-6+1=14$.
	} 
\end{bt}

\begin{bt}%[0D1H1-3]%[ma:d05bebba]
	Cho mệnh đề $P$: \lq\lq $x^2+x+1$ là số dương với $x$ là số thực\rq\rq\ và các  mệnh đề sau.
	\begin{itemize}
		\item \lq\lq $x^2+x+1$  là số không dương với $x$ là số thực\rq\rq.
		\item \lq\lq $x^2+x+1$ là khác $0$ với $x$ là số thực\rq\rq. 
		\item \lq\lq $x^2+x+1$ là số âm với $x$ là số thực\rq\rq. 
		\item \lq\lq $x^2+x+1$  là số không âm với $x$ là số thực\rq\rq.
	\end{itemize}
	Có bao nhiêu phát biểu là phủ định của mệnh  đề $P$?
	\shortans{1} 
	 \loigiai{
		Phủ định của mệnh đề $P$ là $\overline{P}$: \lq\lq $x^2+x+1$ là số không dương  với $x$ là số thực\rq\rq.\\
		Vậy  có $1$ mệnh đề thỏa. Các mệnh đề còn lại không phải là phủ định của mệnh đề $P$.
	 }
\end{bt}

\begin{bt}%[0D1V2-2]%[ma:8c3a2e78]
	Lớp $10$A có $10$ học sinh giỏi Toán, $10$ học sinh giỏi Lý, $11$ học sinh giỏi hóa, $6$ học sinh giỏi cả Toán và Lý, $5$ học sinh giỏi cả Hóa và Lý, $4$ học sinh giỏi cả Toán và Hóa, $3$ học sinh giỏi cả ba môn Toán, Lý,  Hóa. Tính học sinh giỏi ít nhất một trong ba môn (Toán, Lý, Hóa) của lớp $10$A?
	 \shortans{18}
	\loigiai{ 
		Theo giả thiết đề bài cho, ta có  biểu đồ Ven
		\begin{center}
			\begin{tikzpicture}[>=stealth,line join=round,line cap=round,font=\footnotesize,scale=1]
					\draw [thick,green!60!black] (0,0) circle (1.5) 
						(1.5,-0.5)  circle (1)
						(1.5,0.5) circle  (0.8);
					\draw (1,.8) node {$6$} 
						 (1.2,.2) node {$3$}
						(2,.8) node {Lý} 
						 (1.8,.1) node {$5$}
						(1.8,-1)  node {Hoá}
						(.8,-.8)  node {$4$}
						(-1,0) node  {Toán};
				\end{tikzpicture}
		\end{center}
		Dựa vào biểu đồ Ven, ta có học sinh giỏi ít nhất một trong ba môn (Toán, Lý, Hóa)  của lớp $10$A là
		\begin{itemize}
			\item Số học  sinh giỏi Toán $6+4+3=13$.
			\item Số học sinh  giỏi Lý $6+5+3=14$.
			\item Số học  sinh giỏi Hóa $4+5+3=12$.
		\end{itemize}
		Ta lại có 
		\begin{itemize}
			\item Số học  sinh giỏi cả Toán và Lý $6$.
			\item Số học sinh giỏi cả Toán và Hóa $4$. 
			\item Số học sinh giỏi cả Hóa  và Lý $5$.
			\item Và số học  sinh giỏi cả Toán, Lý và Hóa là $3$.
			\item Số học sinh giỏi hơn một môn  là $4+6+5+3=18$.
		\end{itemize}
	 }
\end{bt}

\begin{bt}%[0D1V3-4]%[ma:bd954b33]
	Cho các tập hợp $A=\left(2;+\infty \right)$ và $B=\left[m^2-7;+\infty \right)$ với  $m > 0$. Số giá trị nguyên của $m$ để $A\setminus B$ là một khoảng $\left(a; b\right)$ thoả mãn $b-a$ thuộc đoạn $\left[3; 16\right]$ là?
	\shortans{2} 
	 \loigiai{
		Điều kiện để $A\setminus  B\ne \varnothing$ là $\heva{&m^2-7> 2\\
		&m > 0} \Leftrightarrow  \heva{&m^2> 9\\
		&m > 0} 
		\Leftrightarrow m  > 3$.\\
		 Khi đó $A\setminus B=\left(2;m^2-7\right)\Rightarrow a=2$, $b=m^2-7$.\\
		Ta có $b-a$ thuộc đoạn $\left[3; 16\right]\Leftrightarrow 3\le m^2-9\le 16\Leftrightarrow 12\le  m^2\le 25$.\\
		Do $m > 3$ và $m\in \mathbb{Z}$ nên $m\in \left\{4; 5\right\}$.\\ Vậy có  $2$ giá trị $m$ nguyên thoả mãn.
	} 
\end{bt}

\begin{bt}%[0D1K1-3]%[ma:7f419606]
	Cho mệnh đề $P\colon$ \lq\lq Bình phương của một số thực lớn hơn $ 2 $ \rq\rq, tìm số số  nguyên thỏa mãn mệnh đề phủ định của mệnh đề $P$.
	\shortans{3} 
	 \loigiai{
		Ta  có: $P\colon$ \lq\lq $x^2>2$\rq\rq.
		Suy ra $\overline{P}\colon$  \lq\lq $ x^2 \leq 2 $\rq\rq.\\
		Vì $x \in \mathbb{Z}$ nên $x$ là $ -1 $; $ 0 $;  $ 1 $.\\
	Vậy có $ 3 $ số nguyên  thỏa mãn mệnh đề $\overline{P}$.}
\end{bt}

\Closesolutionfile{ansbt}

\hetDe{102}

\phanDapAn{102}
\par\noindent\textbf{PHẦN I.}\par
\begin{flushleft}
\begin{Solution}{1}
B
\end{Solution}
\begin{Solution}{2}
C
\end{Solution}
\begin{Solution}{3}
D
\end{Solution}
\begin{Solution}{4}
A
\end{Solution}
\begin{Solution}{5}
C
\end{Solution}
\begin{Solution}{6}
D
\end{Solution}
\begin{Solution}{7}
A
\end{Solution}
\begin{Solution}{8}
A
\end{Solution}
\begin{Solution}{9}
A
\end{Solution}
\begin{Solution}{10}
D
\end{Solution}
\begin{Solution}{11}
C
\end{Solution}
\begin{Solution}{12}
D
\end{Solution}

\end{flushleft}
\par\noindent\textbf{PHẦN II.}\par
\begin{flushleft}
\def\writeANS{\TLsai{A}\TLsai{B}\TLdung{C}\TLsai{D}}
\begin{Solbook}{1}
 \writeANS 
\end{Solbook}
\def\writeANS{\TLdung{A}\TLsai{B}\TLdung{C}\TLdung{D}}
\begin{Solbook}{2}
 \writeANS 
\end{Solbook}
\def\writeANS{\TLdung{A}\TLsai{B}\TLsai{C}\TLsai{D}}
\begin{Solbook}{3}
 \writeANS 
\end{Solbook}
\def\writeANS{\TLdung{A}\TLsai{B}\TLsai{C}\TLdung{D}}
\begin{Solbook}{4}
 \writeANS 
\end{Solbook}

\end{flushleft}
\par\noindent\textbf{PHẦN III.}\par
\begin{flushleft}
\begin{loigiaibt}{1}
 14 
\end{loigiaibt}
\begin{loigiaibt}{2}
 1 
\end{loigiaibt}
\begin{loigiaibt}{3}
 18 
\end{loigiaibt}
\begin{loigiaibt}{4}
 2 
\end{loigiaibt}
\begin{loigiaibt}{5}
 3 
\end{loigiaibt}

\end{flushleft}

\newpage
\dauDe{103}

\phanDe{\textbf{PHẦN I. Câu trắc nghiệm nhiều phương án lựa chọn.} Thí sinh trả lời từ câu 1 đến câu 12. Mỗi câu hỏi thí sinh chỉ chọn một phương án.}

\setcounter{ex}{0}

\Opensolutionfile{ans}[Ans/GHK1_12_103_TN]

\begin{ex}%[0D1N2-3]%[ma:b26dadc3]
	Trong các tập hợp sau, tập nào là  tập rỗng?
	 \choice
		 {$B=\{  x \in \mathbb{Q} |~ 3x^2-5x+2=0 \}$}
		{$D=\{ x \in \mathbb{R} |~ x^2+5x-1=0 \}$}
		{\True $C=\{ x\in \mathbb{Z} |~ x^2+x-1=0  \}$}
		{$A= \{ x \in \mathbb{R} |~ x^2+5x-6=0 \}$} 
	 \loigiai{
		Ta có\\ 
		 $x^2+5x-6=0
		\Leftrightarrow \hoac{&x=1 \in \mathbb{R}\\&x=-6 \in \mathbb{R}}$. Vậy $A=\{ -6;1  \} \neq \varnothing$.\\
		$3x^2-5x+2=0 
		\Leftrightarrow \hoac{&x=1 \in \mathbb{Q}\\&x=\dfrac{2}{3} \in \mathbb{Q}}$. Vậy $B=\left\{ 1;\dfrac{2}{3}  \right\} \neq \varnothing$.\\
		 $x^2+x-1=0
		\Leftrightarrow \hoac{&x=\dfrac{-1+\sqrt{5}}{2} \notin \mathbb{Z}\\&x=\dfrac{-1-\sqrt{5}}{2} \notin \mathbb{Z}}$. Vì $x \in \mathbb{Z}$  nên $C=\varnothing$.\\
		$x^2+5x-1=0 
		\Leftrightarrow \hoac{&x=\dfrac{-5+\sqrt{29}}{2} \in \mathbb{R}\\&x=\dfrac{-5-\sqrt{29}}{2} \in \mathbb{R}}$.  Vậy $D=\left\{ \dfrac{-5+\sqrt{29}}{2};\dfrac{-5-\sqrt{29}}{2} \right\} \neq \varnothing$.
	} 
\end{ex}

\begin{ex}%[0D1B1-5]%[ma:ca8f23aa]
	Cho mệnh đề ``$\forall x \in  \mathbb{R} : x^2+x+1>0$''. Mệnh đề phủ định của mệnh đề đã cho là
	\choice{``$\exists x  \in \mathbb{R} :x^2+x+1 < 0$''} 
		 {``$\forall x \in \mathbb{R} :x^2+x+1 = 0$''}
		{``$\forall x \in \mathbb{R} :x^2+x+1 < 0$''}
		{\True ``$\exists x \in \mathbb{R} :x^2+x+1  \le 0$''}
	\loigiai{ Mệnh đề phủ định của mệnh đề ``$\forall x \in \mathbb{R} : x^2+x+1>0$'' là mệnh đề  ``$\exists x \in \mathbb{R} :x^2+x+1 \le 0$''.
	 }
\end{ex}

\begin{ex}%[0D1B2-1]%[ma:f3f80b92]
	Cho tập hợp $A = \left\{n \in \mathbb{Z}\big| \,6 \, \vdots \, n\right\}$. Hãy  viết tập hợp $A$ dưới dạng liệt kê?
	\choice{$A  = \left\{-6;-1;1;6\right\}$}
		 {$A = \left\{1;2;3;6\right\}$}
		{$A = \left\{-6;-3;-2;-1\right\}$} 
		{\True $A =  \left\{-6;-3;-2;-1;1;2;3;6\right\}$}
	 \loigiai{
		$A$ là tập các ước nguyên  của $6$, do đó $A = \left\{-6;-3;-2;-1;1;2;3;6\right\}$.
	} 
\end{ex}

\begin{ex}%[0D1B1-3]%[ma:5c0c30c1]
	Mệnh đề nào sau đây là phủ định của  mệnh đề “Mọi động vật đều di chuyển” ?
	 \choice
		{ Có ít nhất một động vật di  chuyển}
		{ Mọi động vật đều đứng  yên}
		{ Mọi động vật đều không di  chuyển}
		{ \True Có  ít nhất một động vật không di chuyển}
	 \loigiai{
		Theo định nghĩa, phủ định của mệnh đề “Mọi động vật đều di chuyển” là mệnh đề “Có ít nhất một động vật không di  chuyển”.
	 }
\end{ex}

\begin{ex}%[0D1B3-1]%[ma:a97e4b36]
	Cho hai tập hợp $A=\left\{1;2;3;4\right\}$, $B=\left\{ 2;4;6;8 \right\}$. Tập hợp nào sau đây  bằng tập hợp $A\cap B$?
	 \choice
		{$\left\{1;3 \right\}$}
		 {\True$\left\{ 2;4  \right\}$}
		{$\left\{6;8  \right\}$}
		 {$\left\{1;2;3;4;6;8 \right\}$}
	\loigiai{ 
		Ta có $A\cap B=\left\{ 2;4 \right\}.$ 
	 }
\end{ex}

\begin{ex}%[0D1B3-3]%[ma:5f619861]
	Để thành lập đội tuyển học sinh giỏi khối $ 10 $, nhà trường tổ chức thi chọn các môn Toán, Văn, Anh trên tổng số $ 111 $ học sinh. Kết quả  có: $ 70 $ học sinh giỏi Toán, $ 65 $ học sinh giỏi Văn, $ 62 $ học sinh giỏi Anh. Trong đó có $ 49 $ học sinh giỏi cả hai môn Văn và Toán, $ 32 $ học sinh giỏi cả hai môn Toán và Anh, $ 34 $ học sinh giỏi cả hai môn Văn và Anh. Xác định số học sinh giỏi cả ba môn Văn, Toán, Anh. Biết rằng có $ 6 $ học sinh không đạt yêu cầu cả ba môn.
	\choice 
		 {$ 34  $}
		{$ 36 $}
		{\True $ 23 $}
		 {$  32 $}
	 \loigiai{
	 \immini
		{ 
		Có $ 111-6=105 $ học sinh thi đạt ít nhất $ 1 $ môn.\\ 
		Gọi $ A $ là số học sinh giỏi môn  Toán và Tiếng Anh nhưng không giỏi Văn.\\
		Gọi $ B $ là số học sinh giỏi môn Toán và Văn nhưng không giỏi Tiếng Anh.\\ 
		Gọi $ C  $ là số học sinh giỏi môn Văn và Tiếng Anh nhưng không giỏi Toán.\\
		Gọi $ D $ là  số học sinh giỏi cả ba môn. Ta có hệ:\\
		$\begin{cases}
		 B+D=49\\
		A+D=32\\ 
		 C+D=34\\
		70+65+62-(A+B+C+2D)=105 
	\end{cases}\\
	\Rightarrow  92=32-D+49-D+34-D+2D\\
	\Rightarrow D=23$.\\ 
	Vậy có $  23 $ học sinh giỏi cả ba môn.
	 }
	{ 
	\begin{tikzpicture}[scale=.8]
			\def\radius{2cm} 
			\coordinate  (ceni);
			\coordinate[xshift=.8*\radius,yshift=.2*\radius] (cenii); 
			\coordinate[xshift=.6*\radius,yshift=-.8*\radius] (ceniii); 
			\draw (ceni)  circle (.9*\radius);
			\draw (cenii)  circle (\radius);
			\draw (ceniii)  circle (\radius);
			 \node[xshift=.3*\radius,yshift=.25*\radius] at (ceni) {A};
			\node[xshift=.1*\radius,yshift=-.5*\radius] at (ceni) {B}; 
			\node[xshift=.45*\radius,yshift=-.25*\radius] at (ceni) {D}; 
			\node[xshift=.1*\radius,yshift=-.55*\radius] at (cenii)  {C};
		\end{tikzpicture}
	} 
} 
\end{ex}

\begin{ex}%[0D1H2-3]%[ma:e60745e5]
	Tìm tất cả các giá trị của  tham số $m$ sao cho $(m;+\infty)\subset(1;+\infty)$?
	 \choice
		{\True $m\geq 1$} 
		 {$m<1$}
		{$m\leq 1$} 
		 {$m>1$}
	 \loigiai{
		Do $(m;+\infty)\subset(1;+\infty)$  nên $m\geq 1$.
	} 
\end{ex}

\begin{ex}%[0D1H1-5]%[ma:9a292bba]
	Trong các mệnh đề dưới đây, mệnh đề nào  là mệnh đề đúng?
	 \choice
		{\lq\lq $\exists x\in  \mathbb{R}\mid x^2< 0$\rq\rq} 
		{\lq\lq $\exists x\in \mathbb{R}\mid x^2+1=0$\rq\rq}
		{\lq\lq $\exists x\in \mathbb{N}\mid x^2-2=0$\rq\rq}
		{\True \lq\lq $\exists x\in \mathbb{N}\mid 2x^2-1<  0$\rq\rq} 
	 \loigiai{
		Ta  có $x^2\ge 0\Leftrightarrow x^2+1\ge 1$, $\forall x\in \mathbb{R}$.
		Nên  $\exists x\in \mathbb{R}\mid x^2+1=0$ là mệnh đề sai.\\
		Ta có $x^2\ge 0$, $\forall x\in \mathbb{R}$. Nên $\exists x\in \mathbb{R}\mid x^2< 0$ là  mệnh đề sai.\\
		$2x^2-1< 0\Leftrightarrow x^2< \dfrac{1}{2} \Leftrightarrow-\dfrac{\sqrt{2}}{2} < x < \dfrac{\sqrt{2}}{2}$, mà $x\in \mathbb{N}\Rightarrow x=0$. Vậy $\exists x\in \mathbb{N}\mid 2x^2-1< 0$ là mệnh  đề đúng.\\
		$x^2-2=0\Leftrightarrow x=\pm \sqrt{2} \notin \mathbb{N}$. Vậy $\exists x\in  \mathbb{N}\mid x^2-2=0$ là mệnh đề sai.
	 }
\end{ex}

\begin{ex}%[0D1H3-5]%[ma:1d662c9d]
	Lớp 10A có $45$ học sinh trong đó có $25$ em học giỏi môn Toán, $23$ em học giỏi  môn Lý, $20$ em học giỏi môn Hóa, $11$ em học giỏi cả môn Toán và môn Lý, $8$ em học giỏi cả môn Lý và môn Hóa, $9$ em học giỏi cả môn Toán và môn Học. Hỏi lớp $10$ A có bao nhiêu bạn học giỏi cả ba môn Toán, Lý, Hóa, biết rằng mỗi học sinh trong lớp học giỏi ít nhất một trong $3$ môn Toán, Lý, Hóa?
	\choice 
		{$6$} 
		 {$4$}
		 {$3$}
		{\True $5$} 
	 \loigiai{
		\begin{center}
			\begin{tikzpicture}[scale=.8,font=\footnotesize,line join=round,line cap=round,>=stealth]
					\coordinate  (A) at (0,0);
					 \coordinate (B) at (2,0);
					\coordinate  (C) at (1,-1.5);
					\begin{scope}
						 \clip (C) circle (1.5);
						 \clip (A) circle (1.5);
						\fill[pattern=north west lines]  (B) circle (1.5);
					\end{scope}
					\draw (A)  circle (1.5);
					\draw (B)  circle (1.5);
					\draw (C) circle  (1.5);
					 \path
						 (A)++(90:.6) node[fill=white,above]{$L$}
						(B)++(90:.6)  node[fill=white,above]{$T$}
						(C)++(-90:1)  node[fill=white,above]{$H$};
				\end{tikzpicture}
		\end{center}
		Gọi $T$, $L$, $H$ lần lượt là tập hợp các học sinh giỏi môn Toán,  Lý, Hóa.\\
		Khi đó 
		\begin{eqnarray*}
			|T \cup L \cup H|&=&|T|+|L|+|H|-|T \cap L|-|L \cap H|-|H \cap T|+|T \cap L  \cap H| \\
			 45&=&25+23+20-11-8-9+|T \cap L \cap H|\\
			|T  \cap L \cap H|&=&5
		\end{eqnarray*}
		Vậy có $5$ học  sinh giỏi cả $3$ môn.
	 }
\end{ex}

\begin{ex}%[0D1H3-3]%[ma:76b2c2a1]
	Tổng các  số nguyên của tập $[-3;2]\setminus(1;2)$ bằng
	 \choice
		 {$4$}
		 {\True$-3$}
		 {$-5$}
		{$6$} 
	 \loigiai{
		Ta có $[-3;2]\setminus(1;2)=[3;1]\cup\{2\}$ nên chứa các số nguyên  là $-3;-2;-1;0;1;2$.\\
		Vậy tổng các số nguyên là $-3$. 
	 }
\end{ex}

\begin{ex}%[0D1K1-3]%[ma:a990ffed]
	Phát biểu mệnh đề  phủ định của mệnh đề "Tồn tại một loài thú có nhiệt độ thân thể nhỏ hơn $35^\circ$C".
	\choice 
		{\True "Mọi loài thú có nhiệt độ thân thể  không nhỏ hơn $35^\circ$C"}
		{"Mọi loài thú có nhiệt độ thân thể lớn hơn $35^\circ$C"}
		{"Mọi loài thú có nhiệt độ thân thể  nhỏ hơn $35^\circ$C"}
		{"Tồn tại một loài thú có nhiệt  độ thân thể lớn hơn $35^\circ$C"} 
	\loigiai{Mệnh đề phủ định của mệnh đề "Tồn tại một loài thú có nhiệt độ thân thể nhỏ hơn $35^\circ$C" là  mệnh đề "Mọi loài thú có nhiệt độ thân thể không nhỏ hơn $35^\circ$C".}
\end{ex}

\begin{ex}%[0D1K1-5]%[ma:2b89aaa2]
	Trong các mệnh đề sau,  mệnh đề nào \textbf{sai}?
	\choice 
		{Tồn tại số nguyên  tố chia hết cho $5$}
		{$\exists n\in \mathbb{N}$, $n^2+11n+2$ chia  hết cho $11$} 
		{$\exists n\in \mathbb{Z}$, $2x^2-8=0$}
		 {\True $\exists n\in \mathbb{N}$, $n^2+1$ chia hết cho $4$}
	 \loigiai{
		\begin{itemize}
			\item Xét  mệnh đề ``$\exists n\in \mathbb{N}$, $n^2+11n+2$ chia hết cho $11$''.\\
			Khi $n=3$ thì  giá trị của $(n^2+11n+2)$ bằng $44\vdots 11$ nên mệnh đề đúng.
			\item Xét mệnh đề ``$\exists n\in \mathbb{N}$, $n^2+1$  chia hết cho $4$''.\\
			Khi  $n=2k,k\in \mathbb{N} \Rightarrow n^2+1=4k^2+1$ không chia hết cho $4$, $k\in \mathbb{N}$.\\
			Khi $n=2k+1,k\in \mathbb{N} \Rightarrow n^2+1=(2k+1)^2+1=4k^2+4k+2$ không  chia hết cho $4$, $k\in \mathbb{N}$.
			\item Xét mệnh đề ``Tồn tại  số nguyên tố chia hết cho $5$''.\\
			Tồn tại số nguyên tố $5$ chia hết cho $5$  nên mệnh đề đúng.
			\item  Xét mệnh đề ``$\exists n\in \mathbb{Z}$, $2x^2-8=0$ ''.\\
			Phương trình $2x^2-8=0 \Leftrightarrow x^2=4 \Leftrightarrow x=-2;x=2\in  \Bbb{Z}$ nên mệnh đề đúng.
		\end{itemize}
	 }
\end{ex}

\Closesolutionfile{ans}

\phanDe{\textbf{PHẦN II. Câu trắc nghiệm đúng sai.} Thí sinh trả lời từ câu 1 đến câu 4. Trong mỗi ý a), b), c), d) ở mỗi câu, thí sinh chọn đúng hoặc sai.}

\setcounter{ex}{0}

\Opensolutionfile{ansbook}[Ansbook/GHK1_12_103_DS]

\begin{ex}%[0D1H1-3]%[ma:f76fc24c]
	Cho  mệnh đề $P(x):$ \lq\lq $\forall x\in \mathbb{R},x^2-4x+1> 0$\rq\rq.
	 \choiceTF
		{$P(x)$: \lq\lq Có ít nhất một số thực $x$ thoả mãn $x^2-4x+1> 0$\rq\rq} 
		{$P(x)$ là mệnh đề chứa biến} 
		{\True  $\overline{P(x)}$: \lq\lq $\exists x\in \mathbb{R}, x^2-4x+1\le 0$\rq\rq}
		{Mệnh đề $\overline{P(x)}$ là mệnh  đề sai}
	 \loigiai{
		\begin{itemchoice}
			\itemch  SAI.\\
			Ta có phát biểu đúng của mệnh đề $P(x):$\lq\lq Với mọi số thực  $x$ thoả mãn $x^2-4x+1> 0$\rq\rq.
			\itemch SAI.\\ 
			$P(x)$ là mệnh đề sai nên không phải là  mệnh đề chứa biến.
			\itemch  ĐÚNG.\\
			Mệnh đề phủ định  của $P(x)$ là $\overline{P(x)}:$ \lq\lq $\exists x\in \mathbb{R}, x^2-4x+1\le 0$\rq\rq.
			\itemch SAI.\\ 
			Ta có $\overline{P(x)}:$ \lq\lq $\exists x\in \mathbb{R}, x^2-4x+1\le 0$\rq\rq. Mệnh đề $\overline{P(x)}$ đúng vì giả sử ta  chọn $x=1$ thì $1^2-4.1+1=-2< 0$.
		\end{itemchoice}
	} 
\end{ex}

\begin{ex}%[0D1H1-2]%[ma:bff4f470]
	Cho $\Delta ABC$ có $A'$, $B'$, $C'$ lần lượt là các trung điểm của các  cạnh $BC$, $CA$, $AB$. Khi đó
	 \choiceTF
		{\True $B{C'}=C' A=A'{B'}=\dfrac{AB}{2}{\rm}$} 
		{Hai vectơ $\overrightarrow{BC'}$,  $\overrightarrow{A'B'}$ ngược hướng}
		{\True $\overrightarrow{B{C'}}=\overrightarrow{C' A} =\overrightarrow{A'B'}$} 
		{\True $\overrightarrow{B'  C'}=\overrightarrow{CA'}$}
	\loigiai{ 
		\begin{center}
			\begin{tikzpicture}[scale=1,line join=round,line cap=round,font=\footnotesize,>=stealth]
					\def\r{3} 
					\path (180:\r)  coordinate (B)
						(0:\r)  coordinate (C)
						(120:\r) coordinate  (A)
						($(B)!0.5!(C)$) coordinate  (A')
						($(C)!0.5!(A)$) coordinate (B') 
						($(B)!0.5!(A)$)  coordinate (C');
					\draw  (A)--(B)--(C)--cycle;
					\foreach \x/ \goc in  {A/135,B/-135,C/-45,A'/-90,B'/30,C'/180}
					 \fill (\x) circle (1pt)
						 ($(\x)+(\goc:3mm)$) node {$\x$};
				\end{tikzpicture}
		\end{center}
		\begin{itemchoice}
			\itemch Ta có $C'$ là trung điểm của $AB$ và $A'B'$ là đường trung bình của tam giác $ABC$ ứng với  cạnh đáy $AB$ nên $BC'=C'A=A'B'=\dfrac{AB}{2}$.
			\itemch  Vì Hai vectơ $\overrightarrow{BC'}$, $\overrightarrow{A'{B'}}$ cùng hướng.
			\itemch Vì $C'$ là trung điểm của $AB$ và $A'B'$ là đường trung bình của tam giác $ABC$ ứng với cạnh đáy $AB$ nên $\overrightarrow{BC'}=\overrightarrow{C'A}  =\overrightarrow{A'B'}$.
			\itemch Vì $A'$ là trung điểm của $BC$ và $B'C'$ là đường trung bình  của tam giác $ABC$ ứng với cạnh đáy $BC$ nên $\overrightarrow{B'C'}=\overrightarrow{CA'}$.
		\end{itemchoice}
	 }
\end{ex}

\begin{ex}%[0D1V2-1]%[ma:161eac9d]
	 Cho các tập hợp sau: A các số nguyên tố nhỏ hơn 11; $B=\left\{\left. x\in \mathbb{R}\right|3x^2-4x+1=0\right\}$;
	$C=\left\{\left. x\in  \mathbb{N}\right|\left(x^2-5x+6\right)\left(2x+1\right)=0\right\}$; $D=\{\left. x\in \mathbb{Z}\right|\left|x+1\right|<3\}$.
	 \choiceTF
		{\True Tập hợp $A$ có $4$ phần tử} 
		{Tập hợp $B$ có  $3$ phần tử}
		{Tập hợp  $C$ có $3$ phần tử}
		{\True Tập  hợp $D$ có $3$ phần tử}
	 \loigiai{
		\begin{itemchoice}
			\itemch Ta  có các số nguyên tố nhỏ hơn $11$ là $2$; $3$; $5$; $7$. \\
			Vậy $A=\{2; 3; 5; 7\}$. 
			\itemch Ta có $3x^2-4x+1=0\Leftrightarrow \hoac{ 
			&x=1\in \mathbb{R} \\ 
			 &x=\dfrac{1}{3}\in \mathbb{R}.}$ \\
			Vậy $B=\left\{\dfrac{1}{3};  1\right\}$.
			 \itemch Ta có $\left(x^2-5x+6\right)\left(2x+1\right)=0\Leftrightarrow \hoac{
			& x^2-5x+6=0 \\ 
			& 2x+1=0}\Leftrightarrow  \hoac{
			 & x=2\in \mathbb{N} \\
			& x=3\in \mathbb{N}  \\
			&  x=\dfrac{-1}{2}\notin \mathbb{N}.} $\\
			Vậy  $C=\left\{2; 3\right\}$.
			 \itemch Ta có $\heva{
			&x\in \mathbb{Z}  \\
			 &|x+1|<2}\Rightarrow \hoac{
			&x=-2  \\
			 &x=-1 \\
			 &x=0.} $\\
			Vậy $D=\{-2;-1; 0\}$. 
		\end{itemchoice}
	} 
\end{ex}

\begin{ex}%[0D1V1-2]%[ma:d97640b0]
	Cho hình thang $ABCD$ vuông tại $A$ và có $AB=AD=\dfrac{1}{2}DC=a$. Gọi $BF$ là  đường phân giác trong của tam giác $ABD$ ($F\in AD$). Khi đó
	 \choiceTF
		{\True  $CA^2=DA^2+DC^2$}
		{$\left|\overrightarrow{CA}\right|=a\sqrt  3 $}
		 {$\widehat{ABF}=45^\circ$}
		{$\left|\overrightarrow{BF}\right|\approx 2{,}08a$} 
	\loigiai{\begin{center}
		\begin{tikzpicture}[>=stealth,line join=round,line cap=round,font=\footnotesize,scale=0.8]
				\def\a{2.4} 
				\path 
					 (0,4) coordinate (A)
					(0,0) coordinate  (D)
					(8,0) coordinate (C) 
					(4,4)  coordinate (B)
					($(A)!1/\a!(D)$) coordinate  (F);
				 \draw (A)--(B)--(C)--(D)--(A)
					 (A)--(C)
					(F)--(B)--(D); 
				\foreach \x/\g in  {B/60,A/120,C/30,D/-60,F/-140} \fill (\x) circle (1pt) +(\g:3mm) node[scale=.8] {$\x$};
				 \draw pic[thin,draw,angle radius=2mm] {right angle= B--A--D};
				 \draw pic[thin,draw,angle radius=2mm] {right angle= A--D--C};
			\end{tikzpicture}
	\end{center}
	\begin{itemchoice}
		\itemch Áp dụng định lí Pytago cho tam giác $ADC$ vuông  tại $D$.
		\itemch Áp dụng  định lí Pytago cho tam giác $ADC$ vuông tại $D$, ta có\\
		$C{A^2}=D{A^2}+D{C^2}=a^2+(2a)^2=5a^2$\\ 
		 $\Rightarrow \left|\overrightarrow{CA}\right|=CA=a\sqrt 5$.
		\itemch Ta có  $\Delta ABD$ vuông cân tại $A$, do đó\\
		$\widehat{ABD}=45^\circ  \Rightarrow \widehat{ABF}=22{,}5^\circ$.
		\itemch Xét  $\Delta ABF$ vuông tại $A$, ta có\\
		$ \left| \overrightarrow{BF} \right| =BF=\dfrac{AB}{\cos \left(  ABF \right)}=\dfrac{a}{\cos{22{,}5^\circ}}\approx 1{,}08a$.
	\end{itemchoice}
 }
\end{ex}

\Closesolutionfile{ansbook}

\phanDe{\textbf{PHẦN III. Câu trắc nghiệm trả lời ngắn.} Thí sinh trả lời từ câu 1 đến câu 5.}

\setcounter{bt}{0}

\Opensolutionfile{ansbt}[Ansbook/GHK1_12_103_TLN]

\begin{bt}%[0D1H1-6]%[ma:deff5465]
	Cho phát biểu $\forall x\in\mathbb{R}$, $x^{2}+a-5>0$ với $a$ là số thực cho trước. Có bao giá trị  nguyên của $a<20$ để phát biểu đã cho là mệnh đề đúng?
	\shortans{14} 
	 \loigiai{
		Mệnh đề $\forall x\in\mathbb{R}$, $x^{2}+a-5>0$ đúng $\Leftrightarrow  x^{2}>5-a$ đúng $\forall x\in\mathbb{R}$.\\
		Điều này xảy ra khi và chỉ khi  $\max(5-a, 0) < 0$.\\
		Vì $x^2 \ge 0$, nên ta cần  $5-a$ phải nhỏ hơn $0$.\\
		 $5-a<0\Leftrightarrow a>5$.\\
		Kết hợp  với điều kiện $a<20$, ta có $5<a<20$.\\
		Các giá trị nguyên của $a$ là $6, 7, 8, \dots,  19$.\\
		Số giá trị nguyên của  $a$ là $19-6+1=14$.
	} 
\end{bt}

\begin{bt}%[0D1H1-3]%[ma:d05bebba]
	Cho mệnh đề $P$: \lq\lq $x^2+x+1$ là số dương với $x$ là số thực\rq\rq\ và các  mệnh đề sau.
	\begin{itemize}
		\item \lq\lq $x^2+x+1$  là số không dương với $x$ là số thực\rq\rq.
		\item \lq\lq $x^2+x+1$ là khác $0$ với $x$ là số thực\rq\rq. 
		\item \lq\lq $x^2+x+1$ là số âm với $x$ là số thực\rq\rq. 
		\item \lq\lq $x^2+x+1$  là số không âm với $x$ là số thực\rq\rq.
	\end{itemize}
	Có bao nhiêu phát biểu là phủ định của mệnh  đề $P$?
	\shortans{1} 
	 \loigiai{
		Phủ định của mệnh đề $P$ là $\overline{P}$: \lq\lq $x^2+x+1$ là số không dương  với $x$ là số thực\rq\rq.\\
		Vậy  có $1$ mệnh đề thỏa. Các mệnh đề còn lại không phải là phủ định của mệnh đề $P$.
	 }
\end{bt}

\begin{bt}%[0D1K1-3]%[ma:7f419606]
	Cho mệnh đề $P\colon$ \lq\lq Bình phương của một số thực lớn hơn $ 2 $ \rq\rq, tìm số số  nguyên thỏa mãn mệnh đề phủ định của mệnh đề $P$.
	\shortans{3} 
	 \loigiai{
		Ta  có: $P\colon$ \lq\lq $x^2>2$\rq\rq.
		Suy ra $\overline{P}\colon$  \lq\lq $ x^2 \leq 2 $\rq\rq.\\
		Vì $x \in \mathbb{Z}$ nên $x$ là $ -1 $; $ 0 $;  $ 1 $.\\
	Vậy có $ 3 $ số nguyên  thỏa mãn mệnh đề $\overline{P}$.}
\end{bt}

\begin{bt}%[0D1V2-2]%[ma:8c3a2e78]
	Lớp $10$A có $10$ học sinh giỏi Toán, $10$ học sinh giỏi Lý, $11$ học sinh giỏi hóa, $6$ học sinh giỏi cả Toán và Lý, $5$ học sinh giỏi cả Hóa và Lý, $4$ học sinh giỏi cả Toán và Hóa, $3$ học sinh giỏi cả ba môn Toán, Lý,  Hóa. Tính học sinh giỏi ít nhất một trong ba môn (Toán, Lý, Hóa) của lớp $10$A?
	 \shortans{18}
	\loigiai{ 
		Theo giả thiết đề bài cho, ta có  biểu đồ Ven
		\begin{center}
			\begin{tikzpicture}[>=stealth,line join=round,line cap=round,font=\footnotesize,scale=1]
					\draw [thick,green!60!black] (0,0) circle (1.5) 
						(1.5,-0.5)  circle (1)
						(1.5,0.5) circle  (0.8);
					\draw (1,.8) node {$6$} 
						 (1.2,.2) node {$3$}
						(2,.8) node {Lý} 
						 (1.8,.1) node {$5$}
						(1.8,-1)  node {Hoá}
						(.8,-.8)  node {$4$}
						(-1,0) node  {Toán};
				\end{tikzpicture}
		\end{center}
		Dựa vào biểu đồ Ven, ta có học sinh giỏi ít nhất một trong ba môn (Toán, Lý, Hóa)  của lớp $10$A là
		\begin{itemize}
			\item Số học  sinh giỏi Toán $6+4+3=13$.
			\item Số học sinh  giỏi Lý $6+5+3=14$.
			\item Số học  sinh giỏi Hóa $4+5+3=12$.
		\end{itemize}
		Ta lại có 
		\begin{itemize}
			\item Số học  sinh giỏi cả Toán và Lý $6$.
			\item Số học sinh giỏi cả Toán và Hóa $4$. 
			\item Số học sinh giỏi cả Hóa  và Lý $5$.
			\item Và số học  sinh giỏi cả Toán, Lý và Hóa là $3$.
			\item Số học sinh giỏi hơn một môn  là $4+6+5+3=18$.
		\end{itemize}
	 }
\end{bt}

\begin{bt}%[0D1V3-4]%[ma:bd954b33]
	Cho các tập hợp $A=\left(2;+\infty \right)$ và $B=\left[m^2-7;+\infty \right)$ với  $m > 0$. Số giá trị nguyên của $m$ để $A\setminus B$ là một khoảng $\left(a; b\right)$ thoả mãn $b-a$ thuộc đoạn $\left[3; 16\right]$ là?
	\shortans{2} 
	 \loigiai{
		Điều kiện để $A\setminus  B\ne \varnothing$ là $\heva{&m^2-7> 2\\
		&m > 0} \Leftrightarrow  \heva{&m^2> 9\\
		&m > 0} 
		\Leftrightarrow m  > 3$.\\
		 Khi đó $A\setminus B=\left(2;m^2-7\right)\Rightarrow a=2$, $b=m^2-7$.\\
		Ta có $b-a$ thuộc đoạn $\left[3; 16\right]\Leftrightarrow 3\le m^2-9\le 16\Leftrightarrow 12\le  m^2\le 25$.\\
		Do $m > 3$ và $m\in \mathbb{Z}$ nên $m\in \left\{4; 5\right\}$.\\ Vậy có  $2$ giá trị $m$ nguyên thoả mãn.
	} 
\end{bt}

\Closesolutionfile{ansbt}

\hetDe{103}

\phanDapAn{103}
\par\noindent\textbf{PHẦN I.}\par
\begin{flushleft}
\begin{Solution}{1}
C
\end{Solution}
\begin{Solution}{2}
D
\end{Solution}
\begin{Solution}{3}
D
\end{Solution}
\begin{Solution}{4}
D
\end{Solution}
\begin{Solution}{5}
B
\end{Solution}
\begin{Solution}{6}
C
\end{Solution}
\begin{Solution}{7}
A
\end{Solution}
\begin{Solution}{8}
D
\end{Solution}
\begin{Solution}{9}
D
\end{Solution}
\begin{Solution}{10}
B
\end{Solution}
\begin{Solution}{11}
A
\end{Solution}
\begin{Solution}{12}
D
\end{Solution}

\end{flushleft}
\par\noindent\textbf{PHẦN II.}\par
\begin{flushleft}
\def\writeANS{\TLsai{A}\TLsai{B}\TLdung{C}\TLsai{D}}
\begin{Solbook}{1}
 \writeANS 
\end{Solbook}
\def\writeANS{\TLdung{A}\TLsai{B}\TLdung{C}\TLdung{D}}
\begin{Solbook}{2}
 \writeANS 
\end{Solbook}
\def\writeANS{\TLdung{A}\TLsai{B}\TLsai{C}\TLdung{D}}
\begin{Solbook}{3}
 \writeANS 
\end{Solbook}
\def\writeANS{\TLdung{A}\TLsai{B}\TLsai{C}\TLsai{D}}
\begin{Solbook}{4}
 \writeANS 
\end{Solbook}

\end{flushleft}
\par\noindent\textbf{PHẦN III.}\par
\begin{flushleft}
\begin{loigiaibt}{1}
 14 
\end{loigiaibt}
\begin{loigiaibt}{2}
 1 
\end{loigiaibt}
\begin{loigiaibt}{3}
 3 
\end{loigiaibt}
\begin{loigiaibt}{4}
 18 
\end{loigiaibt}
\begin{loigiaibt}{5}
 2 
\end{loigiaibt}

\end{flushleft}

\newpage
\dauDe{104}

\phanDe{\textbf{PHẦN I. Câu trắc nghiệm nhiều phương án lựa chọn.} Thí sinh trả lời từ câu 1 đến câu 12. Mỗi câu hỏi thí sinh chỉ chọn một phương án.}

\setcounter{ex}{0}

\Opensolutionfile{ans}[Ans/GHK1_12_104_TN]

\begin{ex}%[0D1B3-1]%[ma:a97e4b36]
	Cho hai tập hợp $A=\left\{1;2;3;4\right\}$, $B=\left\{ 2;4;6;8 \right\}$. Tập hợp nào sau đây  bằng tập hợp $A\cap B$?
	 \choice
		{$\left\{1;2;3;4;6;8 \right\}$}
		 {$\left\{6;8  \right\}$}
		{\True$\left\{ 2;4  \right\}$}
		 {$\left\{1;3 \right\}$}
	\loigiai{ 
		Ta có $A\cap B=\left\{ 2;4 \right\}.$ 
	 }
\end{ex}

\begin{ex}%[0D1B1-3]%[ma:5c0c30c1]
	Mệnh đề nào sau đây là phủ định của  mệnh đề “Mọi động vật đều di chuyển” ?
	 \choice
		{ Mọi động vật đều đứng  yên}
		{ Mọi động vật đều không di  chuyển}
		{ Có ít nhất một động vật di  chuyển}
		{ \True Có  ít nhất một động vật không di chuyển}
	 \loigiai{
		Theo định nghĩa, phủ định của mệnh đề “Mọi động vật đều di chuyển” là mệnh đề “Có ít nhất một động vật không di  chuyển”.
	 }
\end{ex}

\begin{ex}%[0D1N2-3]%[ma:b26dadc3]
	Trong các tập hợp sau, tập nào là  tập rỗng?
	 \choice
		 {$A= \{ x \in \mathbb{R} |~ x^2+5x-6=0 \}$}
		{\True $C=\{ x\in \mathbb{Z} |~ x^2+x-1=0  \}$}
		{$B=\{  x \in \mathbb{Q} |~ 3x^2-5x+2=0 \}$}
		{$D=\{ x \in \mathbb{R} |~ x^2+5x-1=0 \}$} 
	 \loigiai{
		Ta có\\ 
		 $x^2+5x-6=0
		\Leftrightarrow \hoac{&x=1 \in \mathbb{R}\\&x=-6 \in \mathbb{R}}$. Vậy $A=\{ -6;1  \} \neq \varnothing$.\\
		$3x^2-5x+2=0 
		\Leftrightarrow \hoac{&x=1 \in \mathbb{Q}\\&x=\dfrac{2}{3} \in \mathbb{Q}}$. Vậy $B=\left\{ 1;\dfrac{2}{3}  \right\} \neq \varnothing$.\\
		 $x^2+x-1=0
		\Leftrightarrow \hoac{&x=\dfrac{-1+\sqrt{5}}{2} \notin \mathbb{Z}\\&x=\dfrac{-1-\sqrt{5}}{2} \notin \mathbb{Z}}$. Vì $x \in \mathbb{Z}$  nên $C=\varnothing$.\\
		$x^2+5x-1=0 
		\Leftrightarrow \hoac{&x=\dfrac{-5+\sqrt{29}}{2} \in \mathbb{R}\\&x=\dfrac{-5-\sqrt{29}}{2} \in \mathbb{R}}$.  Vậy $D=\left\{ \dfrac{-5+\sqrt{29}}{2};\dfrac{-5-\sqrt{29}}{2} \right\} \neq \varnothing$.
	} 
\end{ex}

\begin{ex}%[0D1B3-3]%[ma:5f619861]
	Để thành lập đội tuyển học sinh giỏi khối $ 10 $, nhà trường tổ chức thi chọn các môn Toán, Văn, Anh trên tổng số $ 111 $ học sinh. Kết quả  có: $ 70 $ học sinh giỏi Toán, $ 65 $ học sinh giỏi Văn, $ 62 $ học sinh giỏi Anh. Trong đó có $ 49 $ học sinh giỏi cả hai môn Văn và Toán, $ 32 $ học sinh giỏi cả hai môn Toán và Anh, $ 34 $ học sinh giỏi cả hai môn Văn và Anh. Xác định số học sinh giỏi cả ba môn Văn, Toán, Anh. Biết rằng có $ 6 $ học sinh không đạt yêu cầu cả ba môn.
	\choice 
		 {$ 34  $}
		{\True $ 23 $}
		{$ 36 $}
		 {$  32 $}
	 \loigiai{
	 \immini
		{ 
		Có $ 111-6=105 $ học sinh thi đạt ít nhất $ 1 $ môn.\\ 
		Gọi $ A $ là số học sinh giỏi môn  Toán và Tiếng Anh nhưng không giỏi Văn.\\
		Gọi $ B $ là số học sinh giỏi môn Toán và Văn nhưng không giỏi Tiếng Anh.\\ 
		Gọi $ C  $ là số học sinh giỏi môn Văn và Tiếng Anh nhưng không giỏi Toán.\\
		Gọi $ D $ là  số học sinh giỏi cả ba môn. Ta có hệ:\\
		$\begin{cases}
		 B+D=49\\
		A+D=32\\ 
		 C+D=34\\
		70+65+62-(A+B+C+2D)=105 
	\end{cases}\\
	\Rightarrow  92=32-D+49-D+34-D+2D\\
	\Rightarrow D=23$.\\ 
	Vậy có $  23 $ học sinh giỏi cả ba môn.
	 }
	{ 
	\begin{tikzpicture}[scale=.8]
			\def\radius{2cm} 
			\coordinate  (ceni);
			\coordinate[xshift=.8*\radius,yshift=.2*\radius] (cenii); 
			\coordinate[xshift=.6*\radius,yshift=-.8*\radius] (ceniii); 
			\draw (ceni)  circle (.9*\radius);
			\draw (cenii)  circle (\radius);
			\draw (ceniii)  circle (\radius);
			 \node[xshift=.3*\radius,yshift=.25*\radius] at (ceni) {A};
			\node[xshift=.1*\radius,yshift=-.5*\radius] at (ceni) {B}; 
			\node[xshift=.45*\radius,yshift=-.25*\radius] at (ceni) {D}; 
			\node[xshift=.1*\radius,yshift=-.55*\radius] at (cenii)  {C};
		\end{tikzpicture}
	} 
} 
\end{ex}

\begin{ex}%[0D1B1-5]%[ma:ca8f23aa]
	Cho mệnh đề ``$\forall x \in  \mathbb{R} : x^2+x+1>0$''. Mệnh đề phủ định của mệnh đề đã cho là
	\choice{\True ``$\exists x \in \mathbb{R} :x^2+x+1  \le 0$''} 
		 {``$\forall x \in \mathbb{R} :x^2+x+1 = 0$''}
		{``$\exists x  \in \mathbb{R} :x^2+x+1 < 0$''}
		{``$\forall x \in \mathbb{R} :x^2+x+1 < 0$''}
	\loigiai{ Mệnh đề phủ định của mệnh đề ``$\forall x \in \mathbb{R} : x^2+x+1>0$'' là mệnh đề  ``$\exists x \in \mathbb{R} :x^2+x+1 \le 0$''.
	 }
\end{ex}

\begin{ex}%[0D1B2-1]%[ma:f3f80b92]
	Cho tập hợp $A = \left\{n \in \mathbb{Z}\big| \,6 \, \vdots \, n\right\}$. Hãy  viết tập hợp $A$ dưới dạng liệt kê?
	\choice{$A = \left\{1;2;3;6\right\}$}
		 {\True $A =  \left\{-6;-3;-2;-1;1;2;3;6\right\}$}
		{$A = \left\{-6;-3;-2;-1\right\}$} 
		{$A  = \left\{-6;-1;1;6\right\}$}
	 \loigiai{
		$A$ là tập các ước nguyên  của $6$, do đó $A = \left\{-6;-3;-2;-1;1;2;3;6\right\}$.
	} 
\end{ex}

\begin{ex}%[0D1H3-5]%[ma:1d662c9d]
	Lớp 10A có $45$ học sinh trong đó có $25$ em học giỏi môn Toán, $23$ em học giỏi  môn Lý, $20$ em học giỏi môn Hóa, $11$ em học giỏi cả môn Toán và môn Lý, $8$ em học giỏi cả môn Lý và môn Hóa, $9$ em học giỏi cả môn Toán và môn Học. Hỏi lớp $10$ A có bao nhiêu bạn học giỏi cả ba môn Toán, Lý, Hóa, biết rằng mỗi học sinh trong lớp học giỏi ít nhất một trong $3$ môn Toán, Lý, Hóa?
	\choice 
		{$6$} 
		 {\True $5$}
		 {$3$}
		{$4$} 
	 \loigiai{
		\begin{center}
			\begin{tikzpicture}[scale=.8,font=\footnotesize,line join=round,line cap=round,>=stealth]
					\coordinate  (A) at (0,0);
					 \coordinate (B) at (2,0);
					\coordinate  (C) at (1,-1.5);
					\begin{scope}
						 \clip (C) circle (1.5);
						 \clip (A) circle (1.5);
						\fill[pattern=north west lines]  (B) circle (1.5);
					\end{scope}
					\draw (A)  circle (1.5);
					\draw (B)  circle (1.5);
					\draw (C) circle  (1.5);
					 \path
						 (A)++(90:.6) node[fill=white,above]{$L$}
						(B)++(90:.6)  node[fill=white,above]{$T$}
						(C)++(-90:1)  node[fill=white,above]{$H$};
				\end{tikzpicture}
		\end{center}
		Gọi $T$, $L$, $H$ lần lượt là tập hợp các học sinh giỏi môn Toán,  Lý, Hóa.\\
		Khi đó 
		\begin{eqnarray*}
			|T \cup L \cup H|&=&|T|+|L|+|H|-|T \cap L|-|L \cap H|-|H \cap T|+|T \cap L  \cap H| \\
			 45&=&25+23+20-11-8-9+|T \cap L \cap H|\\
			|T  \cap L \cap H|&=&5
		\end{eqnarray*}
		Vậy có $5$ học  sinh giỏi cả $3$ môn.
	 }
\end{ex}

\begin{ex}%[0D1H1-5]%[ma:9a292bba]
	Trong các mệnh đề dưới đây, mệnh đề nào  là mệnh đề đúng?
	 \choice
		{\lq\lq $\exists x\in  \mathbb{R}\mid x^2< 0$\rq\rq} 
		{\True \lq\lq $\exists x\in \mathbb{N}\mid 2x^2-1<  0$\rq\rq}
		{\lq\lq $\exists x\in \mathbb{N}\mid x^2-2=0$\rq\rq}
		{\lq\lq $\exists x\in \mathbb{R}\mid x^2+1=0$\rq\rq} 
	 \loigiai{
		Ta  có $x^2\ge 0\Leftrightarrow x^2+1\ge 1$, $\forall x\in \mathbb{R}$.
		Nên  $\exists x\in \mathbb{R}\mid x^2+1=0$ là mệnh đề sai.\\
		Ta có $x^2\ge 0$, $\forall x\in \mathbb{R}$. Nên $\exists x\in \mathbb{R}\mid x^2< 0$ là  mệnh đề sai.\\
		$2x^2-1< 0\Leftrightarrow x^2< \dfrac{1}{2} \Leftrightarrow-\dfrac{\sqrt{2}}{2} < x < \dfrac{\sqrt{2}}{2}$, mà $x\in \mathbb{N}\Rightarrow x=0$. Vậy $\exists x\in \mathbb{N}\mid 2x^2-1< 0$ là mệnh  đề đúng.\\
		$x^2-2=0\Leftrightarrow x=\pm \sqrt{2} \notin \mathbb{N}$. Vậy $\exists x\in  \mathbb{N}\mid x^2-2=0$ là mệnh đề sai.
	 }
\end{ex}

\begin{ex}%[0D1H2-3]%[ma:e60745e5]
	Tìm tất cả các giá trị của  tham số $m$ sao cho $(m;+\infty)\subset(1;+\infty)$?
	 \choice
		{$m>1$} 
		 {\True $m\geq 1$}
		{$m\leq 1$} 
		 {$m<1$}
	 \loigiai{
		Do $(m;+\infty)\subset(1;+\infty)$  nên $m\geq 1$.
	} 
\end{ex}

\begin{ex}%[0D1H3-3]%[ma:76b2c2a1]
	Tổng các  số nguyên của tập $[-3;2]\setminus(1;2)$ bằng
	 \choice
		 {\True$-3$}
		 {$4$}
		 {$-5$}
		{$6$} 
	 \loigiai{
		Ta có $[-3;2]\setminus(1;2)=[3;1]\cup\{2\}$ nên chứa các số nguyên  là $-3;-2;-1;0;1;2$.\\
		Vậy tổng các số nguyên là $-3$. 
	 }
\end{ex}

\begin{ex}%[0D1K1-5]%[ma:2b89aaa2]
	Trong các mệnh đề sau,  mệnh đề nào \textbf{sai}?
	\choice 
		{$\exists n\in \mathbb{Z}$, $2x^2-8=0$}
		{Tồn tại số nguyên  tố chia hết cho $5$} 
		{$\exists n\in \mathbb{N}$, $n^2+11n+2$ chia  hết cho $11$}
		 {\True $\exists n\in \mathbb{N}$, $n^2+1$ chia hết cho $4$}
	 \loigiai{
		\begin{itemize}
			\item Xét  mệnh đề ``$\exists n\in \mathbb{N}$, $n^2+11n+2$ chia hết cho $11$''.\\
			Khi $n=3$ thì  giá trị của $(n^2+11n+2)$ bằng $44\vdots 11$ nên mệnh đề đúng.
			\item Xét mệnh đề ``$\exists n\in \mathbb{N}$, $n^2+1$  chia hết cho $4$''.\\
			Khi  $n=2k,k\in \mathbb{N} \Rightarrow n^2+1=4k^2+1$ không chia hết cho $4$, $k\in \mathbb{N}$.\\
			Khi $n=2k+1,k\in \mathbb{N} \Rightarrow n^2+1=(2k+1)^2+1=4k^2+4k+2$ không  chia hết cho $4$, $k\in \mathbb{N}$.
			\item Xét mệnh đề ``Tồn tại  số nguyên tố chia hết cho $5$''.\\
			Tồn tại số nguyên tố $5$ chia hết cho $5$  nên mệnh đề đúng.
			\item  Xét mệnh đề ``$\exists n\in \mathbb{Z}$, $2x^2-8=0$ ''.\\
			Phương trình $2x^2-8=0 \Leftrightarrow x^2=4 \Leftrightarrow x=-2;x=2\in  \Bbb{Z}$ nên mệnh đề đúng.
		\end{itemize}
	 }
\end{ex}

\begin{ex}%[0D1K1-3]%[ma:a990ffed]
	Phát biểu mệnh đề  phủ định của mệnh đề "Tồn tại một loài thú có nhiệt độ thân thể nhỏ hơn $35^\circ$C".
	\choice 
		{"Tồn tại một loài thú có nhiệt  độ thân thể lớn hơn $35^\circ$C"}
		{"Mọi loài thú có nhiệt độ thân thể  nhỏ hơn $35^\circ$C"}
		{"Mọi loài thú có nhiệt độ thân thể lớn hơn $35^\circ$C"}
		{\True "Mọi loài thú có nhiệt độ thân thể  không nhỏ hơn $35^\circ$C"} 
	\loigiai{Mệnh đề phủ định của mệnh đề "Tồn tại một loài thú có nhiệt độ thân thể nhỏ hơn $35^\circ$C" là  mệnh đề "Mọi loài thú có nhiệt độ thân thể không nhỏ hơn $35^\circ$C".}
\end{ex}

\Closesolutionfile{ans}

\phanDe{\textbf{PHẦN II. Câu trắc nghiệm đúng sai.} Thí sinh trả lời từ câu 1 đến câu 4. Trong mỗi ý a), b), c), d) ở mỗi câu, thí sinh chọn đúng hoặc sai.}

\setcounter{ex}{0}

\Opensolutionfile{ansbook}[Ansbook/GHK1_12_104_DS]

\begin{ex}%[0D1H1-3]%[ma:f76fc24c]
	Cho  mệnh đề $P(x):$ \lq\lq $\forall x\in \mathbb{R},x^2-4x+1> 0$\rq\rq.
	 \choiceTF
		{$P(x)$: \lq\lq Có ít nhất một số thực $x$ thoả mãn $x^2-4x+1> 0$\rq\rq} 
		{$P(x)$ là mệnh đề chứa biến} 
		{\True  $\overline{P(x)}$: \lq\lq $\exists x\in \mathbb{R}, x^2-4x+1\le 0$\rq\rq}
		{Mệnh đề $\overline{P(x)}$ là mệnh  đề sai}
	 \loigiai{
		\begin{itemchoice}
			\itemch  SAI.\\
			Ta có phát biểu đúng của mệnh đề $P(x):$\lq\lq Với mọi số thực  $x$ thoả mãn $x^2-4x+1> 0$\rq\rq.
			\itemch SAI.\\ 
			$P(x)$ là mệnh đề sai nên không phải là  mệnh đề chứa biến.
			\itemch  ĐÚNG.\\
			Mệnh đề phủ định  của $P(x)$ là $\overline{P(x)}:$ \lq\lq $\exists x\in \mathbb{R}, x^2-4x+1\le 0$\rq\rq.
			\itemch SAI.\\ 
			Ta có $\overline{P(x)}:$ \lq\lq $\exists x\in \mathbb{R}, x^2-4x+1\le 0$\rq\rq. Mệnh đề $\overline{P(x)}$ đúng vì giả sử ta  chọn $x=1$ thì $1^2-4.1+1=-2< 0$.
		\end{itemchoice}
	} 
\end{ex}

\begin{ex}%[0D1H1-2]%[ma:bff4f470]
	Cho $\Delta ABC$ có $A'$, $B'$, $C'$ lần lượt là các trung điểm của các  cạnh $BC$, $CA$, $AB$. Khi đó
	 \choiceTF
		{\True $B{C'}=C' A=A'{B'}=\dfrac{AB}{2}{\rm}$} 
		{Hai vectơ $\overrightarrow{BC'}$,  $\overrightarrow{A'B'}$ ngược hướng}
		{\True $\overrightarrow{B{C'}}=\overrightarrow{C' A} =\overrightarrow{A'B'}$} 
		{\True $\overrightarrow{B'  C'}=\overrightarrow{CA'}$}
	\loigiai{ 
		\begin{center}
			\begin{tikzpicture}[scale=1,line join=round,line cap=round,font=\footnotesize,>=stealth]
					\def\r{3} 
					\path (180:\r)  coordinate (B)
						(0:\r)  coordinate (C)
						(120:\r) coordinate  (A)
						($(B)!0.5!(C)$) coordinate  (A')
						($(C)!0.5!(A)$) coordinate (B') 
						($(B)!0.5!(A)$)  coordinate (C');
					\draw  (A)--(B)--(C)--cycle;
					\foreach \x/ \goc in  {A/135,B/-135,C/-45,A'/-90,B'/30,C'/180}
					 \fill (\x) circle (1pt)
						 ($(\x)+(\goc:3mm)$) node {$\x$};
				\end{tikzpicture}
		\end{center}
		\begin{itemchoice}
			\itemch Ta có $C'$ là trung điểm của $AB$ và $A'B'$ là đường trung bình của tam giác $ABC$ ứng với  cạnh đáy $AB$ nên $BC'=C'A=A'B'=\dfrac{AB}{2}$.
			\itemch  Vì Hai vectơ $\overrightarrow{BC'}$, $\overrightarrow{A'{B'}}$ cùng hướng.
			\itemch Vì $C'$ là trung điểm của $AB$ và $A'B'$ là đường trung bình của tam giác $ABC$ ứng với cạnh đáy $AB$ nên $\overrightarrow{BC'}=\overrightarrow{C'A}  =\overrightarrow{A'B'}$.
			\itemch Vì $A'$ là trung điểm của $BC$ và $B'C'$ là đường trung bình  của tam giác $ABC$ ứng với cạnh đáy $BC$ nên $\overrightarrow{B'C'}=\overrightarrow{CA'}$.
		\end{itemchoice}
	 }
\end{ex}

\begin{ex}%[0D1V2-1]%[ma:161eac9d]
	 Cho các tập hợp sau: A các số nguyên tố nhỏ hơn 11; $B=\left\{\left. x\in \mathbb{R}\right|3x^2-4x+1=0\right\}$;
	$C=\left\{\left. x\in  \mathbb{N}\right|\left(x^2-5x+6\right)\left(2x+1\right)=0\right\}$; $D=\{\left. x\in \mathbb{Z}\right|\left|x+1\right|<3\}$.
	 \choiceTF
		{\True Tập hợp $A$ có $4$ phần tử} 
		{Tập hợp $B$ có  $3$ phần tử}
		{Tập hợp  $C$ có $3$ phần tử}
		{\True Tập  hợp $D$ có $3$ phần tử}
	 \loigiai{
		\begin{itemchoice}
			\itemch Ta  có các số nguyên tố nhỏ hơn $11$ là $2$; $3$; $5$; $7$. \\
			Vậy $A=\{2; 3; 5; 7\}$. 
			\itemch Ta có $3x^2-4x+1=0\Leftrightarrow \hoac{ 
			&x=1\in \mathbb{R} \\ 
			 &x=\dfrac{1}{3}\in \mathbb{R}.}$ \\
			Vậy $B=\left\{\dfrac{1}{3};  1\right\}$.
			 \itemch Ta có $\left(x^2-5x+6\right)\left(2x+1\right)=0\Leftrightarrow \hoac{
			& x^2-5x+6=0 \\ 
			& 2x+1=0}\Leftrightarrow  \hoac{
			 & x=2\in \mathbb{N} \\
			& x=3\in \mathbb{N}  \\
			&  x=\dfrac{-1}{2}\notin \mathbb{N}.} $\\
			Vậy  $C=\left\{2; 3\right\}$.
			 \itemch Ta có $\heva{
			&x\in \mathbb{Z}  \\
			 &|x+1|<2}\Rightarrow \hoac{
			&x=-2  \\
			 &x=-1 \\
			 &x=0.} $\\
			Vậy $D=\{-2;-1; 0\}$. 
		\end{itemchoice}
	} 
\end{ex}

\begin{ex}%[0D1V1-2]%[ma:d97640b0]
	Cho hình thang $ABCD$ vuông tại $A$ và có $AB=AD=\dfrac{1}{2}DC=a$. Gọi $BF$ là  đường phân giác trong của tam giác $ABD$ ($F\in AD$). Khi đó
	 \choiceTF
		{\True  $CA^2=DA^2+DC^2$}
		{$\left|\overrightarrow{CA}\right|=a\sqrt  3 $}
		 {$\widehat{ABF}=45^\circ$}
		{$\left|\overrightarrow{BF}\right|\approx 2{,}08a$} 
	\loigiai{\begin{center}
		\begin{tikzpicture}[>=stealth,line join=round,line cap=round,font=\footnotesize,scale=0.8]
				\def\a{2.4} 
				\path 
					 (0,4) coordinate (A)
					(0,0) coordinate  (D)
					(8,0) coordinate (C) 
					(4,4)  coordinate (B)
					($(A)!1/\a!(D)$) coordinate  (F);
				 \draw (A)--(B)--(C)--(D)--(A)
					 (A)--(C)
					(F)--(B)--(D); 
				\foreach \x/\g in  {B/60,A/120,C/30,D/-60,F/-140} \fill (\x) circle (1pt) +(\g:3mm) node[scale=.8] {$\x$};
				 \draw pic[thin,draw,angle radius=2mm] {right angle= B--A--D};
				 \draw pic[thin,draw,angle radius=2mm] {right angle= A--D--C};
			\end{tikzpicture}
	\end{center}
	\begin{itemchoice}
		\itemch Áp dụng định lí Pytago cho tam giác $ADC$ vuông  tại $D$.
		\itemch Áp dụng  định lí Pytago cho tam giác $ADC$ vuông tại $D$, ta có\\
		$C{A^2}=D{A^2}+D{C^2}=a^2+(2a)^2=5a^2$\\ 
		 $\Rightarrow \left|\overrightarrow{CA}\right|=CA=a\sqrt 5$.
		\itemch Ta có  $\Delta ABD$ vuông cân tại $A$, do đó\\
		$\widehat{ABD}=45^\circ  \Rightarrow \widehat{ABF}=22{,}5^\circ$.
		\itemch Xét  $\Delta ABF$ vuông tại $A$, ta có\\
		$ \left| \overrightarrow{BF} \right| =BF=\dfrac{AB}{\cos \left(  ABF \right)}=\dfrac{a}{\cos{22{,}5^\circ}}\approx 1{,}08a$.
	\end{itemchoice}
 }
\end{ex}

\Closesolutionfile{ansbook}

\phanDe{\textbf{PHẦN III. Câu trắc nghiệm trả lời ngắn.} Thí sinh trả lời từ câu 1 đến câu 5.}

\setcounter{bt}{0}

\Opensolutionfile{ansbt}[Ansbook/GHK1_12_104_TLN]

\begin{bt}%[0D1H1-3]%[ma:d05bebba]
	Cho mệnh đề $P$: \lq\lq $x^2+x+1$ là số dương với $x$ là số thực\rq\rq\ và các  mệnh đề sau.
	\begin{itemize}
		\item \lq\lq $x^2+x+1$  là số không dương với $x$ là số thực\rq\rq.
		\item \lq\lq $x^2+x+1$ là khác $0$ với $x$ là số thực\rq\rq. 
		\item \lq\lq $x^2+x+1$ là số âm với $x$ là số thực\rq\rq. 
		\item \lq\lq $x^2+x+1$  là số không âm với $x$ là số thực\rq\rq.
	\end{itemize}
	Có bao nhiêu phát biểu là phủ định của mệnh  đề $P$?
	\shortans{1} 
	 \loigiai{
		Phủ định của mệnh đề $P$ là $\overline{P}$: \lq\lq $x^2+x+1$ là số không dương  với $x$ là số thực\rq\rq.\\
		Vậy  có $1$ mệnh đề thỏa. Các mệnh đề còn lại không phải là phủ định của mệnh đề $P$.
	 }
\end{bt}

\begin{bt}%[0D1H1-6]%[ma:deff5465]
	Cho phát biểu $\forall x\in\mathbb{R}$, $x^{2}+a-5>0$ với $a$ là số thực cho trước. Có bao giá trị  nguyên của $a<20$ để phát biểu đã cho là mệnh đề đúng?
	\shortans{14} 
	 \loigiai{
		Mệnh đề $\forall x\in\mathbb{R}$, $x^{2}+a-5>0$ đúng $\Leftrightarrow  x^{2}>5-a$ đúng $\forall x\in\mathbb{R}$.\\
		Điều này xảy ra khi và chỉ khi  $\max(5-a, 0) < 0$.\\
		Vì $x^2 \ge 0$, nên ta cần  $5-a$ phải nhỏ hơn $0$.\\
		 $5-a<0\Leftrightarrow a>5$.\\
		Kết hợp  với điều kiện $a<20$, ta có $5<a<20$.\\
		Các giá trị nguyên của $a$ là $6, 7, 8, \dots,  19$.\\
		Số giá trị nguyên của  $a$ là $19-6+1=14$.
	} 
\end{bt}

\begin{bt}%[0D1V3-4]%[ma:bd954b33]
	Cho các tập hợp $A=\left(2;+\infty \right)$ và $B=\left[m^2-7;+\infty \right)$ với  $m > 0$. Số giá trị nguyên của $m$ để $A\setminus B$ là một khoảng $\left(a; b\right)$ thoả mãn $b-a$ thuộc đoạn $\left[3; 16\right]$ là?
	\shortans{2} 
	 \loigiai{
		Điều kiện để $A\setminus  B\ne \varnothing$ là $\heva{&m^2-7> 2\\
		&m > 0} \Leftrightarrow  \heva{&m^2> 9\\
		&m > 0} 
		\Leftrightarrow m  > 3$.\\
		 Khi đó $A\setminus B=\left(2;m^2-7\right)\Rightarrow a=2$, $b=m^2-7$.\\
		Ta có $b-a$ thuộc đoạn $\left[3; 16\right]\Leftrightarrow 3\le m^2-9\le 16\Leftrightarrow 12\le  m^2\le 25$.\\
		Do $m > 3$ và $m\in \mathbb{Z}$ nên $m\in \left\{4; 5\right\}$.\\ Vậy có  $2$ giá trị $m$ nguyên thoả mãn.
	} 
\end{bt}

\begin{bt}%[0D1V2-2]%[ma:8c3a2e78]
	Lớp $10$A có $10$ học sinh giỏi Toán, $10$ học sinh giỏi Lý, $11$ học sinh giỏi hóa, $6$ học sinh giỏi cả Toán và Lý, $5$ học sinh giỏi cả Hóa và Lý, $4$ học sinh giỏi cả Toán và Hóa, $3$ học sinh giỏi cả ba môn Toán, Lý,  Hóa. Tính học sinh giỏi ít nhất một trong ba môn (Toán, Lý, Hóa) của lớp $10$A?
	 \shortans{18}
	\loigiai{ 
		Theo giả thiết đề bài cho, ta có  biểu đồ Ven
		\begin{center}
			\begin{tikzpicture}[>=stealth,line join=round,line cap=round,font=\footnotesize,scale=1]
					\draw [thick,green!60!black] (0,0) circle (1.5) 
						(1.5,-0.5)  circle (1)
						(1.5,0.5) circle  (0.8);
					\draw (1,.8) node {$6$} 
						 (1.2,.2) node {$3$}
						(2,.8) node {Lý} 
						 (1.8,.1) node {$5$}
						(1.8,-1)  node {Hoá}
						(.8,-.8)  node {$4$}
						(-1,0) node  {Toán};
				\end{tikzpicture}
		\end{center}
		Dựa vào biểu đồ Ven, ta có học sinh giỏi ít nhất một trong ba môn (Toán, Lý, Hóa)  của lớp $10$A là
		\begin{itemize}
			\item Số học  sinh giỏi Toán $6+4+3=13$.
			\item Số học sinh  giỏi Lý $6+5+3=14$.
			\item Số học  sinh giỏi Hóa $4+5+3=12$.
		\end{itemize}
		Ta lại có 
		\begin{itemize}
			\item Số học  sinh giỏi cả Toán và Lý $6$.
			\item Số học sinh giỏi cả Toán và Hóa $4$. 
			\item Số học sinh giỏi cả Hóa  và Lý $5$.
			\item Và số học  sinh giỏi cả Toán, Lý và Hóa là $3$.
			\item Số học sinh giỏi hơn một môn  là $4+6+5+3=18$.
		\end{itemize}
	 }
\end{bt}

\begin{bt}%[0D1K1-3]%[ma:7f419606]
	Cho mệnh đề $P\colon$ \lq\lq Bình phương của một số thực lớn hơn $ 2 $ \rq\rq, tìm số số  nguyên thỏa mãn mệnh đề phủ định của mệnh đề $P$.
	\shortans{3} 
	 \loigiai{
		Ta  có: $P\colon$ \lq\lq $x^2>2$\rq\rq.
		Suy ra $\overline{P}\colon$  \lq\lq $ x^2 \leq 2 $\rq\rq.\\
		Vì $x \in \mathbb{Z}$ nên $x$ là $ -1 $; $ 0 $;  $ 1 $.\\
	Vậy có $ 3 $ số nguyên  thỏa mãn mệnh đề $\overline{P}$.}
\end{bt}

\Closesolutionfile{ansbt}

\hetDe{104}

\phanDapAn{104}
\par\noindent\textbf{PHẦN I.}\par
\begin{flushleft}
\begin{Solution}{1}
C
\end{Solution}
\begin{Solution}{2}
D
\end{Solution}
\begin{Solution}{3}
B
\end{Solution}
\begin{Solution}{4}
B
\end{Solution}
\begin{Solution}{5}
A
\end{Solution}
\begin{Solution}{6}
B
\end{Solution}
\begin{Solution}{7}
B
\end{Solution}
\begin{Solution}{8}
B
\end{Solution}
\begin{Solution}{9}
B
\end{Solution}
\begin{Solution}{10}
A
\end{Solution}
\begin{Solution}{11}
D
\end{Solution}
\begin{Solution}{12}
D
\end{Solution}

\end{flushleft}
\par\noindent\textbf{PHẦN II.}\par
\begin{flushleft}
\def\writeANS{\TLsai{A}\TLsai{B}\TLdung{C}\TLsai{D}}
\begin{Solbook}{1}
 \writeANS 
\end{Solbook}
\def\writeANS{\TLdung{A}\TLsai{B}\TLdung{C}\TLdung{D}}
\begin{Solbook}{2}
 \writeANS 
\end{Solbook}
\def\writeANS{\TLdung{A}\TLsai{B}\TLsai{C}\TLdung{D}}
\begin{Solbook}{3}
 \writeANS 
\end{Solbook}
\def\writeANS{\TLdung{A}\TLsai{B}\TLsai{C}\TLsai{D}}
\begin{Solbook}{4}
 \writeANS 
\end{Solbook}

\end{flushleft}
\par\noindent\textbf{PHẦN III.}\par
\begin{flushleft}
\begin{loigiaibt}{1}
 1 
\end{loigiaibt}
\begin{loigiaibt}{2}
 14 
\end{loigiaibt}
\begin{loigiaibt}{3}
 2 
\end{loigiaibt}
\begin{loigiaibt}{4}
 18 
\end{loigiaibt}
\begin{loigiaibt}{5}
 3 
\end{loigiaibt}

\end{flushleft}
\end{document}
